% ============================================================
%  F. C. Sibbern, Om Elskov eller Kjærlighed imellem Mand og Qvinde
%  Kjøbenhavn: Paa eget Forlag (Trykt hos J. H. Schultz), 1858.
%  "Tredie uforandrede Udgave" — text identical to the 1st ed. (1819).
%  TRANSCRIPTION (Danish)
%  Source: Google Books scan (Univ. of Wisconsin copy),
%    ~/bibliotek/Sibbern, Frederik/Om_elskov_eller_Kjærlighed_imellem_mand.pdf
%  198 PDF pages. Front matter (Fraktur prefaces) + roman body.
%  PAGE MAP: printed page = PDF page - 16  (Første Bog p.1 = PDF 17).
%  See RESUME-NOTES.md for structure, offsets, and transcription progress.
% ============================================================
\documentclass[12pt,a4paper]{book}
\usepackage[utf8]{inputenc}
\usepackage[T1]{fontenc}
\usepackage{libertinus}
\usepackage{libertinust1math}
\usepackage{textalpha}   % Greek words typed directly (e.g. p. 18)
\usepackage[danish]{babel}
\usepackage[protrusion=true,expansion=true]{microtype}
\usepackage{setspace}
\usepackage[margin=1.2in]{geometry}
\usepackage{fancyhdr}
\setlength{\headheight}{28pt}
\addtolength{\topmargin}{-13.5pt}
\usepackage{hyperref}

\hypersetup{
  pdftitle={Om Elskov eller Kjærlighed imellem Mand og Qvinde},
  pdfauthor={F. C. Sibbern},
  hidelinks
}

\pagestyle{fancy}
\fancyhf{}
\fancyhead[LE,RO]{\thepage}
\fancyhead[CE]{\textit{\leftmark}}
\fancyhead[CO]{\textit{\rightmark}}

\setstretch{1.3}

\title{\textbf{Om}\\[6pt]
  \Huge\textbf{Elskov}\\[6pt]
  \large eller\\[6pt]
  \Large\textbf{Kjærlighed imellem Mand og Qvinde}\\[10pt]
  \normalsize af\\[6pt]
  \large Frederik Christian Sibbern.\\[16pt]
  \normalsize Tredie uforandrede Udgave.}
\author{}
\date{Kjøbenhavn 1858.\\
  Paa eget Forlag.\\
  Trykt hos J.\ H.\ Schultz.}

\begin{document}
\maketitle
\thispagestyle{empty}
\cleardoublepage

%% ============================================================
%%  FORTALE TIL ANDEN UDGAVE  (printed pp. III--IV; PDF 11--12)
%%  Fraktur in the original. Transcribed and image-verified.
%% ============================================================
\chapter*{Fortale til anden Udgave.}
\markboth{Fortale til anden Udgave}{Fortale til anden Udgave}

I flere Aar har nærværende Bogs første Oplag været udsolgt, og der
har i nogle Aar været Tale om, at jeg skulde besørge en ny Udgave
deraf. Men hver Gang jeg til denne Ende tog Bogen frem, for at
gjennemgaae den, var det mig imod at læse den, og jeg kunde ikke
overvinde mig dertil. Endelig tilbød en ung Videnskabsmand sig til at
gaae den igjennem for mig, for at rette, hvad der maatte være at rette
med Hensyn til Udtryk og Retskrivning, fornemmeligen for at alt Det
kunde komme til at staae med store Bogstaver, hvad jeg nu, for den
større Tydeligheds og Bestemtheds Skyld, skriver saa, men forhen ikke
skrev saaledes; forresten tænktes der ei paa at gjøre nogen Forandring
i Stil og Foredrag, hvorpaa jeg ikke vilde indlade mig\footnote{At der
paa mangt et Sted til Ordet „han“ kunde være tilføiet „eller hun“,
ville Læserinderne vel snart bemærke.}. Saaledes er jeg da efter mange
Aars Forløb først kommen til igjen at læse denne Bog i Correcturarkene,
og maa jeg tilstaae, at Adskilligt nu forekom mig ikke at være saa
meget Ueffent, men Adskilligt var mig aldeles ikke tilpas, skjønt det
maatte blive staaende, naar Bogen skulde udgaae igjen; thi den skulde
blive som den var; Forandringer --- Smaating fraregnede --- følte jeg
mig ikke i Stand til; de vilde vel ogsaa stukket af mod det Øvrige.
Saa maa jeg da gjentage, hvad jeg seer, at jeg allerede i Slutningen af
Fortalen til første Udgave har sagt, og nu føler mig dobbelt opfordret
til at sige. En Ansamling af psychologiske Bemærkninger og religiøse
og moralske Hentydninger vil man dog finde her, og enkelte Partier af
større Værd. Ved nogle Steder, især i tredie Bog, har jeg tænkt mig,
at, ligesom Romaner oftere ere blevne ledsagede af Kobbere, der, naar
de kom fra en Chodowieckys Haand, kunde kaldes sande værdifulde
Illustrationer, eller Fremhævelser, hvorved det Fremsatte gjennem en
anden Konst fik forhøiende Opklaring og Virkning, saaledes vilde
Adskilligt i denne Bog kunne være egnet til at faae Novelledigtningens
Illustrationer, eller kunne afgive et Slags Themaer til
Novellebehandling. Men endnu er der Noget, der maaskee tør kunne give
Adskilligt, især i første Bog, en særegen Betydning. Det er Dette, at
Forfatteren udtog det ligefrem af hvad han i sit Indre havde ei blot
oplevet men ogsaa levet, medens han i saa Henseende gik i den
skjønneste Skole. Der har stundom været Tale om, at ogsaa han skulde
forsøge at give Skildringer af sit Liv. Saa meget han oftere i fordums
Tid har glædet sig ved Tanken om engang at ville kunne det, især for at
sætte mange herlige Mænd og Qvinder af sin Omverden et taknemmeligt
Minde, tør han dog ei haabe, at han nogensinde vil naae til at begynde
herpaa. Men i nærværende Bog, ligesom i sin psychologiske Pathologie,
og maaskee mangt andet Sted i sine Skrivter (foruden i hans to
Gabrielisser, hvorom der her ei skal tales) har han givet mangen
Livsbekjendelse, og overhovedet tør han vel, omdannende nogle bekjendte
Linier af Goethe, sige:

\begin{verse}
Was ich liebte, was ich lebte,\\
Was ich schaute, was ich strebte,\\
Wollte ausgesprochen seyn;\\
So vertraut' ich oft sub rosa.\\
Beichtend in ganz schlichter Prosa,\\
In der Forschung erstem Hain.
\end{verse}

\noindent Kjøbenhavn, den 1ste November 1853.
\begin{flushright}
Fr.\ Chr.\ Sibbern.
\end{flushright}

%% ============================================================
%%  FORTALE TIL FØRSTE UDGAVE  (printed pp. V--VIII; PDF 13--16)
%%  Fraktur in the original. Transcribed and image-verified.
%% ============================================================
\chapter*{Fortale til første Udgave.}
\markboth{Fortale til første Udgave}{Fortale til første Udgave}

Den nærmeste Anledning til at udarbeide nærværende Skrivt have de
Foredrag givet mig, jeg i adskillige Aar har holdt ved Universitetet
under Navn af Forelæsninger over Lyksalighedslæren. I disse Foredrag
knyttede jeg efter en vis bestemt Tankegang en Række Betragtninger til
hinanden, hvilke gik ud paa et høiere Livs Nydelse under det jordiske
Livs Former og Forhold, uden at jeg dog, ved nogen streng videnskabelig
Skjelnen imellem det til Formaalet egentligen Henhørende og det
Ikke-Henhørende, lod mig afholde fra, ved nogle endog meget omfattende
Sidebetragtninger udførligen at oplyse adskillige af de meest
interessante blandt de mødende Gjenstande fra alle Sider, saavidt Tid
og Foredragets Hovedbestemmelse vilde tillade det. Herved skeede det
da, at nogle Gjenstande afhandledes paa en langt mere omfattende Maade,
end andre i samme Plan medoptagne og dertil væsentligen henhørende. I
et til den roligere Betragtning gjennem Trykken overleveret og for
bestandigen fixeret videnskabeligt Arbeide, indsaae jeg derimod nok, at
en saadan ved visse Hovedpartier indtrædende Afvigelse fra en bestemt
til Grund liggende Hovedtanke ikke vilde tage sig alt for vel ud, endnu
mindre en dermed følgende Ulighed i de forskjellige Deles Behandling.
Hine med særdeles Udførlighed omhandlede Gjenstande viiste sig derfor
mere og mere for mig som de, der langt bedre kunde fremstilles i
særegne for sig bestaaende Udarbeidelser. Saaledes er da nærværende
Skrivt kommet i Stand som et mindre Hele for sig, og saaledes haaber
jeg ogsaa med Tiden særskilt at kunne levere mine philosophiske
Betragtninger over det akademiske Studium, der udgjorde en Deel med af
hine akademiske Forelæsninger, saa og, under Guds Bistand, der derved
dobbelt maa anraabes, mine philosophiske Betragtninger over
Christendommens Hovedlærdomme. Hvad der siden kan blive tilbage, som
passende Stof for en i altid lige Retning fremadskridende Udvikling af
en Lære om menneskelig Lyksalighed, samt om overhovedet en saadan skal
komme i Stand ved mig, vil vel efterhaanden vise sig for mig. Dette
foreløbigen til Svar for dem af mine Tilhørere og Venner, der
velvilligen undertiden have opfordret mig til at udgive mine fornævnte
Forelæsninger, hvilke jeg iøvrigt endnu ikke har udarbeidet skrivtlig.

De her, just ikke i nogen egentlig videnskabelig Form, (hverken hvad
ydre eller indre Form angaaer), meddeelte Betragtninger over Elskov
have af sig selv deelt sig i trende Hoveddele. Disse forekom det mig
meest passende at kalde Bøger, da Benævnelsen af Afsnit, Capitler,
eller Deslige, mere vilde antyde en her ikke anvendt systematisk
Fremgangsmaade; men rigtignok ere de af et noget vel ringe Omfang for
at kunne fortjene hiint Navn. En fjerde allerede for en stor Deel
udarbeidet Bog om den i Elskov medvirksomme Naturdrivt har jeg af
adskillige Grunde holdt tilbage. Den største Deel deraf hørte heller
ikke hen under Bogens Titel, hvorimod den vil have sin rette Plads i
den psychologiske Pathologie, jeg som anden Deel af min Psychologie
haaber engang at levere. Betragtningerne over den ulykkelige Elskov i
tredie Bog ere faldne kortere ud, end jeg fra først af havde tænkt.
Overhovedet har jeg ved dem ofte følt, hvor langt bedre de, da de
individuelle Forhold komme saa meget i Betragtning derved, kunne
oplyses ved en Fremstilling af ganske anden Art. Et Forsøg til en
saadan tør jeg maaskee forelægge Publicum engang\footnote{Herved er
sigtet til de da allerede for største Delen skrevne „Gabrielis's
efterladte Breve“, der udkom første Gang i 1826.}.

Iøvrigt har jeg, uagtet jeg maa tilstaae, at jeg i sin Tid ikke uden
sand Lyst og Glæde nedskrev og flere Gange gjennemgik de fleste af
Hoveddelene af dette Skrivt, dog siden under Trykningen saa vel seet og
følt, hvor langt det er fra at være, hvad det efter Ideen om en saadan
Fremstilling skulde være, at jeg ikke kan Andet end bede Læseren om,
ikke at gaae til dets Læsning med store Fordringer.

\noindent Kjøbenhavn, den 12te December 1819.
\begin{flushright}
Forfatteren.
\end{flushright}

%% ------------------------------------------------------------
%%  EFTERSKRIVT VED TREDIE UDGAVE  (printed p. VIII; PDF 16)
%% ------------------------------------------------------------
\bigskip
\noindent\textit{Efterskrivt ved tredie Udgave,} den 26de April 1858:
Nogle Smaaforandringer, foretagne under Correcturen, have ei kunnet
hindre mig i at kalde nærværende Udgave en uforandret. Den følger ei
blot Side for Side men og Linie for Linie med den foregaaende, der kun
var paa 350 Exemplarer og derfor i en Tid af 4½ Aar har kunnet sælges.
Det nærværende Oplag er dobbelt saa stort, og Prisen er nedsat til 48
Sk.\ fra 80 Sk. --- Iøvrigt vil jeg kun bemærke, at af en Bog, som
denne, skal Enhver tage sig Det, som kan være ham eller hende til Nytte
og Fremme, til Qvægelse og Frydelse og lade det Øvrige ligge.

%% ============================================================
%%  FØRSTE BOG  (printed pp. 1--...; PDF 17--...)
%%  Fraktur in the original. Transcribed and image-verified.
%% ============================================================
\chapter*{Første Bog.}
\markboth{Første Bog}{Første Bog}

Med uudsigelige Følelser fylder den opvaagnende Elskov enhver inderligt
elskende Sjæl, og i et nyt, herligt og rigt Liv udfolder sig den i
lykkelig Elskov gjenfødte Tilværelse. Men kun altfor ofte skjænkede
Naturen Mennesket denne jordiske Velsignelses rigeste Fylde ved en
Lykke, den har bestemt for Enhver, uden at den Lykkelige lod Det, der
tildeeltes ham, saaledes gaae ind i sig og blive sit, at det kunde
vorde ham til en ægte Skat i Livet og til et Glædens bestandig rindende
Væld. Kun alt for ofte gaaer Elskovs skjønne Blomstretid forbi som en
Drøm, og bliver ikke, hvad den kunde blive, Indgangen til en
uomskiftelig Tilværelse. Og dog er ingen jordisk Vei til det sande Liv
skjønnere, end denne, og intet jordiskt Lys lyser klarere for den ægte
Stræben, end Elskovs kraftigt vækkende og inderligt qvægende Lys. Med
indtrængende Stemme forkynder Naturen i denne Følelse Livets indre
Betydning, og kalder ved den ud af det Synliges vidt udbredte
Mangfoldighed ind i det Usynliges dybere Helligdom. For Den, der vil
trænge ind i Livets og Naturens indre Sandhed og Væsenhed, er den
skjønneste Vielse beredet her. Men skal den vorde ham til en saadan
Indvielse, skal af den en varende Glæde ved Tilværelsen udvikle sig hos
ham, saa maa han fatte den i dens indre Væsen og tilegne sig den af
Hjertet ved en dyb Erkjendelse. Ikke at han, eftersporende og sondrende
Følelsernes vexlende Række, saaledes som de maaskee paa en egen og
mærkelig Maade synes ham at fremstille sig i ham, skulde, grublende
over sig selv, falde ud af Livets fulde Fornemmelse, og nu
sønderskille sig det Liv, han forhen levede heelt og som af eet Stykke.
Men fatte og erkjende maa han, hvad der er skjænket ham, for netop at
kunne leve og nyde det inderligen og tilfulde, idet han føler dets
indre i alle vexlende Momenter sig rørende Dybde, og lader denne komme
til fuld Klarhed og til en gjennemtrængende Udbredelse i hans hele
Tilvær, hvilket ikke skeer uden en Erkjendelse. Ved denne maa Elskoven
hæves frem til at vise sin Magt som et ægte Livet besjælende og
dannende Princip, og ligeledes maa den derved finde sin faste Grund,
hvori den kan rodfæstes, og sit høiere Hjem, hvori den kan vorde
forklaret, i det Eviges Sphære. Uden at finde denne Grundvold og at
naae til denne Forklarelse, vil Elskoven blomstre af, som den
blomstrede frem. Menneskets hele aandelige Liv har sin Rod i Kjærlighed
til den Evige og i Troen paa hans uendelige Væsen, og denne Tro og
denne Kjærlighed maa i enhver Henseende udgjøre Hjertet, Livsfølelsernes
indre varmende og bevægende Middelpunct, i alle det aandelige Livs
Lemmer og Dele. --- Derfor, baade for at naae til denne Erkjendelse, og
tillige for overhovedet at fornemme den Naturens Røst, der taler til os
i denne den skjønneste blandt de til en blot jordisk Velsignelse
virkende Livets Magter, ville vi gaae ind i en Eftertænkning over
Elskovs Natur og Indhold. Vi følge kun Naturens Vink, naar vi i
Elskoven, saaledes som den lever i lykkelige, rene og kraftige Sjæle,
og udfolder sig i deres Samliv, søge et Billede paa det ægte Liv. I
denne dens skjønneste og heldigste Skikkelse ville vi da betragte den
først.

De tvende yderste, hinanden meest modsatte, Puncter i Elskovs hele
Fylde, de Holdningspuncter tilmed, hvorom, som om tvende Poler, hele
Betragtningen derover maa dreie sig, ere: det aldeles Jordiske i samme
paa den ene Side, det aldeles mod det Evige Henvendte deri paa den
anden Side. Hiint er den Naturdrivt, hvorved Elskovs Længsler
fremkaldes. Dette er Glæden ved det Skjønne og Guddommelige, der
ligeledes er i Bevægelse, hvor Elskov kommer frem. Ifølge den første
søger Individet sin Halvdeel, i og under den sidste sin Guds Nærhed. I
ægte Elskov finder Den, hvis Stræben er en lykkelig, begge i Eet, i sin
Halvdeel sin Guds herlige Billede. Hvad der her som indre
Naturtilskyndelse ikkun synes at lede til en jordisk Tilfredsstillelse,
kan hos Mennesket, fordi han er Menneske, ikke finde en menneskelig
Fyldestgjørelse uden at føre ham sit evige Udspring usigelig nær og
aabne for ham den Allerhelligstes Tempel.

I et ungdommeligt, for Livets mangfoldige Indvirkninger aabent Sind
udspringer Elskov næsten altid af en Beskuelsens Lyst og Glæde og søger
i denne sin første Næring. Enten saa dens Begyndelse i et Individ er en
pludselig Bevægelse i Sindet og en Henrykkelse, der viser hvor stærkt
Hjertet er truffet, eller den langsomt gaaer frem af et venligt Samqvem
og en skjøn Omgang, hvorunder lidt efter lidt tvende Sjæle udfolde sig
for hinanden: Det, hvad den ene skuer og med Lyst idelig paa Ny glæder
sig ved at beskue i den anden, er i begge Tilfælde Det, der vækker
Hengivenheden og drager det ene Hjerte til det andet. Det er vistnok
Drivten, som ligger i Baggrunden, det er den dybt sig rørende Natur,
som stemmer Sjælen til Tilbøielighed og til Elskovs blødgjørende
Længsel, det er en fra Individets Inderste fremstigende Trang og
Trængsel, der vækker og oplader Øiet for det Velgjørende i et andet
Individ, og gjør, at den Ene kan tiltrækkes saa stærkt af den Anden.
Men saa meget den sig rørende Natur end kan sætte Mennesket i urolig
Bevægelse, saa kan han dog, som Den, hos hvem Sandsen for noget Høiere
er vaagen og overalt er rede til at træde til, ikke, naar han ei
fornægter sin Natur aldeles, knytte sig til sin Attraas Gjenstand, uden
at finde Velbehag i den, og det et sjæleligt Velbehag. Derved bliver
nu, idet den mødende Gjenstand sammenholdes med en indre Fordring, den
høiere Skjønhedsidee vakt til Liv. Længselen bliver nu en Længsel efter
det Skjønne og Tækkelige og en Attraa efter inderligen at forenes med
et skjønt og indtagende menneskeligt Væsen. Og saaledes maae da --- for
at laane et Billede fra den nordiske Mythe om Freias af Katte trukne
Vogn, der saa skjønt er bleven tydet af Oehlenschläger i hans Vaulundurs
Saga --- de Lyster, der ellers overalt i Naturen saa ofte ere vilde og
uregjerlige, nu hos det sundt og reent følende Menneske drage den
høieste Kjærligheds Guddom frem, medens de dog beherskes af den.
Naturens stærke Bevægelser fremkalde en dyb Hjertets Trang og Længsel,
en Trang til Sjælens Fyldestgjørelse og til Sjæleglæde. Denne bliver nu
fremherskende, og hvad der i Naturen i organisk Henseende var det
Første, vedbliver ei mere at være det, men træder tilbage. Elskovs
Længsel bliver først da fængslet ved en bestemt Person, og glæder sig
ved at have funden sin Gjenstand, naar den opvakte Sands bliver greben
og udfyldt ved det inderligt Tiltrækkende og Herlige, der oplader sig
for Sjælens Øie ved Beskuelsen af det andet Væsen. Først da opstaaer
Elskovs levende, Sands, Sind og Hjerte fyldestgjørende Følelse. Den
viser sin første Kraft i den inderlige Fryd, den uudslukkelige Lyst,
hvormed det ene Individ dvæler og hviler i, ja taber sig i, Beskuelsen
af det andet, og deraf som af en uudtømmelig Kilde øser en bestandig
frisk Vederqvægelse, saa at Hjertet bevæges i en levende Glæde,
hvergang den elskede Skikkelse viser sig paa Ny. Og er det ogsaa
stundom ikke Beskuelsens Lyst, hvormed Elskov begynder, er det det
Tilhold, som et Hjerte har fundet hos et andet Hjerte, og den deraf
udspringende Tiltro og Fortrolighed, der har fremlokket Hengivenheden
og Kjærligheden, saa er det dog naturligt, at denne Sympathie nu
efterhaanden drager Contemplationen frem og Glæden ved denne, og først
da kommer det til en egentlig Elskov. Naar et Hjertes inderlige Værd
føles og erkjendes, saa vil det ogsaa snart beskues med sand Lyst i
alle dets Yttringer; men da vil, naar Elskovs Drivt og Længsel derved
sættes i Bevægelse, snart ogsaa den hele ydre Skikkelse, hvori et
saadant Væsen aabenbarer sig, beskues med inderligt Velbehag, og det
gaaer da paa den vederqvægeligste Maade i Opfyldelse, at Kjærlighed
gjør seende.

Men saa meget end denne Glæde ved den Elskede knytter sig til det meest
Eiendommelige og Personlige, og netop derved viser sig som en egentlig
Elskov, der gaaer ud paa en bestemt Person og paa ingen anden end
denne, saa maa den dog have sit Grundlag i en Erkjendelse af Noget, der
i og for sig og uden Hensyn til al personlig Tilbøielighed er godt og
fortrinligt hos det elskede Individ. Enhver, der elsker af Hjertet,
elsker fordi han skuer eller troer at skue noget i sig selv Elskværdigt
hos den Anden; men elskværdigt er et Individ ikkun ved det Høiere, det
Guddommelige, hvilket det fremstiller i sig, eller hvis Præg det bærer
og af hvis levende Kraft det er opfyldt. Netop fordi et Individ synes
os i sig selv at være smukt, tækkeligt, yndigt, af reen eller dyb
Natur, eller overhovedet af høiere Præg, føle vi os trukne hen dertil.
Elskov gaaer frem af et inderligt Velbehag, hvilket forudsætter et
hjerteligt Bifald. Men dette Bifald er kun muligt, naar vi, henførende
Det, vi skue, til en i os selv levende Idee, føle og sande, at det ret
dybt treffer sammen med denne. Ja ikkun ved denne kunne vi overhovedet
erkjende det Fortrinlige udenfor os og fatte dets Værd. Rørte sig den
guddommelige Fylde ikke i vort Inderste, saa vilde vi ei kunne finde
den, og glæde os ved at finde den, udenfor os. Saaledes er da Det, som
allerede ved den blotte Beskuelse drager tvende Væsener til hinanden,
Noget, der er høiere end begge. Ligesom, hvor Tvende ere forsamlede i
Christi Navn, Christus er midt iblandt dem, saaledes er det, hvor
Tvende ere forenede ved en ægte Kjærlighed, det Guddommelige, der
forener dem, hvilket er tilstede i dem Begge, og hvilket den Ene skuer
i den Anden, naar han skuer ind i den Andens Sjæl, som i den for ham
opladte Himmel. Hvor det ei er dettes mægtige Tilstedeværelse, der
holder begge forenede ved en sand Glæde ved hinanden, der er heller
ikke Elskoven bygget paa en varig Grund. Et Blændværk har trukket Begge
til hinanden; thi den Ene har da skuet i den Anden, hvad der ei var i
denne. Ogsaa kan ikkun det i sig selv Gode, som det Ægte og Bestandige,
vække og bevare den dybe Tillid, uden hvilken Elskoven er Intet. Den
blotte Tanke, at det andet Væsen maaskee ei er, hvad det synes, et
reent Kar for den høiere Aand, nager strax paa Baandet mellem Begge, og
Visheden derom er som en ond Dæmon, der staaer imellem dem. Og beholde
end maaskee de ydre Tillokkelser deres Magt, saa er dog den egentlige
Elskov sønderbrudt, og den Tanke, at Det, der endnu holder os fængslet,
er i sit indre Væsen uægte og usandt, qvæler endog den ydre Beskuelses
Lyst. Hvor Tvende troe sig forenede i Elskov, og ikke ere det ved den
alene i Sandhed enigende Magt, et høiere Livs begge besjælende Aand,
ligner Samlivet for Enhver af dem et Liv med en huul Gjenganger, der
synes at være et fyldigt Menneske. Den Gudopfyldte elsker i Alt, hvad
der opfylder ham med dyb Kjærlighed, sin Guds uendelige Nærhed, og hvo
der overhovedet glæder sig dybt og inderligen ved noget Herligt, i ham
rører sig en fra Sjælens Dyb fremkommende Følelse af noget Usynligt og
Evigt, om hvis himmelske Herlighed han ligesom mindes, idet det Samme
taler til ham udaf den ydre Gjenstand. Derfor føler han sig indtil sit
Inderste qvæget, ved Det, han skuer som ved en evig Skjønheds uendelige
Tilstedeværelse i hans Sjæl. Denne er det, vi prise, og hvortil vi
knytte os, hos den Elskede, og Glæden ved den Elskede er da allerede en
sand Glæde ved Herren.

Imidlertid er dog det i begge de Elskende levende Guddommelige og
sammes gjensidige Erkjendelse alene ikke Nok til at fængsle den i
Bevægelse værende Drivt og Længsel, og til derved at fremvirke og
grundfæste Elskov. Inderlig Agtelse og Beundring, ja en Kjærlighed af
en anden Art end Elskov, kan derved vækkes og befæstes. Men først naar
en Glæde ved den personlige Individualitet, den særegne
Eiendommelighed, kommer til, naar vi ei blot blive henrykte og
inderligen glædede ved det Fortrinlige og Herlige i den anden Person,
men ogsaa ved denne Person selv og dens af det Guddommelige forklarede
Personlighed, kan virkelig Elskov udvikle sig af Beskuelsens Lyst, og
Elskovs Længsel finde sin Fylde. Hvad det saa er, der hos den Elskede
kan have været det første Tiltrækkende og Vækkende og tilmed det første
Incitament til Elskov, være det Formens Skjønhed, Trækkenes Udtryk
eller Bevægelsernes Ynde, Retsindighedens Energie og det sikkre Blik
paa Livet, eller Sindets friske Livelighed og Phantasiens aanderige
Spil, den dybere Sands for Tilværelsens indre Væsen, eller en stille i
sig selv beroende og hvilende Fred, eller hvad det endnu ellers monne
have været, der først har truffet og tiltrukket Hjertet, og der maaskee
ogsaa for den fjærnere Iagttager viser sig som det meest Fremspringende
i den elskede Individualitet: saa er dog Det, der nu i Følelsernes
Fylde indtager Elskeren og har hans Hu og Hjerte i sin Besiddelse, ikke
dette eller hiint Enkelte og Særlige alene, men er den hele
Individualitet selv, saaledes som den i et heelt og fuldt Udtryk
udtaler sig udaf den Elskedes hele Existents og alle Yttringer. Det
første Indtryk kan være Virkningen af disse eller hine Tillokkelser
hvilke have henrevet den ungdommelige aabne Sands og hentaget Sindet i
Henrykkelse ved deres opflammende eller oplivende Magt. Et eller andet
Enkelt er ofte Det, der først som Incitament drager Øiet til sig og
vækker Attraaen og Længselen i den efter Elskovs søde Lyst
længselsfulde Sjæl. Derved henrives da Sindet og hengiver sig til
Beskuelsens Glæde, men den enkelte Tillokkelse bliver nu Det, der
ansporer til at trænge ind i den os mødende Sjæls indre Heelhed og
dybere Væsen, den bliver Det, hvortil vi ligesom støtte os og som leder
os i den higende Lyst efter, gjennemtrængende at gjennemskue den
elskede Sjæl og at forfølge den i alle dens Yttringer. Men nu gjelder
det, om dens indre Heelhed saaledes svarer til vor Længsel, at den kan
stille og fylde denne. Ikkun da, naar ei blot den i Øieblikket
henrivende Tillokkelse, men ogsaa den i sin Dybde roligt levende og
virkende Fylde i det elskede Væsen udøver hos os en vedvarende
Indflydelse, kommer det til en virkelig ægte Elskov. Den Elskedes altid
sig selv lige indre Væsen, der med sin blivende Grundcharacteer udgjør
det Individualitetens indre Fond eller Væld, hvoraf Livets ydre
Mangfoldighed og Bevægelse gaaer frem, dette indre Væsen er det, vi
maae have skuet og erkjendt, og hvoraf vi maae føle os inderligen
tiltrukne, forinden vi kunne siges i Sandhed at elske. Hvortil
udfordres, at den første heftige Bevægelse maa dulme, og Sindet maa
komme til en saadan Besindighed og Klarhed, at Elskeren kan fatte sig
og sige sig selv, om det kun er den friske, frodige Livsyttring, eller
det er selve den aandelige Grundskikkelse og indre Fylde hos den
Elskede, til hvilken han med saa stor en Magt føler sig hendraget, og
hvilken han med idelig frisk Lyst længes efter at skue og atter skue.
Sand Elskov kan overhovedet ei udvikle sig af forbigaaende Indtryk, men
ikkun af et fortvarende Liv, hvori Tvende, idet de meget leve med
hinanden, komme til at oplade deres hele indre Eiendommelighed for
hinanden. Man maa leve sig ind i Kjærligheden.

Det er altsaa selve Sjælen, vi, naar Elskoven er ægte, elske i et andet
Individ, eller det indre besjælende aandelige Hele, af hvilket
Individets hele Liv og Eiendommelighed gaaer frem. Dog yttrer Sjælens
indre Herlighed den stærkeste Indflydelse i den Maade, paa hvilken den
udtaler sig i Følelserne og de af disse udspringende Sindsbevægelser og
Sindsstemninger. Den i disses hele Mangfoldighed sig udtalende Sjæl
udgjør det egentlige Hu, Hjerte, Hjertelag, Gemyt. Hermed betegne vi
den indre Grundbevægelse i Sjælen, af hvilken Livsvarmen gaaer frem,
det indre dybe Væld for Følelsernes til alle Sider sig udbredende
Strømme. Erkjendelsens Dybde og Klarhed og Villiens Fasthed og Energie,
de største Fortrin i Henseende til Talent og Characteerstyrke, virke
endnu ei ægte inderlig Elskov; hvorimod et godt Hjerte, et dybtfølende
og reent Gemyt, hvormed en god, skjønt just ikke en kraftig og stærk
Villie altid følger, ofte er Elskeren fuld Erstatning for manglende
Aandskraft og Naturgaver.

Naar Elskov knytter sig til Gemytet og overhovedet til
Individualiteten, til selve Sjælen i det andet Væsen, knytter den sig
til det Blivende og Uforgjængelige, og et uudtømmeligt Velbehag kan da
gaae frem af den. Sjælen er det evige Unge i Mennesket, ja kan, sig selv
lig, endnu udaf den forældede ydre Form straale Elskeren imøde i sin
uforandrede Grundskikkelse. Men heel skue vi den ikke i noget Moment af
den Andens Liv. Kun naar dette Liv i sine mangfoldige Yttringer udbreder
sig for Blikket, lyser Sjælen i sin indre Heelhed frem deraf. Man kunde
sammenligne denne indre sjælelige Individualitet med et Væsen, der,
indhyllet i et Klædebon, under utallige Bevægelser saaledes forraadte
sin Skikkelse for den skuelystne Iagttager, at den tilsidst, uden at
træde frem af sit Dække, stod heel for dennes Phantasie.
Individualiteten udtaler sig i Momenter, men dens totale Udtryk er kun
hele Individet selv. Derfor er her for Elskerens Øine den idelige Nyhed
i den Elskedes mangfoldige Bevægelser og Yttringer, naar Livet hos
denne ved Elskoven er opvakt og bestandig friskt. Den elskede
Grundskikkelse glæder han sig ved at skue, ideligen den samme og dog
ideligen ny, gjentaget i de utallige Livsbevægelser, og finder heri en
Lyst, der er ligesaa uudtømmelig, som Livet og Individualiteten selv.

Den Elskedes livlige Yttringer ere lutter Tillokkelser, da de alle
række det i Elskeren liggende Grundbillede, der er hans Hjerte saa
kjært, og derved vække hans Lyst og Glæde op paa Ny. Jo dybere og
klarere Elskoven er, jo mere inderligt Velbehaget, desto mere udfolder
sig det kjære Billede i vort Indre. Hos Elskere med dyb Følelse og dybt
Sind udvikler sig af deres Glæde ofte et sandt Studium, et Ord, der maa
tages mere egentlig, end det ved første Øiekast kan synes. Er jo dog
alt Studium desto ægtere og desto mere velgjørende, jo mere vi saa
ganske fordybe os i Gjenstandens Beskuelse, at vi derved lade den i sin
fulde Heelhed gaae ind i os og udpræge sig i os, hvorved ogsaa først
tillige vor egen Aands Fylde udvikler sig indenfra.

Elskeren er i sin Kjærlighed en Seer, der formaaer igjennem det ydre
Dække at skue ind til Hjertet i det andet Væsen. Et lykkeligt Seerblik
hører til at vække den ægte Elskov. Men enhver ungdommelig bevæget Sjæl
har lettelig saadanne geniale Øieblikke. Dog kan der stundom virkelig
høre Held til, at et andet Væsens Indre kommer til at oplade sig for os
og en Straale fra den Himmel, som deri ligger beredet for os, treffer
vort Øie i en Stund, da vort Sind er frit og aabent nok til at modtage
et reent Indtryk med Reenhed. Ei sjeldent levede Tvende længe ved Siden
af hinanden, inden et lykkeligt Øieblik bragte dem til at oplade deres
inderste Væsen for hinanden og gjensidigen at erkjende hinandens Værd,
saa at nu begge følte sig som rørte af noget Guddommeligt. Stundom laae
vel ogsaa i et Menneskehjerte den himmelske Gehalt tilsluttet og skjult
som en Diamant i Klippen, indtil en lykkelig Finder erkjendte dens Værd
og kaldte den frem til Dagen, for at den nu i Kjærlighedens Sol kunde
udstraale, det er gjenstraale, en uudsigelig Herlighed af Lys. Ogsaa
skeer det ofte, at et Individ, der syntes Andre, og bliver ved at synes
dem, lidet betydeligt, elskes med en Hjertelighed og et Velbehag og en
vedvarende Tilfredshed og Lyst, der for meget beviser Elskovens indre
Ægthed og Dybde, til at der ikke til Grund for samme skulde ligge noget
særdeles Ægte, men hvilket ikkun Elskeren, der lykkeligen har fattet og
skuet den dybere Grundskikkelse i sin Elskede, har været i Stand til at
erkjende og tilegne sig saa hjerteligen. Vistnok kan Kjærlighed, uagtet
den efter sin indre Natur og Væsen gjør seende, dog undertiden bringe
et Individ til at hildes i Vildfarelse og Blindhed, ligesom en
philosophisk Grunden og Stræben ei sjeldent førte til Sværmerie eller
Ideeforvirring, uagtet Philosophie og philosophisk Grunden dog netop er
Noget, der oplader Aandens høiere Øie og gjør ægte seende. Men
ingenlunde er altid den Elskov fremkaldet ved et Blændværk, der synes
Andre saa, fordi de i den Elskede ikke see noget saa Elskværdigt, at
det kan synes dem i Stand til at yde saa megen Lyst og Fyldestgjørelse.
En menneskelig Sjæl gjemmer ofte meget Mere i sit indre Dyb, end hvad
der lader sig tilsyne for Andre, end Den, hvem Elskoven gjør til Seer.
Det Guddommelige lever jo dog virkeligen i ethvert godt Menneske og
giver Mennesket Væsenhed, være det end overskygget eller trykket
tilbage formedelst jordiske Mangler, hvilke dog letteligen svinde hen
for Den, der erkjender og skuer hiint Guddommelige i dets indre
Reenhed.

Intet er mægtigere til at fremkalde Elskov, end den Legemets Skjønhed,
der er det fulde, i den sandselige Form reent og klart udprægede Udtryk
af et skjønt og himmelsk Indre; thi Skjønheden udfylder vor hele
menneskelige Sands, den høiere tilligemed den lavere. Det hele Menneske
føler sig oplivet og qvæget ved den. Meest virker dog i den ydre Form
Det, i hvilket Sjælen udtaler sig meest umiddelbart, Øiets, Trækkenes
og den i den hele Individualitet udtrykte Bevægelsernes Skjønhed,
hvilken sidste er den over det hele Individ udbredte Ynde. Men Skjønhed
er en Lykke, som kun sjeldent er forundt et elskværdigt Individ. Det er
kun sjeldent, at der i Livets Trængsel undes et med ungdommelig Friskhed
fremblomstrende Væsen Rum og Frihed nok til en rolig og alsidigen
gjennemtrængende Udfoldning af den individuelle Fylde, være det i
sjælelig eller legemlig Henseende, ikke at tale om, at denne saa ofte
allerede er eensidigen anlagt, ja reent fortrykket og forkuet, i den
første Spire, eller i den tidligste Fremvæxt, inden Individet endnu er
kommen til ret Bevidsthed om sig selv.

Dog for Den, der med inderlig Lyst tilegner sig sin Elskov og sin
Elskede, er denne altid skjøn. Og det er aldeles naturligt, at den
elskede Skikkelse udøver Skjønhedens Magt over ham. Al Skjønhed beroer
paa, at Noget, som udspringer af Livets høiere Kilde og bærer dens Præg
i sig, træder fuldeligen frem i en ydre sandselig eller timelig Form,
udfyldende og gjennemtrængende denne, saa at denne bærer det fulde
Udtryk af hiint. Men hvad Elskeren skuer som Elskværdigt, tilmed
Himmelsk, i den Elskedes indre Væsen, dette Samme netop seer han med
inderlig Glæde udtrykt og udtalt i den hele elskede Individualitet og
enhver af dens Yttringer. Og saaledes vil den Elskede ogsaa i
Henseende til sit Ydre gjøre Skjønhedens Indtryk hos Elskeren, og hans
Glæde og Velbehag vil være af samme Art som den, Skjønheden fremkalder.
Er det ikke ogsaa en Straale fra den sande Himmel, som falder i hans
Sjæl igjennem hans Elskede, og vækker denne ikke den uendelige Følelse
af det evige Skjønnes himmelske Herlighed? Men hvad der er i Stand til
at vække hos os denne Følelse, staaer derved selv forklaret for os, og
endog hvor den individuelle Fuldendthed og Fylde fattes, prise vi det
Væsen som skjønt, ved hvis Beskuelse en Mindelse om den evige Skjønhed
bestandigen lever op igjen hos os. Derfor udvikler sig da af Elskov et
Velbehag og en Lyst ei blot ved det i sig Fortrinlige og Væsentlige,
men ogsaa ved det Ubetydeligste i den Elskedes Yttringer og Bevægelser,
fordi den Elskende deri ideligen møder den elskede Individualitet, og
forlyster sig ved at see den udpræget indtil den yderste Overflade af
Formen. Dog fornemmes denne rene Glæde fuldest i det totale Indtryk, i
hvilket alt Enkelt og Mangfoldigt smelter sammen i en eneste rolig og
heel Anskuelse, hvori den Elskedes hele Væsen udfylder vor Sjæl.

Usigelig husvalende og oplivende er denne rene Contemplation, denne
rolige udeelte Hensjunkenhed i Beskuelsen af et elskværdigt Væsen, og
den milde, klare Indflydelse, hvormed da Billedet af dette Væsen gaaer
ind i den Beskuendes Sjæl. Vækkende i vort Indre den beslægtede
aandelige Fylde i den rene Følelse af det Eviges inderlige Høihed og
Herlighed, fremkalder den i os et høiere Livs frie, reent aandelige
Fornemmelser. Som i Solens Lys og Varme organisk Liv udfolder sig i
Naturen, saaledes udfolder sig et nyt aandeligt Liv ved Kjærlighedens
varmende Lys. Den med inderlig Lyst erkjendte og beskuede Individualitet
bringer vor egen Aands indre Grundskikkelse og Væsen til Udfoldning.
Det som af en Guddom rørte Sjælens indre Øie oplader sig nu til levende
at skue den under det jordiske Dække tilhyllede Væsenhed i Livet og
Naturen, og som Kjærlighedens vidtudstrakte Rige ligger hele Naturen
udbredt for det aandelige Blik.

Dog, i Elskovs første Vaagnen og Virken hviler den nye Verden, der skal
komme frem af Sjælens Indre, endnu kun i et stort Totalindtryk. De
første Øieblikke kunne vel virke tilskyndende; de Forelskte kunne komme
i megen ydre Bevægelse, vise sig oprømte, færdige til Spøg og
Munterhed; de ere rede til Alt, hvad der kan bringe dem den Elskedes
Person nær, de ere om sig paa alle Maader. Men letteligen skeer det
snart, at den elskende Sjæl, følende det sælsomt Dybe i Det, der bevæger
den, trækker sig tilbage i sig selv. Den begynder at føle Livets høie
Alvor; det er, som om noget Guddommeligt (θεῖόν τι) havde berørt den;
alt Personligt, al Begjær forstummer; Følelsernes vexlende Spil smelter
hen i en usigelig dyb og stille Længsel efter den fraværende, og en
ligesaa dyb og stille Glæde ved den nærværende Elskede. Endnu kun i en
i det Hele oplivet Existents, ei i en bestemt Attraa eller Opfordring,
vover Sjælen at aabne sig for den Elskede. Selv naar den føler sig
usigelig vis paa Gjenkjærlighed, er det dog, som om den frygtede for,
ved et bestemt Udbrud at forstyrres eller foruroliges i Alt, hvad der
endnu saa himmelsk skjønt staaer for den. Dog snart griber den
tiltagende fulde Glæde ved det andet Væsens Tilværelse vækkende og
tilskyndende ind i hele Personligheden og fremkalder den heftigste
Attraa og Stræben med Mod og Selvtillid. Derved kan da ogsaa efterhaanden
det Higende, Brændende træde ind i Elskoven. Men hvor skjøn staaer ikke
i og for Erindringen Elskovens første milde Vaagnen, dens rolige Dvælen
i dens indre Lyst, og den stille, dybe Længsel og Glæde, i hvilken en
anelsesfuld Forudfølelse af et stort endnu slumrende Liv stod i den
Elskendes Indre, og i hvilken et eneste Ord, det Ord, der udtalte sig i
den Elskedes hele Tilvær, for os overgik alle de mangfoldige Ord, hvori
vi senere hen higende forlystede os ved at høre vor Lykke udtalt.

I en uendelig Hengivenhed bryder Kjærligheden ud. Med mægtig Kraft
river den Mennesket udaf den blotte Jegheds snævre Kreds. Skulde end
Individet forhen nok saa meget have betragtet sig selv som sin
Tilværelses nærmeste Middelpunct, ført Alt tilbage til sig selv og
ladet Tanken paa sig selv ledsage al sin Modtagen og Stræben, saa vil
det dog snart, naar Elskov opfylder det, skue sit Livs Middelpunct i
den Elskede. Tanken paa denne, der nu synes overalt at omgive os i en
evig Nærhed, ledsager alle øvrige Tanker. Den Elskede bliver nu Den,
til hvem Jeget fører Alt tilbage, hvad det monne leve, om hvem al dets
Lyst og Attraa dreier sig, ved hvem det hænger med al sin Hu og alle
sine Følelser, saaledes som Mennesket skulde hænge ved Gud. Hvorledes
en ungdommelig Elsker føler, vide vi Alle. Det er for ham, som om den
Elskede er den Guddom, fra hvilken hans Liv gaaer ud og hvortil det
gaaer tilbage, ret som om Livets hele Fylde strømmede til ham fra denne
Kilde alene. I og ved det elskede Væsen føler han sig salig og synes
han sig at leve. Og ikke er det lutter Blændværk og Forvirring, naar
Følelsen lever i ham med denne Magt. Thi fra Livets sande Sol kommer
dog den Straale, der, bringende et nyt Liv, falder fra den Elskede i
Elskerens Sjæl, og hvad der i den Elskedes Væsen henrykker ham saa
usigeligen, er jo selve det Evige, som træder ham imøde, og derved
tillige vækker de dybeste Mindelser hos ham af det samme Eviges i hans
eget Væsen sig rørende Herlighed. Derfor har hans Elskov denne Magt
over ham, at han vilde føle sig som tilintetgjort, om han skulde
fortabe dens Gjenstand, som om han var nedstødt fra sin Himmel.

Dog fornemmer den lykkelige Elsker denne Hengivenhed og denne fremmede
Indflydelse, som nu lever og raaber i ham, ingenlunde som noget
Hildende, Indskrænkende eller Betingende. Hvad Elskeren føler, naar han
kaster al Livets Angest fra sig under det ham heelt opfyldende nye Liv,
og naar hans Elskov bliver ham en skjøn Tilflugt, en herlig Besiddelse,
til hvilken han, i hvad der saa monne treffe ham, trygt iler hen, fordi
han veed den, og i den det Skjønneste, han eier paa Jorden, sikker og
tilforladelig, --- dette er netop Frihedens ægte Følelse, fordi det er
Følelsen af en fra alle lavere Baand løsnet, over det Jordiske sig
hævende, ægte og væsentlig Tilværelse. I Elskov er Aand og Liv; men
Aanden gjør overalt fri. Inderlig at hænge ved et elsket Væsen og at
give sig ganske hen for samme og nyde Livet ikke i det egne Jegs snævre
Kreds, men i det andet Væsens Værd og Fryd og Lykke, dette er at tabe
Livet for at gjenvinde det dobbelt. Men saaledes attraaer Elskeren at
tabe sig i sin Elskede. Al Kjærlighed har denne Jegiskheden
tilintetgjørende og døbende Magt, naar man elsker af et fuldt, udeelt
Hjerte. Ingensinde føles det Jegiskes dybe Armod snarere.

Hvo der elsker, er af Hjertet ydmyg. Hvor lidet forekommer den
Elskende, i sig Følelsen af sin Mangelfuldhed, værdig til at elskes af
et tilbedet Væsen, værdig til den Lykke, han har funden, som om han var
udkaaret fremfor Utallige. Det er naturligt, at det ofte er et
lykkeligt elskende Væsen ganske utroligt, at en Lyksalighed, der,
henrykkende Phantasien, kun syntes hjemme i Digtningens Verden, nu
virkeligen er bleven dets egen, og at det lever midt i Besiddelsen af
Det, hvorefter et ungdommeligt Hjerte udstrækker sine inderligste
Længsler. Det er naturligt, at det stundom kan være for den Elskende,
som om det Alt kun var en Drøm. Og dog føles Elskoven atter som den
fuldeste, virkeligste Besiddelse, med inderlig Tilforladelighed, Tryghed
og Vished, ikke en saadan, der, støttende sig paa disse eller hine
Overveielser, blot vilde være en middelbar Vished, nei med en Vished
saa umiddelbar og total, som Elskoven selv.

Den dybe Beredvillighed til at give Alt hen for den Elskede og den dybe
Ydmyghedsfølelse, naar vi betænke den usigelige Herlighed, der, som en
reen Himmelens Gave, er aabenbaret paa os i vor Elskov, denne Følelse
svarer netop ganske til Menneskets indre Natur, ligesom ogsaa netop den
ligger til Grund i Elskovs sødeste Fornemmelser. Den afsondrede, paa
sig selv beroende Existents er vel vor første, og er paa Jorden endnu
saa meget fremherskende; men den ægte Tilvær er dog den, hvori et fra
en høiere Kilde til os nedstrømmende almeent Liv træder ind i os og
udfylder os. Heri ophører og forsvinder den Følelse af Endelighed, der
ligger i al Personlighedsfølelse hos det Væsen, der ei har Livets
Udspring og Kilde i sig selv. Individet føler sig da, som om det svandt
hen, eller ligesom døde hen og løstes op i det almene evige Liv, og
føler sig derfor i Modsætning til dette almene Liv som Lidet eller
Intet, da det føler sig som modtagende saa uendelig Meget. Dog forgaaer
heri ikke Personligheden, men gjennemtrænges og udfyldes, og hæves og
forklares tilmed. Derved er da denne usigelige Hengivenhed og Ydmyghed
saa inderlig dyb og sød; den er den fuldeste Følelse af den høiere
Indflydelses rene Tilstedeværelse i os, hvilken Intet hos os
modstræber. Netop naar i Ydmyghedens dybe Følelse Selvfølelsen
overvældes og synker hen, føle vi os inderlig rørte og qvægede, og føle
os hævede, just idet vi føle os i vor Lidenhed og Ringhed. Denne
Ydmyghedens søde Følelse ligger da nu ogsaa i Elskoven og hører
væsentligen med til denne, ja netop her er den saa indtrængende og sød,
fordi her ogsaa Sandsen og den jordiske Lyst bliver udfyldt og qvæget.
Dog ligger det i Kjønnenes Natur, at Qvinden uden Tvivl føler denne
hjertelige Beredvillighed til at hengive Alt og denne Ydmyghed med en
Heelhed og tilmed en Inderlighed, hvilken en Mand maaskee neppe
formaaer at fatte tilfulde, hvorfore da ogsaa Qvindens Elskov uden
Tvivl kan naae en langt høiere indre Styrke, end Mandens. Hvad Fabelen
om Tiresias udtaler, gjelder sikkert ogsaa i den høieste aandelige
Henseende. Just fordi Qvinden er den Underordnede, og tillige er Den,
hvis hele Hu og Stræben kan gaae op i den egteskabelige og huuslige
Existents, hvorimod Manden endnu søger en Virkekreds udenfor Huset og
tilmed vender sin Interesse til tvende Steder, just derfor maa Qvinden,
og kan ogsaa, give sig hen med en Absoluthed og en Heelhed, der maa
give hendes Følelser en saa usigelig Intensitet.

Den ene Elskende er for den anden som en Gud, saa aldeles synes sig den
ene opfyldt med nyt Liv og opløftet til et høiere Himmelens Rige blot
ved den andens herlige Nærhed, ved Det, denne yder os ved sin blotte
Existents. Men hvad den ene saa inderligen tilbeder i den anden, er dog
i sig selv den Guddom, som overalt, hvor Tvende ere samlede i dens
Navn, er tilstede midt imellem dem. Derfor kunne Begge føle sig saa
uendeligen ydmyge, saa meget de end ere for hinanden. Jo mere den
Elskende ophøies og prises, desto dybere føler han sig som Den, der
Intet er, og selv maa tage Det som en uendelig Himmelens Gave, at han
kan være sin Elskede saa Meget. Derfor fylder ogsaa Kjærlighed med den
sødeste Taknemmelighed, fordi vi som reen Gave modtage saa usigelig
Meget, en Taknemmelighed, der ei, saaledes som den ofte ellers kan
have, har noget Betingende ved sig, men der er saa total og uendelig,
som selve den Elskovens velgjørende Magt, der fremkalder den. Aandelige
Gaver ere overalt herligere, end jordiske, og ere desto mere
velgjørende, jo mere de uden disse eller hine Hensyn i deres fulde
Heelhed gaae ind i Modtagerens Sjæl. Men af Alt, hvad en ypperlig Sjæl
kan skjænke, er Intet saa velgjørende, som hvad den kan give
umiddelbart ved sig selv, udtrykkende det Høieste og Guddommeligste i
sin umiddelbare Tilvær. Men hvad den Elskede saaledes skjænker
Elskeren, nu endda ganske hengivende sig til, for ham især at være
dette Alt, det er af alle Gaver den, der vækker den inderligste
Taknemmelighed, saa lidet her end er Tanke paa en Forbindtlighed. Deraf
ogsaa det inderligt Rørende i Elskovs Følelser, hvilke stundom
fremkalde en reen Glædes Taarer. --- Og her viser det sig da ogsaa, at
det bekjendte Ord: „det er saligere at give, end at modtage“, ikke
gjelder overalt. Tvertimod turde vel Den, der kjender alle de Følelser,
som bevæge Menneskehjertet, ligesaa ofte finde Anledning til at sige,
at det er saligere at modtage end at give, hvilket ogsaa ganske stemmer
overeens med Menneskets inderste Natur og Væsen, ifølge hvilket den
totale Hensjunkenhed i et høiere Liv er den høieste Leven. Hvor Det,
der gives, ikkun er noget Jordiskt, kan det vel ofte være hildende for
Sindet at modtage; men ogsaa at give er da ikke sjeldent ligeledes
hildende for Sindet. Hvor derimod Det, der gives, er noget Høiere,
eller hvor overhovedet Noget gives paa en hjertelig Maade og af fuld og
sand Kjærlighed, og nu derfor et fuldt, udeelt Hjerte kan træde Gaven
imøde og aabne sig for den, der vil som oftest Den, der modtager, nyde
Livets høieste Fylde, nyde Kjærligheden, inderligere og dybere, end
Den, der giver. Allerede Det, at Den, der giver, maa virke og handle,
hvorimod Den, der modtager, udeelt kan fornemme og tilmed udeelt og
ganske give sig hen, gjør, at denne Sidste paa en mere total Maade kan
føle sig udfyldt og grebet. Men saliggjorte ere vi kun, naar vi ret ere
hentagne og grebne. Enhver kan blive vaer, at det Livsens Ord, vi
hørende henrives af, er mere besaligende, naar vi ganske henrives
deraf, end det samme Ord, naar vi give det til Andre, om det end gives
med sin fulde Kraft. Den Glæde, hvilken den i os vakte, undfangne og
fremspirende Idee fører med sig, er en ganske anden end den, hvilken
Ideen bringer, naar vi udtale den og vække den hos Andre. Ja selv da
træder Glæden ikkun ret tilfulde ind, naar vi, idet vi give, føle os
selv saa henrevne, at vi i denne Henrevenhed selv modtage en fra en
høiere Kilde til os kommende Livets Strøm. --- Og saaledes har da ogsaa
al den Kjærlighed, der ydes os i Elskovsforholdet, og hvilken vi med
inderlig Lyst modtage, noget Fiint og Sødt, Qvægende og Inderligt i
sig, til hvilket den Kjærlighed, vi yde, ei altid kommer.

Af Elskovens dybe Hengivenhed udfolder sig dens skjønne Tjenstfærdighed
og den Hylding, der især i Middelalderens Riddertid antog saa skjøn en
Skikkelse. Til den Elskedes Fødder nedlægger Elskeren enhver Priis, han
vinder, og tillægger hende Æren for enhver Seier. Og med Rette; thi hun
er Den, der besjælede ham med høit Liv, og fra hende kom det Mod og den
Kraft til ham, ved hvilken han seirede. En Muse er hun ham, der
beaander ham ved alle Granskninger, under alle Bestræbelser. Hvorimod
hendes Hu staaer til, istedetfor alle tilvundne Tegn paa Seier eller
udførte Værker, at lægge sig selv ned ved den Elskedes Fødder, og
hengive sig ganske til at være ham en tro Tjenerinde. Hvad det er for
en inderlig Glæde og Lyst at kunne opoffre for den Elskede, det føler
ethvert lykkeligt elskende Hjerte. Møie, Arbeide og Besværlighed blive
vort Hjerte inderligt kjære, naar vi have dem for Den, vi elske, ja en
Trang og Fornødenhed blive de os, for at vi kunne tilfredsstille og
ligesom svale vor henstræbende Attraa efter en total Hengivenhed. Kun
den sælsomste Forvirring i Forestillingerne og Følelserne, Følgen af en
til en ulyksalig høi Grad voxen Afvigelse fra Naturen og dens
betydningsfulde Tilsagn, kan, især hos Qvinden, have været Det, der har
kunnet bringe saa mange Elskende til at lade sig af Dem, som maae synes
at elske uden at kunne det, forlede til at troe, at de, formedelst en
over Andres naturlige Levemaade dem ophøiende Fornemhed, maatte give
Slip paa Elskovs skjønne Tjenstfærdighed, den søde Omhu for et elsket
Individs personligste Fornødenheder, og den Inderlighed, hvormed
Elskoven føles og nydes, naar vi for den arbeide, opoffre, lide. Elskov
er inderlig øm. Af dyb Omhyggelighed for, at den Elskede kan have det
vel, have Alt paa det Bedste, forsager den selv af Hjertet gjærne, ja er
endog fiin og varsom. Ømhed er dyb Omhu, der holder det Elskede i Hævd
og helligt og netop deri nyder sin Følelses Dybde og Sødme. --- Men
atter bevæges Hjertet igjen af inderlig Lyst til at give sig hen til
sød Nydelse for den Elskede, at denne kan føle sin Besiddelses fulde
Velbehagelighed, idet han selv med Ømhed gjør af sin Elskede. Deraf den
Elskovs Kjælenhed, der vil nyde ved at nydes.

Men som Elskov er en ubegrændset Hengivenhed, saa er den ogsaa en
ubegrændset Attraa. Ligesaa inderlig Elskeren taber sig i sin Elskede
og hænger ved denne, lige saa indtrængende staaer hans Hu til at møde
samme Hengivenhed og ganske at kunne tilegne sig det tilbedte Væsen og
tage det i Besiddelse. Begge Dele, Hengivenheden og Attraaen efter
udelukkende Forening, ligge i den af Elskovens Inderste fremgaaende
Drivt til Sammensmeltning og fuldkommen Eenhed. Som Trang og Længsel og
Begjær er Elskoven lige saa udelukkende og færdig til at fortrænge alt
Fremmedt, ja ligesaa uslukkelig, som dens sig selv hengivende og
hentabende Lyst er uudtømmelig. Baade som en Lyst til i sig at optage,
ja ligesom fortære den Elskede, saa og som en Lyst til at løses op og
forgaae, ja ligesom at døe hen i den Anden, yttrer sig Elskovens
guddommelige Afsindighed (μανία θεία) i tusinde endog mimiske Udbrud.
Ved denne udelukkende Attraa kan nu ogsaa det Heftige, Brændende,
Urolige, i Fryd og Qvide Vexlende, komme ind i Elskovens Følelser, og
den Elskende komme til at prøve, hvad Platon (Symposion, 403) skildrer
som Eros's Vilkaar: At han, født af Poros og Penia, Overflod og Armod,
som et Mellemvæsen mellem Guderne og de Dødelige, bærer Præg af sin
dobbelte Herkomst. Først --- saa hedder det om ham --- er Eros „altid
fattig og langt fra at være fiin og smuk, som de Fleste troe; derimod
er han raa, af forvildet Udseende, løber om uden Skoe og uden Hjem,
ligger paa den bare Jord uden Dække, sover ved Dørene og paa Veiene
under aaben Himmel; han har sin Moders Natur, boer altid sammen med
Trangen. Men atter efter Faderen efterstræber han det Gode og Skjønne,
er mandig, kjæk, drivtig, en drabelig Jæger, altid pønsende paa
Konstgreb, begjærlig efter Indsigt og fuld af Paafund; han
philosopherer hele Livet igjennem, er stærk i Trolddom og i at blande
Trylledrikke, en vældig Sophist; han er hverken som en Udødelig eller
som en Dødelig af Natur; men snart paa selvsamme Dag blomstrer og lever
han, naar det gaaer ham vel, snart døer han hen, men lever atter op
igjen, ifølge Faderens Natur; hvad han erhverver, rinder bestandig bort
igjen, saa at Eros hverken er rig eller fattig, og ligeledes staaer
midt imellem Viisdom og Uforstand.“

En dyb om Individets eget Jeg sig dreiende Følelse, en dyb jegisk
Følelse altsaa, gaaer i Elskoven i Eet med en dyb sympathisk. Allerede
i blot organisk Henseende er Driften i sit Udspring en individuel og
jegisk, thi den er en umiddelbar Fortvirkning af den til bestandig
Selvuddannelse og individuel Formation virkende organiske Tendents. Og
siden sættes ved denne Drivt ikke blot Sympathien i Bevægelse, idet den
drager det ene Individ til det andet, men dernæst ogsaa reiser sig
Følelsen af dybe sjælelige Savn. Her viser sig en inderlig stærk Trang,
som kræver Fyldestgjørelse. Bevæget af inderlig Glæde over det Skjønne
og Guddommelige, det Tiltrækkende og Qvægende, der skues i et andet
Individ, føler Mennesket sig snart som Den, hvis isolerede Liv kun er et
halvt, altsaa som Den, der, for at komme til fuldelig Existents, maa
voxe sammen med et andet Væsen, for nu i et uafbrudt Samliv med dette,
delende Alt med samme, at nyde Livet som en uafbrudt Elsken. Derfor
attraaer han en fast og inderlig og udelukkende Forening; ikke i denne
eller hiin Henseende alene, men i alle Henseender. Han begjærer nu af
al Magt, ikke blot at hengive sig selv, men ogsaa at tilegne sig den
Elskede og udelukkende at besidde denne; ei blot at elske, men og at
elskes igjen med samme Inderlighed.

Men hvad den Elskende attraaer, er, saa himmelsk end Elskovens Salighed
monne være, dog altid et jordiskt Gode. Thi det er et jordiskt Væsens
Gjenkjærlighed og en jordisk Forening. Det er overhovedet en Lykke, som
ei altid undes selv Den, der længes meest derefter; den kan ogsaa gives
og fortabes igjen. Hvor afhængig er Elskovs Lykke af jordiske
Sammenstød? Derfor kan det ved lykkelig Elskovs dybe Følelser opløste
og hensvundne Selv ikke blive ved at være fortabt i et høiere Livs
udeelte og ustandsede Henstrømmen. Jeghedsfølelsen kan ikke hvile. Paa
tusinde Maader ægges den op igjen, og under en idelig Stræben efter at
tilegne sig, holde fast og bevare det Tilegnede, eller forvisse sig om
den udeelte Bevaren af dets Besiddelse, optændes nu letteligen alt det
Hede og Bryndefulde i Sjælen, der ligger i al jegisk Existents, naar
Individet kommer til skarpt at fornemme, at det har at gjøre med en
muligen modstridende Omverden. Ved den store Opæggelighed
(Irritabilitet) hos det ved Elskov blødnede Væsen reiser sig da altfor
let midt under den himmelske Glæde en fra Skinsygens frygtelige Helvede
kommende Fornemmelse. Ved al jordisk Besiddelse er Fare. Tilmed er
Frygten tilrede med de fjendtligste Følelser. Allerede ved sin Tilvær
alene kan Omverdenen da bringe heftige Gemyter i lidenskabelig
Bevægelse; og endog om herfra intet Fiendtligt ægges op i Sjælen, kan
letteligen den Individualiteternes Kamp, som maa opstaae, naar tvende
allerede til en vis Grad fixerede Organisationer skulle smelte hen og
løsnes, for at begge kunne voxe sammen til et nyt begge forenigende
Liv, fremkalde urolig Gjæring og optænde Harme. Ja det kan maaskee være
godt, om denne Gjæring indtræder tidlig under Følelsernes endnu rigeste
Fylde og dybeste Inderlighed, naar Elskoven endnu er frodigst til at
fortrænge al Standsning og at udfolde sit indre Væsen til ydre
Skjønhed. Saaledes som Menneskets Vilkaar nu eengang ere paa Jorden,
kan den rene, i sin indre Lyst heelt hensjunkne Elskov neppe holde ud,
endog i de uskyldigste Hjerter, uden at de engang komme til at fornemme
og erkjende, at det skumle Mørkes glædefiendske Vælde ogsaa kan gribe
ind i den jordiske Lyksaligheds meest himmelske Sphære, naar Mennesket
ei vogter sit Hjerte. Men hvad et dybtfølende Gemyt lider, naar det
første Gang føler Muligheden af et Brud paa Samlivets rene Heelhed, er
sikkert ofte langt Mere, end Nogen troer. Dog, vi ville siden finde
Stedet til videre at oplyse Elskovs, om end stundom ganske naturlige,
dog atter igjen ofte saa høist unaturlige, Lidenskabelighed. Endnu ville
vi blive ved med at betragte Elskoven i den skjønne Skikkelse, hvori
den lever i tvende lykkeligt elskende Væseners Samliv.

Ægte Elskov knytter, naar den er en lykkelig Elskov, Hjerte til Hjerte,
og gjør af tvende Væsener eet. Jo mere den voxer og vældig griber om
sig i begges Indre, desto mere stiger Samlivet til den sødeste Enighed.
Med usigelig Tryghed og Tilforladelighed hænge begge ved hinanden. En
dyb, Hjertet inderlig qvægende Meddelelse og Fortrolighed er imellem
dem. Den Enes hele Liv og Indre ligger aabent for den Andens Øine med
den reneste Sandhed. Allerede saasnart vi begynde at elske og haabe
Gjenkjærlighed, føle vi den inderligste Attraa til at udfolde Alt, hvad
der lever i os, for den Elskende uden Tilbagehold eller Forbeholdenhed.
Fordi vi føle alt Himmelskt og Skjønt som noget os Beslægtet og
Paarørende, føler Hjertet sig draget hen dertil, og længes efter at
aabne sig ganske for samme, trukket af en dyb Sympathie, bevæget af en
inderlig Tillid. Usigelig qvægende og velgjørende er Elskovs dybe
Fortrolighed. En sand frigjørende Kraft er der i denne rene Udstrømmen
af vor inderste Tilværelse, denne Opløsning af ethvert Baand, der ellers
saa ofte trækker os tilbage i os selv. Saa piinligt det er,
tilbageholden at indeslutte sig i sig selv, saa dybt vederqvægelig er
for Sind og Hjerte denne frie Aabenhed, denne udeelte Henstrømmen af
alle vore Følelser. Det er ikke godt for Mennesket, at han er alene:
saaledes løde Herrens Ord allerede om det første Menneske.

Hvor meget trænger ikke Mennesket til at meddele sig og aabne sit
Hjerte for en Anden, baade for at han kan besidde og nyde sig selv og
sit indre aandelige Eie med dobbelt Tryghed og Lyst, naar Det, der
levende opfylder ham, meddeelt til en Anden, vender bekræftet klarere
og fuldeligere tilbage til ham selv, saa og for at han kan befrie sig
for Meget, hvori han ellers lettelig hilder sig selv, at han kan lette
sit Hjerte, forvandle den indre Samtale til en ydre, komme ud af mangen
indre Uenighed eller Ufred, hvormed han ellers maaskee vilde søndergnide
sig selv, og kaste et klarere Blik paa sit eget Liv, idet han udfolder
det for en Andens roligere Øie? Men allersødest er denne Fortrolighed
imellem Elskende, som Dem, der ere saaledes forbundne, at hvad der med
Glæde eller Smerte treffer og bevæger den ene, det maa den anden føle,
som om det traf hans egen Tilværelses kjæreste Deel. Derfor længes
Enhver efter at bringe til sin Elskede, hvad der levende opfylder ham
selv, kun da i Stand til ret at nyde det, om det er glædeligt, eller at
bringe det paa det Rene hos sig, om det er foruroligende, naar han
deler det med Den, hans Sjæl elsker. Inderligen bliver Mennesket i
denne skjønne Fortrolighed fortrolig med sig selv, idet han længes efter
at fremlægge sit Indre reent og heelt for et elsket Væsens Beskuelse.
Og mere end en blot middelbar Meddelelse igjennem Ord er Det, som her
er. Umiddelbart, som Begge leve, ligger den Enes Liv og hele Sjæl aaben
for den Anden. Elskovs Fortrolighed er ikke en Fortrolighed i Dette
eller Hiint; det er en grændseløs Fortrolighed, hvilken Ord ei kunne
udtømme, som om Begge havde en uendelig Hemmelighed at betroe hinanden.
Den er Udtrykket af den idelig friske Stræben efter at give den Anden
sit hele Hjerte. Saaledes er i fuldeste Grad Sandhed mellem Begge, da
de inderligen glæde sig ved at vandre og leve for hinandens Øine.

Idet Mennesket nyder sit Liv i et andet Væsen og atter lever Livet for
dette, føler han sig selv ligesom gjentaget eller atter sat i det andet
Væsen. En Eenhed i en Tohed udgjør Grundcharakteren i Menneskets
Person. Den ligger i Bevidstheden. I hvad vi saa leve, leve vi det,
idet vi have med os selv at gjøre og henføre det til os selv.
Bevidsthedens og Personlighedens dybeste Udtryk, Samvittigheden, er
netop en Kjærlighed, hvormed vi have og besidde og omfatte os selv,
inderligen holdende vort sande Væsen fast og bevarende det ufortabt.
Men ogsaa i denne Henseende erholde vi det Hjerte, vi nedlægge i
Elskovs Helligdom, renere og klarere tilbage. Hvem ei hans ægte sande
Vel, en renere indre Væren, af Hjertet var bleven magtpaaliggende
forhen, han maa komme til denne Omhu for sig selv ved Elskovens Magt,
naar den er dyb og ægte. Hvad der var ham en dyb Mindelse om ideligen
at vogte sit Hjerte og bevare sig reen og fri, det bliver ham det nu
med dobbelt Styrke. Ligesom Intet kan være ham mere magtpaaliggende hos
den Elskede, end Hjertets Skyldfrihed og Reenhed, fordi kun den Elskedes
sande Værd kan give Elskoven selv Sandhed og Bestaaen, saaledes maa nu,
naar han elsker ægte, hans eget Væsens Skyldfrihed blive ham usigelig
hellig, da et elsket Væsen skuer i ham sin Guds herlige Billede. Han
vilde ellers selv forderve sin Elskovs rene Skjønhed. Samvittighedens
Uro vilde nu ikke blot være Følelsen af den egne indre Freds
Sønderbrydelse, men vilde nu ogsaa blive en Følelse af, at Elskovens
rene Fred var sønderrevet. Saaledes fordobler Elskoven Kraften i
Samvittighedens indre Stemme. Mennesket er viet til sin Samvittighed,
lader Goethe etsteds Nogen sige, da Talen er om Egtefolk, der, om de end
synes at leve til Besvær for hinanden, dog saa lidet skulle skilles ad,
som Mennesket skal skilles fra sin Samvittighed. Men vi ville ei see
Dette saaledes, som om Egteskabet derved fik et mørkere Udseende, men
tvertimod føle Samvittighedens inderlige Høihed og Skjønhed ved at
tænke os den, som den troe i Liv og Død med os Forbundne, der skal være
Det for Enhver, hvad paa den skjønneste Maade hans Elskede er ham, Den,
for hvis Skyld han med dyb Hengivenhed og Inderlighed føler sig dreven
til at bevare sig skyldfri, paa det han altid kan vandre reen for sin
Kjærligheds Aasyn. Og saa opløfter da Elskov Menneskets inderste Selv,
og just idet han giver sig selv ganske hen, gjenvinder han sig selv
dobbelt; just idet han afspeiler sig selv, som han lever i og for et
elsket Væsens Anskuen, føler han sin egen Existentses Værd og Betydning
dobbelt, i det samme Øieblik, da netop derved tillige Ydmyghedens dybe
Følelse fremkaldes. Thi den ene af disse Følelser ophæver ei den anden.
Begges levende, i den hele Tilvær indvirkende Tilstedeværelse og
Kraftighed hører til fuldelig og kraftig Tilværelse. Og hvor fri gjør
ei Elskovs dybe Følelse den maaskee ellers hildede og bundne Sjæl, naar
Mennesket deri har en tilforladelig Tilflugt under Livets Ængstelser
eller Angreb? Kjærlighed giver Kraft baade til at virke og til at
udholde.

En hellig Tilflugt saavel under Livets ydre Tilskikkelser, som under de
Ængstelser og Sorger hvormed et Gemyt kan have at gjøre i sit eget Indre
under sin Kamp efter en sand Existents, er et reent med inderlig
Kjærlighed os modtagende Menneskehjerte: det skjønneste af alle jordiske
Templer, i hvilket Mennesket træder den Evige nærmere, og kommer til at
føle Livet let og frit, idet han giver alle Byrder hen til Den, der kan
og vil modtage alle Byrder, som blive os for tunge, naar de med inderlig
Tillid nedlægges hos ham. Den skjønneste Helligdom, hvori det betrængte
Hjerte kan finde Ro, det sønderrevne Fred, og det betyngede kan udøse
hvad der trykker det, er et elsket Væsens elskende Sjæl, for hvis
beskuende Øie den Elskede efter sin indre Heelhed staaer renere og
klarere, end for sig selv i sine urolige Øieblikke. Hvor meget og hvor
ofte den ene Elskede bliver den Eviges Repræsentant og Præst eller
Præstinde for den anden, erfarer Enhver, som elsker, men dobbelt føltes
det i en Tid, da efterhaanden næsten alle Templer lukkedes eller tabte
deres Hellighed, og hver Enkelt mere og mere saae sig kastet tilbage
paa sig selv, og saae sig uden et til hans indre Hu svarende Samfund,
fordi Aanden, saa meget den end levede i Enkelte, dog var veget fra
Menighederne. Derfor pristes i denne Tid med dobbelt Styrke Herligheden
i det Livsforhold, der endnu altid maatte blive ved at være poetisk og
at erkjendes i sin Skjønhed, fordi Naturen her taler saa stærkt til
ethvert Menneskehjerte. Ja det kom undertiden saa vidt, at Helligdommens
lovprisede Vogtere tilbades som selve den Guddom, hvis hellige Ild de
opflammede og nærede i Hjerterne.

Men saa høit ægte Elskov løfter os op over det Jordiskes Trængsel og
Livets lavere Sorger og Idrætter, saa nær den bringer os det Evige, ja
lader os føle den høieste guddommelige Skjønheds umiddelbare
Tilstedeværelse paa Jorden og aabner for os den Eviges Himmel, saa
usigeligen dybe de Hengivenhedens Følelser og den Hjerternes
Fortrolighed er, som den fremkalder, saa aandeligt og reent overhovedet
det skjønne Samqvem er mellem Elskende, --- saa slaaer Elskov dog sine
Rødder ret dybt ned i den Jordbund, hvorpaa den hviler, som paa sin
Grundvold, medens den i sine Blomster er ætherisk og dufter op mod
Himmelen.

Og den viser derved sin indre Dygtighed. Thi den fører til et Egteskab,
en Familie grundfæstes ved den og et Huus, hvoraf en ny Slægt skal gaae
frem paa Jorden. Ei aandeligt alene er her Forholdet, men ogsaa ret
sandseligt, og just derfor ret sundt, forenende begge Livets Sphærer
til skjøn og kraftig Eenhed. Vel lever Mennesket ikke af Brød alene,
men det lever dog ogsaa af Brød, og kan ikke leve sundt uden samme. Det
er hvert endeligt Væsens Bestemmelse, der aldrig maa forglemmes, at det
ei blot skal tragte efter at være salig i sin Gud, men ogsaa udfylde
sin Plads i sin Verden og paa den bestemteste Maade voxe ind i denne.
Ellers vanker det veiløst om, kastet ustadigt frem og tilbage i Livets
Mangfoldighed. Kun den nøiagtigst bestemte jordiske Existents giver
Livet fuldkommen Holdning. Thi her er Mennesket en i et Hele indordnet
Deel og maa føle sig som saadan, ja vil desto fuldeligere blive
deelagtig i det Liv, der lever i Alle, jo fastere han griber sin Plads
og knytter sig til denne. Træet kan ei hæve sin Top mod Himmelen uden
at slaae sine Rødder ned i Jorden; hvorfra fik det ellers Næring? Og
Mennesket, der, som man treffende har sagt, bærer sin Rod inde i sig,
maa dyrke Marken og nære sig af dens Frugter, og er derved ligesaa
fuldt indgroet i Jordelivet, som Træet er det i Jordens Skjød. Saaledes
har da ogsaa Elskoven sit sandselige Rodfæste, og kan netop derfor holde
ud et heelt Liv igjennem. Ja dens indre Skjønhed og Kraft viser sig
desto fuldeligere, jo mere derved Livets lavere Syslen og Attraa faaer
dyb Betydning for de Elskende, idet de netop derved dagligen mere og
mere voxe sammen og nyde det skjønne Samqvem med hinanden og prøve og
forfare Elskovens sig altid mere udfoldende Fylde og Rigdom, naar den
ved tusinde Leiligheder gjør Livet kraftigt og sødt. I mangfoldige
Henseender, fra de høieste aandelige til de laveste sandselige, føler
Mennesket, at han ikke er bestemt til at forblive ene.

Den Naturdrivt, der under Elskovs Følelser er i Bevægelse og giver dem
deres eiendommelige Præg, ligger, selv hvor Elskovs Lyst og Længsel har
forklaret sig til den høieste Aandelighed, endnu bestandig til Grund og
spiller med ind i Følelsernes Hele med ikke ringe Styrke, og det er
ikke sjeldent øiensynligt, hvor meget den bidrager til, at tvende for
et heelt Liv forbundne Mennesker, under saa megen Anledning til
Opirring, holdes sammen, og at ethvert Brud saa let jevnes ud igjen.

Vi ville her nærmere overveie, hvad Kjønnets Dobbelthed, og den Drivt,
der knytter begge Kjøn sammen, har at betyde; dog at vi, istedetfor alle
blot physiologiske Betragtninger, alene ville fremkalde den, til hvilken
der dog altid tilsidst maa stiges, den, hvortil et Blik over hele
Naturen fører.

Igjennem hele Naturen gaaer den organiserende Stræben ud paa, at mere
og mere Det, der er det rette Væsen, det indre Aandelige, kan komme til
at træde frem, yttre og vise sig som saadant. Men dertil kommer det
desto mere, jo mere fremtrædende og bestemte (jo mere deciderede)
Sondringerne og Modsætningerne ere i den Mangfoldighed, hvori det indre
Væsen træder frem. Thi desto mere viser dette sig i sin indre Eenhed,
som enigende Princip, tilmed som besjælende og kraftigt, som ideelt og
aandeligt. Saaledes gaaer da i Naturens organiserende Stræben en idelig
sondrende og individualiserende Stræben igjennem det Hele, og hvad der
aldeles hører sammen, træder paa de bestemteste Maader ud fra hinanden.
Naturen sondrer og individualiserer overalt, men for desto fuldeligere
og paa desto høiere og aandeligere Maader at forene, jo bestemtere den
har sondret, paa det at Livets indre Ene, dets indre aandelige Væsen,
desto mere kan træde frem som fri Leven i en rig Fylde. Og nu gaaer da
den individualiserende Sondren endog saa vidt, at, under Organisationens
Stigen, Hovedmodsætningerne i enhver Art træde frem i tvende for sig
bestaaende, men aldeles sammenhørende Individer, i de tvende Kjøn,
hvoraf ethvert kun udtrykker sin Arts Natur halvt, og nu til fuldelig
Existents trænger til det andet. Individerne vise sig her da som de
spredte Lemmer af den store Naturens Organismus, i hvilken den indre
ene Natur under idelige Sondringer yderliggjør sig; de vise sig her som
til hinanden hørende Dele af en høiere Heelhed, men en Heelhed, hvilken
Menneskene ved fri Selvbestemmelse skulle bringe til Virkelighed.
Derfor føle da ogsaa Menneskene sig som saadanne til hinanden hørende
Lemmer, de føle, at den isolerede for sig selv bestaaende Existents er
en utilstrækkelig, og føle sig drevne til, forenende sig med Andre,
under Meddelelser med dem, at udfylde Savnene, ophæve Mangelfuldheden i
deres egen Existents, og komme til at nyde en mere total
Mennesketilvær. Det første, det nærmeste og naturligste Skridt hertil,
det, hvortil Naturen stærkest tilskynder, er nu det, hvorved den
Modsætning, der af alle Modsætninger i Menneskeexistentsen er den
aabenbareste, Kjønnets Modsætning, løses op i en høiere Eenhed, hvori
dog Personligheden ei gaaer til Grunde, men netop forhøies og fremmes
til sin Fuldendelse. Naar Enhver har funden sin Halvdeel, er Enhver
først et fuldendt Menneske.

Begges Sammensmeltning til en saadan Eenhed skeer ved en Forbindelse,
hvorved der fremkommer et nyt, et tredie Individ, der skal være, og
ogsaa mere eller mindre er, et Udtryk af begges sammensmeltede
Individualiteter. Men det nye Individ hører dog nu altid kun til eet af
Kjønnene, og den sluttede Kjæde aabnes igjen; thi ogsaa det tredie
Individ kommer til at trænge til et fjerde, og Indifferentieringen skeer
ikkun ved en Række uden Ende, paa det at Livet ei skal gaae istaae, men
at, tilligemed den individuelle Relativitet og Trang, Kjærlighedens
Fyldestgjørelser ideligen kunne vende tilbage, og Kjærligheden ideligen
forene det Sondrede. Hvorledes Aandeligheden her naaer et høiere Trin i
Naturen, er klart. Allerede i Dyret kommer her en høi Sensibilitet mere
og mere frem, og et dybt Instinct viser sig. Men Sensibilitet og
Instinct ere i Dyret det Sjæleligste, Det, der meest nærme sig til
Aandelighed. Instinctet er en universel Naturs i Dyret talende Stemme.
Men hos Mennesket kommer her det til Yttring, der overalt er det høieste
Aandelige, det høieste Enigende, Kjærligheden.

De tvende ved Elskovs Styrke til hinanden dragne Individers fuldkomne
Forening i Kjærlighed og et af denne Enigelse udspringende Samliv, ved
hvilket Livet stiger i Kraft, Betydning og Fylde for Enhver af dem:
Dette er, endog med Hensyn til Naturens Formaal og i dens egen Økonomie,
altid det Første og Fornemste; Dette er her Hovedsagen. Affødningen
derimod er kun det Andet og Secundære. Alene for den tilstundende Slægts
Skyld kan den nærværende ei være uddannet til Modenhed. Den tilstundende
bliver jo dog selv engang en nærværende; og skulde da atter den
fornemmeligen være til for den næstes Skyld? Nei, det Nærmeste er her
ogsaa det Første. Nutiden hører os til. For selv i sin egen Tilvær at
udtrykke den Livets Fylde, til hvilken da siden ogsaa den næste Slægt
skal komme, for selv at deelagtiggjøres i denne Fylde og nyde den, og nu
tillige, under Jordelivets mangfoldige Stræben, at modnes frem til en
Evighed, derfor lever nærmest den nuværende Slægt. Ei heller gaaer hiin
Tendents til Affødning frem af nogen anden Stræben end Individets
Stræben efter selv at existere tilfulde. Dens Udspring var jo netop den
individuelle Uddannelsesdrivt. Ja saa samstemmende med sig selv er her
Naturen, og saa rigt er Livet, at alt Det, hvorved den næste Slægt
fremkaldes og fremelskes, netop i høieste Grad tjener til at give Livet
Gehalt og Kraft, Skjønhed og Tillokkelse for de nærværende Individer.
Paa Organisationens lavere Trin see vi saa ofte Individets skjønneste
Blomstretid, den individuelle Uddannelses og Fyldes høieste Punct, gaae
umiddelbart foran for det Tidspunct, ja falde sammen med det Tidspunct,
da Spirerne til de nye Individer dannes, men, naar disse ere lagte, maa
nu ogsaa ofte den skjønne Skikkelse falde hen eller blomstre af. Hos
Mennesket er det derimod ingenlunde saa, end ei i organisk og legemlig
Henseende, endsige i aandelig. Mand og Qvinde staae i den frodigste
Væxt og Fylde, ja opnaae den først, medens den næste Slægt spirer op
under deres Øine og til deres Lyst. Og medens Faderen virker og
tilveiebringer og samler, medens Moderen vaager og sysler og bevarer for
Børnene, ja medens Begge endog opoffre og fornægte sig Meget for disse
Børn, komme Begge først til at nyde deres eget Liv, deres egen Samlen og
Virken og Syslen og Livets mangfoldige Eie tilfulde, og de fremelske hos
sig selv et høiere og herligere Liv, medens de fremelske --- en egentlig
Fremelsken, en Elsken altsaa, være det altid --- det samme Livs første
Rørelse hos de Smaae. Hengivende sig til at leve, ei for det egne Selv,
men for en elsket dem tilhørende Kreds, leve de nu først retteligen.
Livet er overalt desto mere en sand og ægte Leven, jo mere det er en
sand og ægte Elsken. Men ved Affødningen udvides den skjønneste
Kjærligheds Sphære.

Saa agte da, naar under Elskovs Følelser nu ogsaa Naturdrivten er i
Bevægelse, Ingen Det for uværdigt eller ringe, hvad der i Naturen har
den dybeste Betydning. Dog behøves det at sige Dette? Kun i en
Tidsalder, i hvilken, fordi Religionens og Poesiens fulde Kraft ei mere
levede i Livets vigtigere og store Forhold, en sælsom Sentimentalitet ei
sjeldent bragte følende Sjæle til, nu overalt at ville smage noget
Helligt og Poetisk ud af Gjenstandene og alene leve af dette overalt i
alle Forhold, kun i en saadan Tidsalder kunde det behøves at sige, at
man ei skulde lade haant ad Elskovens dygtige Naturlighed. Der kunne vel
gives særegne Forhold, hvori det kan være en ret naturlig og derfor
kraftig Beslutning af Elskende, alene at leve aandeligen knyttede til
hinanden. Men kun i en Tid, hvori Sjælene paa de sælsomste Veie søgte
Erstatning for det fortabte Samfund med Gud, uden hvilket dog Ingen kan
holde Livet ud, var det at befrygte, at sygelig Følsomhed havde bragt og
kunde bringe Denne eller Hiin til at holde sig for god til et sundt,
naturligt Liv. Dog, selv et sundt og reent Gemyt kan vel, især under
Ungdommens fyrige Enthusiasmus for det Store og Høie i Menneskelivet,
undertiden komme i Forvirring ved at tænke paa, hvorledes de dybeste og
helligste Følelser staae saa nær ved de sandseligste. Men Ingen see dog
skjævt til Elskoven, fordi den er saa total, fordi dens Kraft strækker
sig fra det Høieste til det Laveste i Menneskets Tilvær, og ikke udstøder
Noget i hans Natur. Just derved er den saa dygtig og herlig, at den let
forener og gjennemtrænger til Heelhed Det, hvad der ellers saa let
falder fra hinanden i Menneskets Liv, og hvad der saa let gjør, at han
lever det Liv i Momenter og som Stykkeværk, der skulde leves heelt og
som af eet Stykke. Thi til de Ting, der forvirre Menneskene meget og
vanskeligen bringes paa det Rene, hører Dette, hvorledes de skulle føre
det høiere Liv, der gaaer op for dem i Ideernes Fylde i deres Indre, ind
i det daglige jordiske Liv, og midt i dette gjenfinde og endnu besidde
og have hiint. Elskoven opløser Opgaven letteligen; thi den kaster sin
himmelske Glands paa hele Jordelivet. Det Laveste og Ringeste i det
daglige Jordeliv bliver skjønt og qvægende for Hjertet, fordi det hører
til det Samliv, hvori de Elskende leve for og med hinanden. Kjærligheden
helliggjør hvad den berører. En ny og dyb Længsel tragter efter at faae
sin Gjenstand, og faaer den denne, bliver Livet i en ny Sphære paa Ny
en uophørlig Sørgen og Virken for elskede Væsener, en ny Kjærlighedens
Tjeneste. Derfor mistyde Ingen den Styrke, hvormed Længselen ofte virker
og Savnet stærkt forkynder sig i de ædleste qvindelige Væseners Indre.
Ei Længselens Styrke gjør her en Forskjel, men Maaden, hvorpaa den
fremstiller sig i Sindets klarere Følelser, og hvad her den bevidste Hu
stunder hen til. Hvad der i ædle Qvindehjerter kommer til bestemt
Fornemmelse og Bevidsthed, og det endda kun seent, det er netop
Længselen efter Naturens skjønne Kald til Qvinden, Længselen efter den
Glæde, der vindes gjennem Smerter og næres i søvnløse Nætter, Glæden
over, at et Menneske er født til Verden, som Christi Ord lyde Joh.\ 16,
21.

Men vistnok maa Naturen her være naturlig. Den er det aldrig, naar den
ei er reen og sædelig; thi den er det ei, naar den i vilde Udbrud
sønderbryder sin egen rolige Bestaaen og indre Fred. Al ægte Elskov er
af Hjertet reen og derfor ogsaa paalidelig, naar den kun er i sig selv
nogenlunde klar. Men det er den ei altid, og saaledes kan der være
Fare, hvor der mindst skulde være den. Dog, da det ellers Naturlige saa
er saa lidet naturligt og lidet menneskeligt, er ogsaa, naar Elskov kun
ikke er alt for ubesindig eller uægte, den selv i høi Grad værnende. I
alle gode Øieblikke føles Elskov som en uendelig Omhu og Baretægt, og i
dens fulde, skjønne Henrykkelse løser enhver Begjær sig op i den
aandeligste. Saadanne Øieblikkes Kraft maa da formedelst indre
Betænkning og Bekræftelse gjennemtrængende udbrede sig til de øvrige. Jo
mere Elskovs dybe Følelser komme til Klarhed, jo fuldeligere den Attraa
er, hvormed Elskeren tilegner sig sin Elskovs hele Fylde og paa alle
Maader knytter sig til den og lever i den, desto større vil Elskovs
indre Reenhed være. Den bevare da ogsaa selv under den nøieste
Forbindelses retlige Frihed denne Reenhed, der holder al ἀκολασία
borte! Thi her er endog da en Grændse, det for Elskovs egen Tryghed og
Ro er farligt at overskride.

Men endnu udspringer af en sædelig Grund noget Finere og Ædlere, der i
dyb Elskov viser sig skjønt, og altid maa holdes i Hævd. Det er dens
Blyhed og Baersomhed, Agtelsens ømmeste Yttring, en vis hellig, til
Naturmodsætningen svarende, selv i den inderligste Nærhed og nøieste
Forbindelse vedblivende Fjærnhed, der holder de Forbundne ude fra
hinanden, paa det at ikke den selvegne Bestaaen og tilmed selve Elskoven
skal gaae til Grunde under den stærke Tiltrækning, og Begge i vild
Henrykkelse fortære hinanden. Alt, hvad der berører os som helligt, og
hvori ligesom en Guddom kommer os nær, har noget Standsende i sig, der
forrasker og tæmmer Sindets frie Lyst. Deels tør Jeget, saa meget det
end tiltrækkes, ikke gribe til, fordi det føler sig som lidet og ringe,
og nu frygter for at støbes tilbage; deels stræber det ogsaa selv imod;
for ei at overvældes, indtil det endeligen dog med ganske Hjerte giver
sig hen og nu gjenvinder sig ved villigen og gjærne at fortabe sig. Ja,
jo stærkere Modstanden var, eller kunde have været, i første Øieblik, jo
stærkere overhovedet Personens Selvegenhed til at holde sig fast, desto
stærkere vil Følelsens Intensitet vorde, naar nu dog Hjertet maa give
sig ganske hen. Men i den indtrædende Vexelstrid imellem Personligheden,
der vil haandhæve den selvegne Bestaaen, og det mødende Guddommeliges
overvægtige Tiltrækkende, maa ofte Hjertet først lukke sig end mere, for
at kunne sunde sig, til at aabne sig ret fuldkommen for Sympathie og
Kjærlighed. Det støber da selv tilbage for at holde sig, imedens det
føler sig forrasket og betaget, og derved føler sig forlegent og ufrit.
Derfor følger, med Blyheden i Omgang imellem begge Kjøn, i Ungdommen
ofte tillige en vis Skyhed og Strenghed, ja en vis indre Stoltheds
Følelse, hvortil Sindet ligesom tyer hen, naar det Tidspunct nærmer sig,
da Hjertet skal føle sig sig hendraget saa stærkt til et andet Hjerte,
at det skal tabe sig selv til dette. Forbærvelig bliver denne Stolthed
kun, naar Hjertet vedvarende vil slutte sig ind i sig selv og lukke sig
for Elskovs skjønne Sympathie, og vil rodfæste sig i sin selvegne
Bestaaen. Da stræber den mod Naturen og mod Livets indre Væsen, der kun
sondrer for desto fuldere at forene. Men selv i den inderligste Enigelse
maa Sondringen og Modsætningen ei gaae til Grunde. Elskov er en Eenhed i
Tohed, og for at den kan leve med Kraft, maae baade Eenheden og Toheden
træde kraftigen frem. Derfor blive og bevare sig baade Selvegenheden og
hiin hellige Fjærnhed under Elskovs Hengivenhed og Sympathie! Selv under
dennes dybeste Magt vil dog et kraftigt Sind haandhæve sin eiendommelige
Natur og sin Personlighed. Denne føle netop den Elskende luttret og
hævet ved sin Kjærlighed! Ja Elskovs Inderlighed er, naar der kun elskes
i Sandhed, saa langt fra at lide under den nu muligen indtrædende
Modstand, at netop ofte Elskov blev mat og lunken, fordi en vis Modstand
eller Modvirkning, en vis Resistance, fattedes. Men tilligemed bevare
sig da netop derfor hiin ømme, endog blye Agtsomhed, der vil holde det
elskede Væsens indre Personlighed i Hævd og Agt, ja holde den, som en
Guddom, hellig. Der er, saa dygtigt end og frit for Tilbagehold og
Blyhed ellers det sunde Forhold viser sig, dog i al Elskov Noget,
hvilket, som den spæde Blomst, fordrer den ømmeste Berøring.

Ogsaa udadtil yttrer sig hiin varsomme Fiinhed, saavelsom hiin i sig
selv sig tilbagetrækkende Skyhed, i Omgangen mellem Elskende, naar ædel
Elskov er sig nogenlunde klar. Jo inderligere dens hellige
Guddommelighed fornemmes, desto mere vil et fiintfølende Hjerte skye
alle saadanne lydelige Udbrud eller ufordulgte Yttringer af sin Elskov,
hvilke kunde profanere den ved at give den blot for den ligegyldige
Tilskuer, eller ved alt for meget at vise den Elskede i Drivtens Magt;
ligesom den finere Følelse ogsaa vil lede til, overalt at være
opmærksom for den Agtelse, der skyldes enhver Tredie, ja for den Uret,
det kan være mod mangt et efter Elskovs søde Glæde længselfuldt Gemyt,
at ægge dets Længsler op ved sin Lykke. Alt Dybt og Helligt overhovedet
vil et ædelt Gemyt, om det end vel lader det træde frem igjennem Ordet,
dog kun sjeldent lade tale i umiddelbare Følelser midt i det daglige
Livs Samqvem, hvor de forskjelligste Tanker og Tilbøieligheder bevæge de
forskjellige Mennesker. Igjennem Ordet omtales det kun, i Følelsen er
det selv tilstede. Og det skal det ei være, hvor det ei kan være det
tilfulde, og kan bevæge Alle med samme Følelser. Ellers gives det
letteligen til Priis paa en Maade, der enten støder eller ægger til
Spot. Der skal megen Yndighed til, for ei at støde an mod al Ynde, naar
Mennesket rives hen af heftige Følelsers overvældende Udbrud. Dog, det
ere just ei saadanne Betænkninger, det er den umiddelbare Følelse af det
Hellige i en dyb Kjærlighed, som gjør rene og ædle Sjæle blye og
tilbageholdne med deres Elskov i selskabelige Samqvem.

Men hvor Elskende ere samlede, kunne de ei letteligen være fra
hinanden. Idelig lystes de at være om hinanden og give sig af med
hinanden. Hvilken Udvei finder da Følelsen der, hvor Hjertets dybere
Bevægelser skye at sees? I tusindfoldig Spøg og Skjemt tilhyller sig
Følelsens Inderlighed. Men medens de Elskende nu gjækkes og synes at
slaae sig selv løs fra deres Kjærligheds Overmagt, nyde de netop det
søde Samqvem med hinanden.

Dog er ogsaa ellers i det Hele Elskovs muntre, ofte saa tækkelige
Overgivenhed lige saa naturlig, som den er snar ved Haanden. Ja netop,
naar i dens fyrige Ungdom Følelsen af dens dybe Alvor er stærkest, er
just Sindet lettest færdigt til, i Øieblikke ligesom at rive sig løs fra
den, rørende sig friskt i de ubundneste Bevægelser, som om det lystedes
at drive et løst Spil med Hjertets bedste Eie og at løse Følelserne
aldeles op. Det er da, som om det endelige Væsen ei kan udholde saa
megen Fylde, men stundom ligesom maa løse sig fra den, for at kunne bære
den. Ja saavidt gaaer denne Lyst til under Spøg at nyde og inderligen
gaae ind i sin Elskov, at det endog er hyppigt og naturligt, at en
Elskende forlyster sig ved at have Noget at belee og at finde komiskt
hos den Elskede, hvilket Tiek i Phantasus endog mener aldeles at høre
med til Elskov. Naturlig er denne Lyst til at finde Noget at lee over
hos den Elskede; thi naturligt er det at sysselsætte sig med den
Elskedes Eiendommelighed indtil dens individuelleste Særegenheder og
idelig at vende tilbage til og dvæle ved dens Betragtning. Men for Den,
der med dette levende Øie, med denne Interesse, beskuer en
Individualitet, er der noget Opæggende og Irriterende i den, fordi der i
enhver bestemt Individualitet er noget Særegent og Charakteristiskt,
anderledes, end hos andre, og det Noget, der, fordi det hører med til
denne Individualitet, træder op som noget Absolut, men nu derfor ligesom
støder og contrasterer besmere, og derved irriterer til at gjækkes
dermed.% [FLAG: "contrasterer besmere" — the word "besmere" is verified from the scan but reads oddly here; verify against another copy]
 Overalt i Venskab og Kjærlighed føle vi just Lyst til hos de
kjæreste Personer at finde noget Sælsomt og Pudseerligt; og det være
Ingen ukjært at være Gjenstand for en Latter, hvormed kun Den belees,
der elskes. Endog Det, der i og for sig er alvorligt og skikket til at
fremkalde en dyb Agtelse, kan blive Gjenstand for denne Spøg, naar det
betragtes som noget Personligt, og nu snart Contrasterne, snart
Sammentrefningerne imellem det aldeles Personlige, Særegne eller
Individuelle, og det Almeengyldige og over al concret Skikkelse
Ophøiede, blive en uslukkelig Forlystelses Gjenstand. Overhovedet ligger
dybt i Menneskets Væsen en, under mangelunde Skikkelser sig yttrende,
Lyst til en Tilintetgjørelse af Alt, hvad der er noget Bestemt, for at,
hensynsfrit og ubundet af al Form, det frie Liv kan leve alene og
strømme ustandset hen efter sin indre Lyst. Alt Spil er overhovedet en
indre Løsning eller Relaxation i Aanden, for at denne, slippende Alt,
hvad der ellers anspænder og drager Virksomheden hen til een Side, kan,
under en Vexel, hvori intet Punct holder det andet fast, og Alt synes at
ville gaae ud fra hinanden i lutter Momenter, føle en løsnet
Tilværelses frie Rørighed. Under Spil og Leeg oplives derfor Aanden og
forfriskes, idet den, ubunden af enhver enkelt og bestemt Retning, rører
sig frit og løst til alle Sider og i alle sig frembydende Retninger,
uden at give sig hen til nogen enkelt Retnings anspændende tilmed
alvorliggjørende Magt. Selv i Elskoven yttrer sig da nu ogsaa denne Lyst
til at slippe det Henrykkende, Rørende og Dybe, og overgive sig til en
indre Relaxation, under hvilken dog de Elskende midt i den øieblikkelige
Følelse af Ubundenhed, som om de nu havde sat sig i Frihed, komme til
med dobbelt Lyst at føle sig som dem, der ere bundne og fortabte og ikke
ere deres egne, men ogsaa til at føle den sødeste Lyksalighed i at være
saaledes fængslede. Dog, ad mange Veie søger Elskoven hen til den muntre
Spøg og Overgivenhed. Thi paa mange Maader ere Elskende Incitament for
hinanden, til Latter og Lystighed, ligesom til Henrykkelse, ja til
Taarer.

Dog vogte Elskov sig for at sætte al sin Lyst i disse Æggelsers idelige
Vexel, som om den kun kunde besidde og nyde sin Fylde i forbisvindende
Øieblikke og under Opæggelsernes Spil. Elskov vogte sig for at brænde
hen i lutter Higen, saavelsom for at søndersprede sig i en idelig
ustadig Forlystelse; ogsaa vogte den sig for, hvad der endnu vilde være
værre, at sønderslides under idelige Opirringer i en Vexel af
Forbittrelser og Udsoninger. Dog, disse Opirringers Forbærvelighed vil
længere hen Betragtningen over Elskovs Lidenskabelighed bringe os til at
overveie nærmere. Livets dybeste Nydelse er Rolighed og stille Fred, den
udeelte rene, fulde Følelse af noget Evigt og Ufortabeligt, der under al
Vexel blivende er tilstede. Det Æggende, saa meget det end opliver og
vækker Lysten, det Henrykkende og Rørende, saa meget det end er den
dybeste Følelse af Sympathien, som deri bevæger og ryster Sindet, begge
lade os her dog føle os selv som endelige og trængende, for at vi kunne
føle os som salige ved vort herlige Eie. Men der gives en høiere
Sjæleglæde, som Jean Paul, ved at omtale den græske Poesies plastiske
Rolighed, saa smukt siger, „end Lysteus Pine, end Henrykkelsens varme
taarefulde Uveir.“ Hvad der ægger eller rører stærkt, virker kun
forbigaaende, og paa enhver stærkere Opvækkelse eller Anspændelse følger
en Neddalen og Afspændelse, ja, naar intet Andet laae i Øieblikkets
Nydelse, end hvad Øieblikket bragte, følger letteligen en Tomhed. Men
for det ideligen Vexlende, ligger i ægte Elskov noget Dybere, selve
Elskovens indre Leven, til Grund. Denne er det ikke Vexlende, det
Varende, Blivende, hvilket giver Glæden indre Uomskiftelighed. Elskov
yder mere end Fryd i tusinde skjønne Øieblikke. Den er en eneste
uafbrudt indre Nyden af en indre Herlighed, der er os nær overalt og i
alle Øieblikke er tilstede. Saa være den da nu ogsaa en saadan; der være
Ro og Hvile, Tryghed og Tilforladelighed i den. Og Vishedens og
Tilforladelighedens trygge Følelse vorde ei idelig irriteret og anfægtet
ved enhver Leilighed, hvor Dette eller Hiint kunde ville anfægte. Hvor
skjøn er ikke i Schillers Wallenstein Theklas indre Tryghed og Fred i
hendes Kjærlighed, naar hun midt under Uro i tre Dage ei seer sin
Elskede? Hvo der elsker, lære dog at elske med dyb Inderlighed uden den
idelige Higen og Jagen, Brænden og Tørsten, der ei engang under os den
længe og dybt fortvirkende Efternydelse, som først gjør Timens Nydelse
til en ægte. Der gives Timer i Kjærlighed, af hvilke man kan leve hele
Dage, men dem nyder man ogsaa kun inderligen, naar man lever hele Dage
af dem. Den bestandige Unøisomhed og Utaalmodighed, der aldrig skatter
den forundte Lyksalighed tilfulde, nyder den ei heller af Hjertens
Grund. Under alle Elskovens vexlende Momenter, under Overgivenhedens
muntre Spil, under fælleds kjære Sysler og sød Fortrolighed, under
Ømheds villige Opoffrelser, under Henrykkelse, Længsel, Beemod og
Smerte, ja selv under Tvivlens bittre Qvaler, under Fryd som under
Taarer, er dog Det, enhver Elsker føler og nyder, altid hans Kjærligheds
Dybde og Styrke og Væsenheden i det Liv, der er gaaet op for Begge i
deres Glæde ved hinanden. Saa være det da nu og saaledes! Saa gaae da
den dybe Følelse af denne Væsenhed, den milde, stille, vederkvægelige
Følelse af et skjønt Godes trygge, rolige Besiddelse, denne gaae
forhøiet, fuldere og inderligere frem af alle Elskovs vexlende
Momenter!

\begin{center}\rule{0.28\linewidth}{0.4pt}\end{center}

Vi have saaledes betragtet Elskovs herlige Magt og Fylde fra de
forskjelligste Sider. Paa alle de Maader, paa hvilke en Gjenstand kan
skjænke os Fryd, viser Elskov sin kraftige og velgjørende Indflydelse.
Meget i Livet yder Velbehag og Glæde ved det Skjønne og Herlige, som
deri umiddelbart fremstilles for den rene Beskuelse alene og for den
blot heraf udspringende Nydelse. Meget atter virker dybt ved at vække
vor Medfølelse, saa at Det, der sætter os i Bevægelse, er en Andens
Glæde, en Andens Værd eller Lykke, ja ogsaa er den blotte Følelse af den
Forening, der knytter os til et andet Væsen, og gjør, at vi give os
ganske hen til dette, og at Intet kan være fremmedt for os, hvad der
rammer eller vedrører dette Væsen. I alle disse Tilfælde er Følelsen
desto stærkere og inderligere, jo mere vi kunne glemme os selv i
Gjenstanden, jo mere alt Personligt kan tie og tabe sig i den fulde
Hengivelse til den bevægende Indflydelse. Paa den anden Side gives der
Meget, der, vækkende Følelsen af en personlig Trang og Fornødenhed
opfylder os med Attraa og Begjær, og Nydelsen bestaaer nu i en
Tilfredsstillelse. Endeligen er Alt, hvad der opliver og forhøier vor
personlige Selvfølelse, af ikke ringe Magt i at fremkalde Glæde. Paa
alle disse Maader virker nu lykkelig Elskovs velgjørende Kraft. Saa
omfattende er den Følelsernes Fylde, hvori Elskov udbreder sig. Her er
qvægende Contemplation; inderlig Sympathie og Hengivenhed, dyb
Fyldestgjørelse, og endelig ogsaa høi Selvfølelse ved Elskovs baade
opløftende og frigjørende Magt, hvorved selv den usigelige Ydmyghed
virker med og giver Følelserne Dybde og Inderlighed.

Men endnu er her en anden Modsætning af Følelser at iagttage. Til Elskov
hører ikke blot en uendelig Tiltrækning, men ogsaa den Fjærnhed mellem
begge Elskende, hvorved den Tohed og Modsætning bevares, uden hvilken
Kjærlighedens enigende og sammensmeltende Kraft ikke kunde være idelig
frisk og ny. Til denne svarer i Følelsen den til al Kjærlighed
væsentligen henhørende Agtelse, der i det elskede Væsen erkjender et
Værd og en Eiendommelighed, der skal holdes i Hævd og skal nyde sin
fulde Ret og Fremme.

Til at leve fuldeligen udfordres dernæst baade Sands for Det, hvad Livet
bringer, og Kraft i at virke tilbage. Fuldest nyder Den Livet, der baade
møder dets oplivende Mangfoldighed med aaben Sands og let Bevægelighed,
saa og tillige med indtrængende Sind fatter dets indre Væsen og Dybde;
der er drivtig, kraftig og udholdende i Daad, men dog under Alt, hvad
der ægger til med Energie at gribe ind i Livet, bevarer Sindets indre
Ligevægtighed og Ro. Disse ere de Sindsstemninger, Enhver stræbe at
drage frem i sig og at bevare. Med et let Sind og med et dybt Sind, med
et kraftigt Mod og med et ligeligt, besindigt Mod maa Livet leves. Men
saaledes lever den lykkeligt Elskende det. I Elskov er Sands for
Tilværelsens Livelighed og for munter Spøg, og Sands for Livets dybe
Alvor; her er Energie og Iver i at ville og at stræbe, og tillige er her
en rolig Tryghed i den fulde Følelse af en herlig Stats tilforladelige
Besiddelse.

Men til denne Kraft og Fylde kommer Elskov kun efterhaanden. Som Alt paa
Jorden maa den fra en Spire af voxe op og udfolde sig lidt efter lidt og
i Tidens Fremgang naae sin fulde Blomstren, hvortil den dog kun kan
komme, idet vi selv drage dens Kraft og Fylde frem, og selv lade dens
indre Varme gjennemtrænge de mangfoldige Ledemod, hvori vort indre Liv
rører sig, saa at dens bevægende og oplivende Kraft kan fra Hjertet
brede sig ud til Extremiteterne.

Al inderlig Elskov er i sit Udspring genial. En dyb Natur rører sig af
sig selv i os, og dybe Ideers Magt ligger til Grund for dens Virken.
Dens indre dunkle Bevægelser fremkalde Elskovs stærke Længsler, inden
den finder sin Gjenstand; og finder den denne, da er Mennesket, førend
han selv endnu ret veed det, midt inde i Kjærligheden, grebet, hentaget
og fængslet af de stærke Følelser, der af sig selv have opfyldt hans
Indre. Han har umiddelbart levet sig ind i sin Elskov, og elsker, fordi
han ikke kan Andet, ifølge en indre uimodstaaelig Nødvendighed, dog en
saadan, der tillige er en høi Frihed og inderligen føles som Frihed,
fordi den føles som den, der gaaer frem af Individets eget inderste
Væsen.

Men hvad der saaledes af sig selv er udsprunget i Mennesket, og, som en
vældig Naturindflydelse af høiere Art, bevæger ham i de stærkeste
Sindsrørelser, dette komme nu til inderlig Klarhed i ham, og han tilegne
sig det af Hjertet og tilfulde! En klar Besindighed træde her, som
overalt, til det Geniale, for at det ved en høiere Beaanding givne Eie
ikke skal besiddes med Blindhed, eller Henrevenhedens skjønne Opløftelse
gaae forbi som en Beruusning. Men det Geniale (τὸ δαιμόνιον, som Platon
kalder det) tilegner sig den Besindige, deels erkjendende, deels
villende.

I en Erkjenden, Følen og Villen har overhovedet det aandelige Liv og
Alt, hvad der lever deri, sin Bestaaen. Derfor maa nu ogsaa til Elskoven
og den Følelsernes Fylde, hvori den udtaler og bevæger sig, komme baade
en Erkjendelse og en Villie. Ved begge bliver den hele Elskov med Alt,
hvad den indeholder, ligesom bekræftet og sat (poneret) paa Ny.
Mennesket maa med Opmærksomhed agte paa sin Elskovs Indhold og
mangesidige Kraft. Han maa sige sig selv, hvad det er, der nu er skeet
med ham, han maa søge med Klarhed at fatte og erkjende, hvad ret er, der
er blevet ham tildeelt, og saaledes virke til, at det kan leve og
udbrede sig i ham med sin fulde Styrke. Da kommer han først til at nyde
det tilfulde og ret at bevare sig det, idet det nu knytter sig til hans
indre Livs hele øvrige Indhold, og, medens det ved den klarere
Bevidsthed selv opklares, nu tillige viser sin velgjørende Magt i alle
Retninger.

Men ligesom paa den ene Side ved en Erkjendelse, saaledes maae vi paa
den anden Side ved en fast Beslutning tilegne os det nye Liv, der gaaer
op for os i vor Elskov. Vi elske, fordi vi maae elske; men hvad vi saa
uimodstaaeligen maae og føle os drevne til, dette maae vi nu ogsaa ville.
Den Elskende vil sin Kjærlighed, saa ville han den da nu ogsaa af
Hjertens Grund, med et sig selv besiddende Sinds indre Frihed. Vi leve
os ind i Elskoven. Men den inderlige udelukkende Forbindelse for hele
Livet skal udtrykkeligen indgaaes. Hvad der er skeet af sig selv, skal
Mennesket med fri Selvbestemmelse bekræfte, derfor ogsaa først prøve sig
selv, om han ogsaa af Hjertens Grund kan ville og virkeligen vil, hvad
han nu selv skal beslutte. Forholdet blive da til et ægte moralskt
Forhold, stiftet ved en fri Villies og Beslutnings faste Bestemmelse! I
denne faste hellige Beslutning, nu ogsaa af Hjertet at ville tilegne sig
den vundne Lykke og holde den fast og ganske give sig hen til den, og
derved at fixere sig selv og den let bevægede Lyst og Drivt, har
Troskaben sin Rod og sit Sæde.

Troskaben maa udgjøre Kjærnen i den herlige Samvær, hvori tvende for et
heelt Liv forbundne Væsener bestandigen voxe mere og mere sammen. I ægte
Elskov tager den, under Samlivets hulde Bane og under Følelsen af Begges
Uundværlighed for hinanden, bestandigen til, og, fæstende sig fast i
begge Hjerter, fæster den begge Hjerter til hinanden. Ved den faaer
Elskov først Væsenhedens Præg. Intet er sødere, dybere og helligere i
Elskoven, end Følelsen af denne Troskab fra vor egen og Visheden om den
fra den Elskedes Side. Hvorimod Tanken paa Utroskab er saa oprørende og
saa frygtelig, fordi Utroskab gjør, at selv det Høiere, Herligere i
Jordelivet viser sig som Noget, der ingen Bestaaen har, ingen Ægthed,
men, undergivet jordisk Omvæltning, er Intet, ligesom alt andet
Jordiskt. Alt, hvad Elskoven kan yde og bringe Frydefuldt og Skjønt, er
det kun under Forudsætning af denne Vished, Tryghed og Paalidelighed.
Derfor føler Mennesket i sin Troskab det Fasteste og Stærkeste, og i
Visheden derom fra den Elskedes Side det meest Magtpaaliggende, i den
hele Elskov. Hvad kunde være frygteligere, end midt under Elskovens
fuldeste Følelse allerede at maatte frygte for den Orm, der, skjult
under Livets opblomstrende Flor, skulde ogsaa undergrave det Aandeligste
og Reneste. I vor inderlige Troskab trykke vi i Aanden vor Elskede til
vort Hjerte, ja bære den Elskede ligesom indesluttet i vort eget Bryst,
som den helligste Deel af os selv. Derfor fornemme vi Elskovs Følelse
som en evig Følelse og kalde den saa. Hvor ødende, ruinerende er derimod
Utroskab. Man tænke sig et ægte, inderligt Elskovsforhold, man tænke sig
den Livets skjønne Fylde, hvori Elskov udbreder sig, hvorledes Alt, hvad
der omgiver Mennesket faaer et nyt livfuldere Udseende og tusinde
Smaating blive Midler for Yttringerne af den Kjærlighed og Hjertelighed,
hvormed ideligen Hjerterne aabne sig mere og mere for hinanden og knytte
sig til hinanden. En rig Stat af de skjønneste Erindringer tager
Mennesket med sig ind i Fremtiden; Livet nydes ei i Momenter alene, den
dybe varige Følelse, der gaaer gjennem alle Momenter, giver hele Livet
Continuitet. Og nu tænke man sig Elskoven forsvunden, Forholdet opløst,
være det ogsaa kun indvortes og sjæleligen. Det er nu, som om en
fiendtlig Dæmon var traadt imellem Mennesket og hans skjønneste
Livsperiode; til alle Erindringer om denne hefter sig den Tanke, at det
alt er forsvundet, og forstyrrer og forbittrer dem. Mennesket er som
afskaaret fra sin Nutids Rod, den skjønne Fortid; den herlige Eiendom,
han besad i en rig Erindring, er borte, hans Liv er som overskaaret, og
dets sødeste Deel kastet bort. Den hele, lange, skjønne Ungdomstid er
endnu mindre end en forsvunden Drøm; thi paa denne kan man endnu tænke
med Lyst, men ikke mere paa hiin. Saa ruinerende er Utroskab, at den
ikke blot opløser det Nærværende til et Intet, men endog sønderriver det
bevarede Billede af den hele Fortid. Vistnok gives der Tilfælde, hvori
et Menneske, for at bevare Troskaben imod sig selv og mod Gud, maa, hvor
meget end Hjertet monne bløde derved, sønderrive et Baand, der knyttede
ham til et andet Menneske. Vi ere den Elskede troe, fordi vi ere det
Sande, Gode og Skjønne troe. Troskaben i Elskov har sin Rod i en høiere
Troskab, og maa falde hen, naar den ei kan næres ved denne, men
modstriver den. Dog, saadanne Tilfælde komme vi til at betragte
nærmere, naar vi komme til at tale om ulykkelig Elskov.

Men endogsaa udenfor disse Tilfælde see vi endnu alt for sjeldent Elskov
at fremkalde den Forbindelse, den kunde og skulde føre med sig, en sand
hjerteqvægende inderlig Sjæleforening for et heelt Liv. Hvor skjønt og
husvalende er det at see Mennesker, der levede med hinanden i mange Aar,
ja til Alderen gjorde dem affældige, endnu at elske hinanden med
inderlig Hengivenhed og dyb Fortrolighed og Tillid, ja endnu med
ungdommelig Lyst og Velbehag at glæde sig ved at skue hinanden. I
Ungdommen haaber Enhver, der elsker med Fyrighed og Enthusiasmus, ja er
sig forvisset om, at hans egen Elskov vil have denne Udholdenhed, han er
vis paa, at den Aand og det Hjerte, han elsker i sin Elskede, ophøiet
over Affældighed, vil endnu udaf den forældede Form tale til ham med
uformindsket Kraft og endnu glæde og henrykke ham. Men imedens Verden
bestandigen er fuld af unge Forelskede, der ere ganske indtagne i
hinanden, er den tillige fuld af Egtefolk, hos hvilke Egteskabet synes
at være blevet, hvad et gjængse Ord endog vil sige, at det pleier at
være, Kjærlighedens Grav.

Det Første i Troskaben er Fastheden og Paalideligheden, der fører med
sig den usigelige Vished, hvormed den Ene forlader sig paa den Anden, og
stoler paa et af alle jordiske Tilstød og Sammenstød uanfægteligt
Sindelag hos den Anden, ifølge hvilket Begge bestandigen blive ved at
være hinanden, hvad de have været hinanden. De føle sig knyttede sammen
deels ved Kraften af det Gode og Herlige, de i hinanden skue og
erkjende, hvorfor ogsaa med slette og forkastelige Væsener en saadan
Pagt ei kan have Bestand eller Tryghed; men deels knyttes de ogsaa
sammen og blive ved at være hinanden Det, de ere hinanden, fordi de have
været det; fordi Livet eengang har samlet dem og ladet dem med og ved
hinanden skue ind i dets dybere Væsen og nyde dets dybere Fylde.
Hvorved det viser sig, at Den, der har, ham gives der fuldtop, idet
ethvert Øieblik, hvorved Erindringen med Glæde dvæler, gjør Foreningen
mere inderlig og fast.

Men endnu er i Elskendes Troskab et Andet at bemærke, hvad der ikke, som
hiint Første, gjelder om enhver Sjæleforbindelse, eller ethvert ægte og
dybt Venskab. Dette Andet er det Udelukkende i Foreningen: at Begge ikke
blot med uanfægtelig Fasthed holde ved hinanden og have et sikkert og
tilforladeligt Hjem, en usvigelig Tilflugt, hos hinanden; men at ogsaa
Ingen af dem attraaer at være, eller bevæges til at være, det Samme for
nogen Tredie; at altsaa Elskov kun forener Tvende paa denne Maade.
Ogsaa her viser sig en fast Beslutnings og Villies Kraft, og ogsaa her
er den fornøden. Det er den til Grund for Elskov liggende Naturdrivt, som
herved kommer i Betragtning. Denne ofte let bevægelige Drivt, der kan
drages til alle Sider maa Mennesket selv fæste eller fixere ved en
hellig Villies Kraft, at den ei, opægget til uordentlige Bevægelser,
skal drive ham ustadig og veiløs om, og sætte Sindet i Uro og Tvedragt
med sig selv. Her er Noget, hvorpaa der maa lægges et Baand og en Tømme,
og hvor en værnende Pligtfølelse, en Følelse af Respect for det Hellige,
maa træde til, næres og bevares. Kun alt for ofte førte Mangel paa
Aandskraft i denne Henseende til Det, der, istedetfor at stille og qvæge
Hjertet og bringe Fylde, lod Hjertet tomt og lagde Sindet øde.

Imidlertid: mon Elskov skal føre med sig, at Livets udbredte Herlighed
skal tabe sin Tillokkelse for det elskende Gemyt? Den Kjærlighed, som
med saa mægtig Kraft vækker Aanden og aabner Øiet for det Skjønne og
Guddommelige i Livet, mon den nu skal føre saavidt, at det samme Øie,
heftet paa en eneste Gjenstand, skal blindes for alle andre, eller lukke
sig for dem? Ingenlunde. Saa hellig end den Forbindelse, hvortil Elskov
fører, i hiin ene Henseende maa holdes som en udelukkende, saa herske
iøvrigt, udenfor Det, der henhører til den her i Bevægelse værende
Natur, det misundelsesfrie Sind, den ἀφθονία, hvilken Grækeren roste saa
hos sine Guder. Der maa elskes med Frihedssands. Saavidt skal ikke det
Udelukkende gaae, at nu al Contemplations Glæde og al aandelig
Meddelelse, al levende Interesse, skal concentrere sig om den Elskede
alene, og den rige vidt udbredte Verden høre op at være enhver af de
Elskende, det den forhen var dem, en Verden, i hvilken Hjertelighed,
Omhu, Interesse og Kjærlighed fandt en viid Sphære for sig, og Hjertet
overalt glædedes ved at møde sin Gud. Misunder dog ikke Eders Elskede
det muntre, livelige, aabne Sind, den friske, altid rede Glæde ved
Livet, I, som elske! Nærer meget mere ogsaa i denne Henseende den
inderlige, søde Omhyggelighed, den hjertelige Baretægt, som vil, at den
Elskede skal have Alt paa det Bedste! Skuer det med Glæde, om en høiere
Glæde ved hele Livet og ved alt Skjønt og Stort, som bevæger sig i
Livet, udspringer af og næres ved Eders Elskedes Kjærlighed! Tillader,
at den Elskedes Lykke bliver en sand Lykke for denne! Elskov skal i
Sandhed ei bornere; og hvad der aabnede det frie Blik paa Livet, skal ei
siden omtaage det. Ei engang den første, fornemste og høieste Gjenstand
for sin Elskedes Høiagtelse, kan den Elskede altid være, saa lidet som
den eneste Gjenstand for den Elskendes hele Fortrolighed og sammes
Tilflugt i alle de Anliggender, der trykke eller bevæge Sindet.
Kjærlighed, som i saa høi Grad kan gjøre seende, skal ingensinde hilde
det rene klare Omdømme. Ingen bliver den Elskede utro, fordi han bliver
sig selv tro, og ei vil nedskue eller forvende, hvad der som Sandhed
taler i hans Indre, eller fordi han stræber at kæmpe den Kamp, vi Alle
have at kæmpe, med et redeligt Sind, der griber til enhver Hjelp og
Støtte, hvortil det trænger. Kun med Hensyn til den ikke altid
tilbørligen ordnede Naturs Rørelser vogte sig Enhver og have Øiet
aarvaagent over sig selv! Thi her er Fare, og hvor noget Uskyldigt og
Godt ligger for nær ved Noget, der ikke er godt, slippe man heller det
Uskyldige og Gode, end at udsætte sig for Fare og Ufred, og glemme ei at
advare sig selv eller hinanden! Uden indre Fred og Ro kommer jo det Gode
os ikke tilgode.

At elske med ægte Frihedssands og lade Egoismen ret gaae til Grunde i
Kjærligheden er imidlertid ikke let at lære. Men een Betragtning kan
maaskee endnu oplyse det Fremsatte og aabne det Indgang til Hjertet: et
Henblik paa det andet Liv.

Enhver, der har prøvet Livet og Elskoven, endog kun nogenlunde, Enhver,
der attraaende og higende har følt den Jeghedens Braad og Spore, der er
i al Attraa, især naar den er heftig, fordi vi under al heftig Attraa og
Higen komme til at føle os selv i skarp Modsætning til en muligen
modstridende Omverden, til en maaskee ei strax for os tilgjængelig
Gjenstand: Enhver, der har prøvet Livet og Elskoven og begges Skjæbner,
og har erfaret, hvad begge give, og hvorledes begge leves og nydes, om
der overalt i dem er klar Ro og reen stille Nøie, eller om de kunne
forurolige begge og hvorledes og hvorfor: --- enhver Saadan spørge sig,
hvorledes han vil tænke sig et Liv i Evigheden, naar han skal tænke det
som et over Jordens Trængsler opløftet, reent og klart og saligt Liv,
men dog, som et Liv, hvori Alt, hvad der elskedes, findes igjen, og
ellers saa ganske menneskeligt, at vi kunne tale paa jordisk Viis derom.
Skal ogsaa det være et Liv, hvori udelukkende Besiddelse skal herske?
Skal ogsaa det salige Liv være et brændende, og ei engang Hjemgangen til
det sande Hjem bringe Ro og Fred? Eller maa ikke snarere Salighedens
Hjem ogsaa være Fredens, og maae vi ikke snarere haabe, at i dette Hjem,
naar Livet har hævet sig til ægte Frihed og fuld Personlighed og
individuel Fylde, nu den Egoismens Braad kan falde ud, hvormed vi her
plagedes, medens vi derved ansporedes, at i Herlighedens Rige Ingen
frier og Ingen føres hjem, fordi der ei er Tale om at besidde? Saa at
vistnok, hvad der her udgjorde et skjønt Samlivs indre Væsen og Fylde,
ingenlunde skal troes at være fortabt eller forbi for et følgende Liv,
nei at tvertimod ogsaa i denne Henseende Livets aandelige her tilvundne
Eie tages med hisset; saa at Elskende forstaae hinanden inderligen og
hænge ved hinanden og elske hinanden, netop fordi de elskedes her, ja at
noget vist Specielt endog kan vedblive i deres Kjærlighed; men at dog
heri eet Væsen ikke fordrer Noget fremfor et andet, eller at overhovedet
et Fremfor ikke høres, eller endog kun kan komme op i Tanken. Thi hvor
et Fremfor tænkes, der rører sig allerede det arme Jeg, det fattige
Selvhensyn, under hvis hildende, standsende Indflydelse vi sandeligen
her paa Jorden have lidt nok. Man spørge sig ogsaa kun, hvorfra hist en
Lyst til noget Udelukkende skulde komme. I den rene Glæde ved
Beskuelsen, den inderlige Nydelse af den usigelige Herlighed, der
oplader sig for vort Øie i et skjønt Væsen, er der dog Intet, hvorfor vi
skulde misunde Nogen, hvilkensomhelst, at lyksaliggjøres derved
hvorsomhelst. Og hvad Sympathien og Meddeelagtigheden angaaer, og den
Hengivenhed, hvormed den Ene giver sig hen til den Anden, da er her for
Den, der nyder den Lykke udeelt at besidde et andet Væsens hele Hjerte,
ingen Grund til misundelig at vredes over, at det samme Hjerte er lige
saa velgjørende for et andet Væsen, og giver sig lige saa reent hen til
dette. Thi det er et reent Hjertes Magt og dets Lykke, at det i sin
fulde Heelhed kan give sig hen til tusinde, udeelt til Enhver, naar det
nemlig giver sig saaledes hen, at det i og med det samme (eo ipso) giver
sig hen til Herren, hvem vi møde paa tusinde Veie. --- Dog, til denne
høie Misundelsesfrihed, denne ἀφθονία, ere kun meget faa rene,
fordringsløse, lykkelige Naturer i Stand til at bringe det i et Liv,
hvori Egoismens Braad saa vanskeligt kan komme til at hvile. Imidlertid
stræbe Enhver dog i det Hele at leve og røres i den Aand og med det
Sind, hvis fulde udeelte Virken han haaber at nyde og besidde i en anden
Verden!

\begin{center}\rule{0.28\linewidth}{0.4pt}\end{center}

%% ============================================================
%%  ANDEN BOG  (printed pp. 70--129; PDF 86--145)
%%  Fraktur in the original. Transcribed and image-verified.
%% ============================================================
\chapter*{Anden Bog.}
\markboth{Anden Bog}{Anden Bog}

Elskov skal og kan efter sin indre Natur føre til en dyb
Hjerteforbindelse, hvorved Tvende, forenede for et heelt Liv,
bestandigen voxe mere inderligen fast til hinanden i et skjønt Samliv
til Velbehag for Gud og Mennesker. Hvorfor fører Elskov da saa sjeldent
til den skjønne Forening, hvortil den kan føre? Elskov er lig et
Frugttræ i Vaar, der bedækket med Blomster udsender en Duft, som
forkynder den kommende Sommers Frugtbarhed og Fylde. Egteskab skulde
være ligt det Frugttræ, der hele Aaret igjennem bærer paa eengang
Blomster og Frugter, begyndende og fuldmodne; thi saadan er Aandens
Kraft i Mennesket, at hans Blomstretid ei er forbi, fordi hans Frugttid
er kommen. Hvorfor er da saa ofte et Egteskab, selv hos ret brave og
gode Mennesker, ligt et Træ, hvis Frugter ere, nogle visnede, og nogle
ormstukne, og næsten ikke een ret fuldkommen moden og tillokkende?

Først og fremmest fordi Menneskene saa ofte gaae besindighedsløse ind i
det vigtige Forhold og saa sjeldent minde sig om, naar det er indgaaet,
at de selv maae agte paa at drage Omsorg for at bevare sig, hvad der er
skjænket dem. Allerede førend den Forening, der skal være gyldig for
hele Livet, indgaaes, vaage man over sit Hjerte og troe ei Alt at være
trygt og paalideligt, fordi de opvakte Følelser med Heftighed og Styrke
concentrere sig om en udkaaret Gjenstand. Skjønheds og Yndes,
overhovedet visse ydre Tillokkelsers Magt, virkende paa Følelserne
gjennem den opvakte Sands, ere som oftest Det, der drager Elskov frem og
fængsler Blikket og Tilbøieligheden. Men allerede Det glemmer man alt
for ofte i Tide at sige sig, at en Elskovsforstaaelse skal være en dyb
Forstaaelse, en ægte og dyb Hjerternes Meddelelse, og at Elskende derfor
bør stræbe at gaae ind i hinandens Naturer. Man kan ei gjentage Dette
for tidt, naar man seer, hvor hyppigt den hele Elsken bliver staaende
ved den Lyst og det Velbehag, hvormed Tvende see hinanden og omfavne
hinanden, hvor hyppigt Mennesker forelskes i hinanden og knytte
Foreningen uden at de sympathisere, uden at de forstaae hinanden af
Hjertet, uden at en dyb og hjerteqvægende sjælelig Meddelelse er imellem
dem, og uden at de tænke paa at fremlokke og fremelske det
Hjertesamqvem, der dog kunde fremkomme. Den af en Beskuelsens Lyst
udspringende, være det og af de stærkeste Drivter bevægede Elskov, der
ei fører til denne dybe Fortrolighed, til den inderlige Hjertelighed og
dybe Tilforladelighed, er ingen ægte Elskov. Dog er det ei sjeldent, at
en vis Skjønheds og Yndes, eller en vis Aandeligheds og Opvaktheds Magt
saaledes drager et Væsen hen til et andet, og saaledes ægger alle
Lyster, alle Drivter, al Tilbøielighed op, at den Ene føler sig betaget
af den heftigste Elskov til den Anden, uden at der er sand Samstemming
og Sympathie, og uden at det kommer til en sand Hjertens Forstaaelse og
Meddelelse. Men fattes den, da gjelder det, at være redelig mod sig selv
og sige sig, at den fattes. Man maa altid have Mod nok til at være
sanddru mod sig selv. Man glemme aldrig, at Elskov ikke blot skal drage
et Øie til et Øie, men og knytte et Hjerte til et Hjerte. Hverken
Beskuelsens inderlige Lyst og Glæde eller den i Bevægelse værende Drivt
henrive Nogen, der attraaer Elskovs fulde og dybe Lyksalighed, førend
ogsaa en ægte Sympathie viser sig, og en dyb og fuldelig Meddelelse er
traadt ind imellem Begge! Saa tung og trykkende kan ikke endog en skarp
og skjærende Følelse af et Savn og en utilfredsstillet Længsel hvile paa
Sjælen, som den Følelse vilde det, at være for bestandigt knyttet til et
Menneske, „som ei forstaaer vort inderste Væsen og som dog uophørligen
fordrer vort Hensyn. Hiin Længsel fylder dog Sjælen, men sligt et
Forhold frembyder os uden Ophør den evindelige Tomhed.“ Men hvor
ubesindigen fare eller drages mange Mennesker ind i slige Forhold. Og
den egne Blindhed understøttes kun alt for ofte ved Andres. Især gives
der mangt et Tilfælde, hvori man maatte ønske en Datter Sjælsstyrke nok
til, ei at gjøre en omhyggelig Moder ulykkelig ved at lade sig selv
gjøre ulykkelig af hende, naar hun af en lidet christelig Ængstelighed
for Fremtiden vil ved Forestillinger bevæge det datterlige Hjerte, der
ei bevæges af nogen Tilbøielighed. Den barnligste Eftergivenhed og
Billighed kjølner ikke Blodet, eller tilfredsstiller ikke Naturen, og
sikkrer ikke for en haard Kamp og for sønderslidte Følelser; og en
ulykkelig Datter er en for stor Straf for en Moders Utaalmodighed til,
at ikke en øm og kjærlig Datter skulde ved standhaftig Modstand
forskaane sin Moder derfor. Ikke at tale om de Tilfælde, da utaalmodige
Forældre, istedetfor en simpelhen ulykkelig Datter, faae en fordærvet
Datter at bære paa Samvittigheden, og tilmed en dobbelt ulyksalig. Et
Egteskab, der indgaaes, efter mange Beraadslagninger, med et halvt
Hjerte, er i Sandhed ikke sluttet i Himmelen, og det skulde dog, efter
et gammelt Ord, ethvert Egteskab være. Hos Manden er iøvrigt
Ubesindigheden og Daarligheden ligesaa stor, naar han egter uden at
elskes, være han endog dreven af den heftigste Elskov. Og er det da at
elske, at gjøre Den ulykkelig, man holder af? Ingen stole paa, at
Gjenkjærligheden nok vil komme. Kan den komme, saa kommer den lettest og
snarest, medens den Elskedes Hjerte endnu er frit. Ligger først en
Pligtfølelses haarde Tryk paa det underkuede Hjerte, kommer den
langsommere. Og hvilken Daarlighed, at springe det Første og Vigtigste
over, for at ile til Det, der dog kun er Noget under Forudsætning af
dette Første, under Forudsætning af Gjenkjærligheden. Og gives der vel
Meget, der er saa umandigt, som det, med en Art af Overvældelse og Tvang
--- ogsaa den sjælelige Tvang er en Tvang --- at bemægtige sig en
Qvinde, der endda modstræber med alle sine Følelser?

Tvende Perioder i Menneskets Liv ere her iøvrigt vel at adskille. I
Livets første ungdommelige Opblomstren kan Intet tilfredsstille uden det
Høieste og Herligste. Den dybeste Fyldestgjørelse for Sind og Hjerte
søges, og en Forbindelse uden fyrig Elskov og inderlig Hentagethed er
her Intet. At give sig hen for udvortes jordiske Hensyn er i alle
Tilfælde en næsten endnu større Daarlighed, end den, at sælge sine
Tænder for lækkre Spiser. Hvad Nydelse skal Nogen have af Spiser, naar
han ingen Tænder har, eller af Rigdom og jordisk Høihed, naar han er
ulykkelig gift? Men i en ungdommelig Periode kan Individet ei engang
give sig hen blot ifølge Erkjendelsen af det andet Væsens Værd alene,
naar en inderlig Tilbøielighed ikke tillige drager Hjertet og bevæger
det med uimodstaaelig Elskov. At give sig hen uden egentlig Elskov, være
det end, at Hjertet ellers bevæges af ægte Agtelses og Velvillies
Følelser, er i den ungdommelige Alder altid betænkeligt. Thi det kan da
letteligen skee, at enten Elskovs fyrige og livlige Opvaagnen og tilmed
dens hele Sjælen opvarmende og oplysende og lyksaliggjørende Kraft
derved undertrykkes og qvæles, saa at Mennesket derved udelukker sig
selv fra den Lykke, der maaskee dog kunde engang have aabnet sig for
ham; eller at endog Faren er større, at en dybere Trang og Lyst engang
virkelig kommer til at vaagne, men at den nu vaagner som en
foruroligende, der bringer Hjertet i Tvedragt med sig selv og med den
dybt sig rørende Natur, som ei kan finde Fyldest, Ro og Qvægelse ved
Det, den rolig overveiende men ei tillige uimodstaaeligen hendragne
Beslutning har grebet. Derimod senere, efterat Naturen har rørt sig i
sin ungdommelige Styrke, efterat Livet er prøvet og Følelserne have
udbredet sig i Individets Indre med den Kraft og i den Art, det os
beskjærede Liv medførte, senere vil, naar Individet i hiin første
Periode ei fandt og kom i Besiddelse af Elskovs fulde Lykke, en anden
Tilstand indtræde. Have Hjertet elsket, men ei funden Forening, have det
overhovedet elsket med mere eller mindre Styrke og Inderlighed og med
denne eller hiin Lykke, eller have det ogsaa ikke elsket egentligen, men
dog under Længslernes indre Magt følt Elskovs dybe Sødme, eller have det
levet roligere hen og ikke været meget bevæget af Elskovs Lyst eller
Længsel: i alle Tilfælde har Livet nu formet og dannet sig og antaget en
bestemtere Skikkelse; Følelsernes Fylde har virket, som den kunde, og
Naturen har gjennemgaaet sit første Stadium. Individet veed nu, hvad
Elskov er. Vistnok, jo bestemtere det veed det, jo klarere det har lært
at overskue sin Tilstand, jo mere det bringer en Fordring med,
istedetfor at den Unge med Glæde tager hvad der gives, som det gives,
desto mindre kan det nu finde en Gjenstand. Hvad der gjør Sagen saa let
i Ungdommen, at Livet selv først skal forme sig aandeligen, og derfor
letteligen former sig med og ved, og tilmed i Harmonie og
Overeensstemmelse med, en fremmed Individualitet, fordi Alt endnu er
blødt og i Vorden, Dette er der nu ikke mere. Jo fastere Livet er
fixeret, desto mindre er det skikket til at løsnes. Imidlertid, lad nu
ogsaa Drivten være bleven holdt nede, først ved den ungdommelige
Enthusiasmus, senere ved den roligere Erkjendelses og Villies Magt, lad
Sindet en Tid lang have dulmet: Naturen dulmer ei, den kræver og
trænger. Dens Bevægelser røre sig op igjen, og Mennesket vil nu let
komme til at sande Pauli Ord: det er bedre at giftes end at lide Brynde.
Nu er han tilfreds, om han, istedetfor ungdommelig Elskovs Himmel, ei
faaer Andet end den jordiske Lyksalighed, hvortil Elskov kan føre. Og er
hans Stræben reen, er han ei blot rolig og besindig, men og tro mod sig
selv og sandru, saa at han vogter sig for at hilde og besnære sig selv,
eller at forhaste og bedøve sig, og lader han enhver Bevægelse, der nu
vil henrive den maaskee uroligt higende Natur, virke langsomt, til den
kommer til Klarhed, og, naar Forholdet først overskues, kan finde den
roligere Betragtnings Samstemming, --- saa kan det vel endnu skee, at
den indre Naturdrivts Styrke saaledes kan løsne Hjertet til at erkjende
og føle alt det Velgjørende, der skues i et andet Væsen, at han endnu i
sin Elskov kan finde den Himmel, uden hvilken dog selv den hele jordiske
Lyksalighed endnu er for Lidet. Saa at endnu en sand Tilbøielighed og
Hjertets Tilknytning til et andet Væsen kan drage den dybe Fortrolighed
og Sympathie, saavelsom den sande Contemplations Lyst og Glæde og en
ægte aandelig Fyldestgjørelse frem. Ingen see skjævt til et
Elskovsforhold, der er knyttet paa denne roligere Maade i denne Livets
Periode, hvori Ungdommens friske, herlige Forhaabninger maae falde hen.
Ingen see skjævt til Nogen, fordi Herren var ham mindre naadig end
Andre, og fordi han lærte at sige sig Dette selv. Ogsaa slige
Egteskaber kunne være --- for igjen at bruge det gamle Ord --- sluttede
i Himmelen, eller sluttede ifølge en saa reen og sikker Følelses Stemme
og tilskyndende Magt, at Individet er sig bevidst, at det vil, fordi det
maa ville, og at Udfaldet beviser Følelsens Ægthed. Kun er det her
dobbelt at erindre, hvad den ungdommeligen bevægede Sjæl endda snarest
føler, at intet Elskovsforhold er Noget, naar ei i samme et Hjerte
knytter sig til et Hjerte, og Kjærlighedsforstaaelsen virkelig er en dyb
Forstaaelse og Meddelelse. Alt Egteskab maa indgaaes saaledes, at en
indre høiere Stemmes tilforladelige Samstemming lader Individerne blive
visse paa, at det ei er Vilkaarlighedens Spil og eget Tykke, der driver
dem. Ethvert Egteskab maa sluttes med et fromt Sind (pio animo). Men det
er herved aldrig at forglemme, at, hvad der kan være paa rette Sted i
Livets anden Periode, ikke derfor strax er det i den første
ungdommelige.

Imidlertid er det Egteskab ikke altid sluttet i Himmelen, hvilket den
heftige Lyst river Tvende hen til at indgaae, inden de komme til at
sunde sig, det er at besinde sig og komme til nogen Klarhed over sig
selv. Med en Forening, der skal gjelde for hele Livet, behøver man jo
dog, især naar Hjertet først føler sig vis i sin Sag, ikke at haste. Der
er dog alligevel saa megen Grund til at gaae roligt og langsomt frem.
Unde man dog et ungt Par og unde det sig selv en livelig
Forelskelsesperiode, og nyde det dog i langsomme, rolige Drag den milde,
qvægende Duft, der, bebudende en underfuld Fremtid, hvis Fylde synes
uudsigelig, gjennemaander Elskovens skjønne Maidage. Det er især for
Qvinden saa velgjørende, at nyde en lang Elskovsvaar; thi hun skal gaae
et Kald imøde og maa have Ro og Tid til at samle Sind og Eftertanke til
den ganske nye Virkekreds, og berede Hjertet til den Tid, som venter. Og
er det ikke ogsaa Synd for et qvindeligt Væsen i Ungdommens første
Blomstren, nu allerede strax, efterat et nyt Liv er gaaet op for hende i
hendes Elskov, at skulle trækkes ind i Egteskabets Sorger, ja bære
Moderens Byrder og udholde de urolige Nætter? Ikke at tale om, at
Forhastelsen let kan føre til det Unaturlige, naar den nye Spire
fremnødes, inden den første Frugt er fuldmoden, saa at Moderen endog i
organisk Henseende endnu selv voxer ud, medens allerede et nyt Væsen
voxer under hendes Hjerte. Hvorved det er at bemærke, at Modenheden,
endog i organisk Henseende, endnu ei strax er der, fordi den fulde Væxt
er der. Hertil kommer, at til den organiske Fuldmodenhed endnu en vis
først senere indtrædende sjælelig Modenhed maa komme og oppebies, paa
det at Manden dog kan være vidt nok for at kunne være Mand og Fader, og
Qvinden for at kunne være Kone og Moder. Vel lære Forældre først af
deres Børn at opdrage Børn, men de maae dog være saa vidt, at deres
Børns Egensind ei kan gjøre dem selv egensindige. Men desværre faaer
ogsaa her ei sjeldent den Forfængelighed, der ruinerer saa Meget, ofte
ogsaa en urolig Utaalmodighed, sin forbærvelige Indflydelse. Naturens
tidligen opbrusende Bevægelser og Tilskyndelser ansee man ei strax for
en Naturens Stemme; thi at denne Natur overhovedet hos de Fleste kommer
alt for tidligt frem, længe førend endog den fulde Væxt er opnaaet, er
bekjendt nok. Dog ogsaa for Elskende i senere Aar er det vigtigt og
velgjørende, ei at haste med en saa vigtig Forbindelse. Og det af en
Grund, der er høist vigtig for alle Individer: at de først maae lære at
leve med hinanden og at skue hinanden ret ind i Sjælen og fulde kjende
hinanden, for at de i Livet selv kunne prøve om de ogsaa høre til
hinanden, eller om den i Ungdommen, ofte ogsaa i den senere Alder, let
opvakte og henrevne Drivt og Lyst maaskee har hildet dem ind i en
Vildfarelse.

Hvad Mennesket egentligen i alle Livets Tilfælde skulde have, og hvad
vi, naar vi ikke ere saa gjenfødte, at vi i alle kunne have det, dog i
alle vigtige Livets Tilfælde skulde med rolig Hengivenhed oppebie, men
naturligvis ei oppebie under en dorsk Ligegyldighed, men under levende
Henvirkning dertil, --- Dette er det især saa vigtigt at oppebie, inden
et Egteskab sluttes: den høiere Dom, hvormed en af vort Inderste
kommende Stemme, der ei er vor egen, siger det at være godt, hvad vi
have fore, eller ogsaa misbilliger det. Netop da, naar Individets Lyst
og Beslutning finder denne indre høiere Samstemming, er, men ogsaa da
først, Det skeet, hvad det nys anførte Ord saa smukt kalder, at
Egteskabet er sluttet i Himmelen. Hvor vilde Menneskene være lykkelige,
om de overalt og i Alt saaledes havde Gud for Øie, at de i daglig Tanke
paa Gud og Bøn til ham vare vante til at bringe deres Anliggender for
ham, og tage Livet som en Herrens Tjeneste, hvori det kom an paa at
ville, hvad han vil. De vilde da vistnok langt mere være, hvad
Menneskene saa ofte ikke ere, og hvad det er vanskeligere at være, end
det kan synes, sandrue og redelige mod sig selv og med sig selv. Kun
altfor ofte ile Menneskene utaalmodigen fremad, uden at samle Sindet og
at gaae ind i sig selv, og redeligen at gjøre sig Rede for deres Hjertes
og Sinds Bevægelser, ja have, selv naar en indre Følelse bestemt advarer
dem, ikke Mod til at høre paa den, ikke Mod til at være oprigtige mod
sig selv. Imidlertid gives der vistnok Tilfælde, da et eneste Moments
usigelige Vished saaledes kan have Vidnesbyrdet om sin Sandhed i sig, at
et Øieblik kan decidere for et heelt Liv. Men disse høist geniale
Øieblikke ere saa sjeldne, som al den Genialitet, der forudgriber Tiden.
Og netop i et saadant Øieblik maae de Tvende, der skue hinanden dybt ind
i Sjælen, føle sig den Evige usigelig nær, idet de føle sig hinanden
usigelig nær; netop i et saadant Moment maa Hjerternes Forening være en
høist religiøs. Og ogsaa endda er det nødvendigt at bie og lade det nye
Liv komme til at sætte sig og komme til en høiere Klarhed. Alt paa
Jorden, endog det høieste Aandelige, naar det træder ind i jordiske
Væsener, er undergiven Væxt, og hvad der undfanges og fødes i Aand og
Hjerte, bliver ligesaa lidet fuldbaaret paa een Dag, som Mennesket selv
blev det. Man betænke, at Livet ogsaa i aandelig Henseende, ligesom det
har sine Kriser, saa og har sine Stadier, og at det er farligt,
besindighedsløst at ile hen over noget. Det er intet Ringe, der skal
skee, naar tvende til en vis Grad fastnede Individualiteter skulle
løsnes for at smelte sammen med hinanden. Hos enhver af dem maa i visse
Henseender den fastere Form opløses, hvori enhver maaskee allerede havde
sat sig fast, for at enhver kan aabne sig for det nye Livs
gjennemtrængende Varme, og begge ret kunne blive til Eet. Det bekjendte:
corpora tantum soluta agunt, kan ogsaa overføres paa det Sjælelige. Men
denne Løsning skeer ei sjeldent lempeligere, mildere og beqvemmere,
inden det faste Baand knyttes, naar Sindet endnu er friere, og ingen
Nødvendighedsfølelse hilder eller ægger. Ingen stole paa, at Det, hvad
der skal komme, men endnu ei er kommet, nok kan komme efter
Forbindelsen. Nei, netop fordi det da skal komme, kommer det undertiden
vanskeligere og langsommere, og ogsaa her kommer det gamle Ordsprog
igjen: at Tiderne ere lige lange, men ikke lige nyttige. Dog er Dette
vistnok ikke i lige Grad anvendeligt paa alle Individualiteter; og naar
midt i Vaaren Sommerens Hede og Lummerhed vil indfinde sig, saa blive
det lige saa gjærne Sommer strax. En meget lang Forlovelsestid er
vistnok endnu mindre ønskelig, medmindre den begyndte i den tidligste
Ungdom. Elskovens første Fyrighed og friske Duft maa for alting ikke
være forbi, naar Egteskabet begyndes; og hellere bevare man sig den ved
Nøisomhed og Maadehold, end at have fortæret det Bedste, inden at det
bedste Tidspunct kommer. Men altid synes det dog at være saare vel, og
skulde aldrig forsømmes, hvad endog Kirkens Bestemmelse forhen har
fordret, at Baandet ved en høitidelig Act knyttes foreløbigen ved en
Forlovelse eller Trolovelse, inden det knyttes tilfulde. Kun maatte
Mellemrummet imellem begge Tidspuncter ikke, som det forhen hos os ved
Love skeede, indskrænkes til en vis ufravigelig Korthed, men netop
udvides til et Tidsrum af en passende Længde, inden hvilket Egteskabet
ei maatte finde Sted. Og den foreløbige Forening maatte kunne hæves ved
den Enes Begjæring alene, fordi den netop betragtedes som en Sjæleprøve.
(En bekjendt Skik i mange Egne af Tydskland og andre Lande,
Vinduebesøgene, kan kun ifølge lignende Betragtning være bleven hævdet
ved den offentlige Mening, der dog ellers paa langt mindre sædelige
Steder ingenlunde pleier at tilstede noget Usædeligt). Den Solennitet,
hvormed Trolovelsen blev indgaaet, den solenne Maade, hvorpaa en
Forlovelse endnu ofte declareres, knytter begge Elskende fastere til
hinanden end forhen; hvad de forhen attraaede, have de nu derved
allerede foreløbigen faaet i Besiddelse, og idet deres Forhold nu
kjendes og respecteres i det selskabelige Samqvem, hvortil de høre,
kunne de nu ogsaa prøve sig selv med Hensyn til dette. Allerede for
hiint første Skridt have ei sjeldent De, der inderligt elske hinanden,
en Slags Angest, som om Det, der i deres skjulte Eensomhed var saa sødt,
nu skulde lide ved at gives til Priis for en af saa mange Lidenskaber
eller Sindsstemninger bevæget Mængdes just ei altid varsomme, stundom
forstyrrende Indflydelser. Derfor unde man ogsaa unge Elskende en Tid
lang det Skjul, under hvilket den opvaagnende Elskov saa gjærne vil
dvæle, og saa gjærne drømmer sig ukjendt. Men da nu en Declaration dog
maa skee, fordi Elskendes Samliv dog engang maa frem for Dagen, og
ligesom maa vove sig ind i Livets urolige Tummel, saa skee den dog først
foreløbigen, inden den skeer slutteligen, eller inden den knytter et
Egteskab! Det foreløbige Baand er dog allerede et Baand, og naar denne
Art af Forbindelse varer saa længe, som den altid burde vare, kan den
blive en ret god Prøve. Det er bedre, at Forlovelsen bliver
Kjærlighedens Grav, end at Egteskabet skulde blive den.

Der gives mangen egoistisk Natur, der, attraaende sin Elskovs
Gjenstand, i lidenskabelig Bevægelse synes, for at tiltvinge sig den, at
kunne opoffre Alt, synes at elske med en Fyrighed, en Energie og
Fasthed, der spaaer det Bedste; men dog snart efter, naar den trygge
Besiddelse er tilvunden, viser sig anderledes. Undertiden skeer det, at
et saadant Individ efterhaanden lunknes og bliver ligegyldigt ved
Opnaaelsen, fordi det blot tog den som en anden Trangs og Lysts
Tilfredsstillelse, altsaa som en Overgang, og ikke som et dybt i
Mennesketilværelsen liggende Savns varige og blivende Fyldestgjørelse.
Ofte derimod see vi, at den egoistiske Natur, der en Tid lang holdtes
nede ved Elskovens Overvælde, kommer op igjen og griber om sig i selve
Elskoven, saa at Individet efterhaanden begynder at plage den Elskede
med sit egensindige, herriske eller selvraadige Sind. Det ere ikke Alle,
der kunne taale Lykken, taale et med Heftighed attraaet Godes trygge,
rolige Besiddelse. Endog den blotte Tanke, at Baandet er et Baand,
plager mangt et Sind; ja det hændes vel, at Den, der viiste stor Omhu og
Opmærksomhed, begynder efterhaanden at efterlade den, blot fordi han
hildes ved den Tanke, at han er ligesom forbunden dertil, medens han dog
selv igjen begynder at betragte som en Ret og Fordring, hvad han forhen
med Glæde og med Elskovs Erkjendtlighed blot tog som et sødt Beviis paa
den Gjenkjærlighed, hvorved han glædede sig saa meget. Ja, saa nær ligge
os de egoistiske Bevægelser i vort Indre, saa snare ere de til at komme
frem i Menneskehjertet, at endog den bedre Individualitet stundom kan
finde sig selv besnæret af lignende, snart utilfredse, snart kolde, ja
ligesom frosne Følelser, næsten Øieblikket efter, at den fuldeste Lykke
og Glæde synes indtraadt. Sligt vide man da, og ændse den uhyggelige
Sindsstemning ikke, til den gaaer over. Men ere Naturerne ret
egoistiske, opkommer let megen Strid, og begge støde desto stærkere imod
hinanden, jo heftigere Tiltrækningen var. I alle Tilfælde er det en
ængstelig Existents, en Existents, der let kan blive ruinerende og
ødende for Aand og Hjerte, at leve under idelig Strid og Disharmonie med
hinanden, være endog Elskoven virkelig heftig og stærk, og følge ogsaa
paa enhver Strid en Elskovsudsonelse. Derfor bør Elskende under en
foreløbig Forbindelse først lære, om de kunne leve med hinanden i den
uadskillelige. Netop først naar Forlovelsen er indgaaet, netop naar
Attraaen har funden sin første Hvile, kan det vise sig, om Forholdet har
Bestand nok, for at det næste Skridt bør voves. Hvo der, indseende, at
Tvende ei høre sammen, ønsker at skille dem ad paa den lempeligste,
mildeste Maade, handler stundom klogt i at give dem sammen ved en
foreløbig Forbindelse, paa det at de nu selv kunne lære at faae Øinene
op og skille sig selv ad; hvorimod en ved et Magtsprog nedslaaet
Tilbøielighed beholder sin Braad og ikke let kommer til Hvile.

Imidlertid, naar endog kun den foreløbige Forbindelse hæves, er Dette
allerede et Brud, og kan gjøre en dyb og smertelig Rivt i Existentsen.
Og dog, dersom ingen Tilknytning skal turde ende med en Gjenafskillelse,
hvor kunne da ret paalidelige Forbindelser komme i Stand, Forbindelser,
som ei blot Øinene, men og Hjerterne af Hjertet knyttede? Derfor være
overhovedet Omgangen mellem begge Kjøn en ret fri, og den Frisindighed,
den Liberalitet i Tænkemaaden herske, der tilsteder, at vakte Haab tør
blive til Intet. Thi hvorledes ville de Følelser, der kunne føre til
Noget, kunne faae Rum til at gribe om sig, naar den Frihed ikke hersker,
at træde tilbage igjen, hvor meget man end har nærmet sig. Mangen Gang
skeede det vel, at Den, der var kommet over den Alder, hvori, hvad der
skal skee, skeer aldeles af sig selv, ikke kom til en Forbindelse med et
andet Væsen, blot fordi han ei vilde vove at vække uvisse Haab, og dog
ei kunde komme til at give et trygt og tilforladeligt, uden at udsætte
sig for at vække det uvisse først.

\begin{center}\rule{0.28\linewidth}{0.4pt}\end{center}

De, der i et heelt Liv skulle leve med hinanden til Velbehag for Gud og
Mennesker, maae mere end holde af hinanden, de maae ogsaa stemme sammen,
de maae inderligen forstaae hinanden og svare til hinanden, og dernæst
ogsaa selv vaage over, at ret dem forundte Gode bliver dem et sandt
Gode.

Grunde, hvorfor Tvende, der hver for sig ere brave og holde af hinanden,
dog ei høre sammen, kunne der i individuelle Tilfælde være mange af. I
Almindelighed ere her trende Hovedpuncter at have for Øie: Kjønnets
Forskjellighed, Alderens og Standens, forsaavidt denne medfører en
Forskjel i begges Cultur og hele aandelige Tilværelsesmaade.

Hvad Kjønnets Forskjellighed angaaer, er vel at overveie, at Manden, saa
meget end hans Kone ellers i Aandsfortrin monne være ham overlegen,
nødvendigvis i visse bestemte Henseender maa have en bestemt Overvægt,
hvis ikke netop Konen skal føle sig ulykkelig under Begges Samliv. Thi
netop for Konen er det en Ulykke, at have en Mand, hun i alle Henseender
kan og tør, vel stundom endog maa, oversee. Ethvert dybt, ædelt og
fiintfølende qvindeligt Hjerte vil sikkert langt hellere vælge en
Tiendedeel af en i Sandhed mandig, kraftig og aandsfuld Mands Elskov,
naar hun kun er fuldkommen og for bestandig vis paa denne ene
Tiendedeel, end al den udeelte Elskov og Tilbedelse, den Mand kan yde
hende, hvem hun maa beherske for at kunne leve med ham. Der er Noget,
Qvinden maa have Respect for hos den Mand, hvem hun skal tilhøre, være
der end meget Andet, hvori han tilstaaer hende at have Fortrinet. Hvad
dette Noget er, er vanskeligt at udtrykke i almindelige Ord, saa bestemt
man end i enkelte givne Tilfælde fornemmer det. Skulle vi udtrykke det,
saa tør vi maaskee sige, at det er hans Villie, betragtet i det Hele som
Villie, eller som Personlighedens inderste og bestemteste Udtryk: det er
hans mandlige Selvstændighed. En Kone, som hører op at betragte sig selv
som Den, der er til for at tjene sin Mand, ja endog kommer til en vis
Følelse af Ringeagt eller Tilsidesættelse imod ham, er en høist
ulykkelig Kone. Alt, hvad der giver Elskov og Elskovsforholdet den Dybde
og Sødme, der endnu i langt høiere Grad synes beskjæret Qvinden end
Manden, er borte, og et unaturligt Forhold er indtraadt. Dette føler
ethvert nogenlunde ordentligt og sundt følende Qvindehjerte. Selv den
Grad af Indflydelse, Konen som oftest ikke blot har, men ogsaa bør have,
og hvilken netop de sindigste og mandigste Mænd lettest tilstaae, ---
selv denne hører hun ikke gjærne nævnet eller bemærket; som det da ei
heller er fornødent, at hun er sig sin Indflydelse bevidst i samme Grad,
som hun udøver den. En dyb Respect for Manden, en vis
Underdanighedsfølelse, ligger i ethvert ædelt qvindeligt Væsen. Den
fiintfølende Qvinde, der agter paa denne Følelse, fornemmer vel dens
Betydning og dens Dybde, og pleier og nærer den med Lyst. Clärchens
Situation ved Egmonts Fødder hos Goethe er derfor for et følelsesfuldt
Qvindehjerte en høist tiltrækkende Stilling; og en uddannet Aands
herligste Rigdom hindrer ikke i, men medfører meget mere, at et
kjærlighedsfuldt qvindeligt Gemyt føler sig dybest i Sjælen
lyksaliggjort ved at kunne vise en høiagtet Mand, hvor ydmyg og
underdanig hun af Hjertet føler sig imod ham. Dog ere disse Følelser af
den ømme og fine Art, at de kun i sjeldne lykkelige Øieblikke kunne
komme til at tale. Men for at Qvinden kan nære disse dybe Følelser mod
Manden, maa hun vistnok ikke i Aandsfortrin være ham altfor overlegen.
Og her ligger da en af Kjønnets Forskjellighed alene udspringende
Aarsag, hvorfor Tvende stundom ei høre sammen, uagtet hver af dem ret
godt kan passe for Andre.

Hvad Manden angaaer, saa er for ham den første Fornødenhed naturligvis
den, at Qvinden er af Hjertet beredt til at være Qvinde, hvortil hører
at være huuslig; thi heri ligger en stor Deel af Qvindens Kald i
Samlivet; hun maa kunne betragte sig som en Tjenerinde for sin Mand og
for den opvoxende Slægt. Iøvrigt forstaaer det sig af sig selv, at her
ikkun tages Hensyn til Det, der er at see paa hos Individer af begge
Kjøn som Egtefæller. Hvad der overhovedet udkræves for at være
retsindige og brave Mennesker, derom er Talen ikke her. Af herriskt
Sindelag maa intet Menneske være, mindst en Qvinde, der skal være
Egtefælle og Moder.

Af Standsforskjelle er der kun een, der er saa væsentlig, at en
Forbindelse, ved hvilken denne Forskjel oversees, naar særdeles
Undtagelser fraregnes, --- et udmærket Individ eller særdeles inderlig
Tilbøielighed kan medføre en grundet Undtagelse, --- bliver en sand
Mesalliance: det er den Standsforskjel, der tillige er en Forskjel i
Henseende til Opdragelse, Dannelse og hele Maaden, hvorpaa Livet
betragtes og nydes, en Forskjel, der fører med sig, at den Ene ikke kan
gaae ind i den Andens Ideekreds, at tilmed imellem Begge ingen sand
aandelig Meddelelse kan finde Sted. Imidlertid er denne Standsforskjel
ikkun en reent menneskelig, ikke en borgerlig. I alle borgerlige
Stænder kunne findes Dannede og Udannede, ædlere og uædlere
Individualiteter. Den sande til alle Tider gyldige Adel, der ogsaa altid
vil haandhæve sin Indflydelse i Staten, kan ethvert Menneske af
Naturanlæg give sig selv, og ethvert Menneske af Menneskesands give
idetmindste sine Børn, naar disse ei for meget fattes Anlæg. Alle
Forskjelle i Henseende til en blot borgerlig Stand betyde kun Lidet,
hvor ægte Elskov findes, især i vor Tid og vort Fædreneland, hvor man
vel neppe, naar enkelte Tilfælde undtages, vil ansee en egteskabelig
Forbindelse imellem Tvende, hvoraf den Ene hører til de høieste Classer
i Staten, den Anden til en langt lavere, for en Mesalliance, eller for
noget Paafaldende, naar kun Begge ikke ere for ulige i Aand og Dannelse.
Det er ogsaa, ikke blot overhovedet fornuftigt, men og i Særdeleshed
klogt, at de høieste Stænder benytte den Frihed, hvilken Tidsaanden har
medført, til at udvide Kredsen for deres Valg, og stundom at forfriske
deres saa ofte temmelig henvisnede Stammer med friskere Safter, og
derfor at nyde godt af den sundere Menneskeexistents, der dog endnu
bestandigen fornemmelig synes at have hjemme i den dannede Middelstand,
især hvad netop det huuslige Liv angaaer. Ikkun Fyrster maae i vor Tid
endnu, ligesom i saa mange andre menneskelige Henseender, saa og i
Henseende til en skjøn og reen Elskov og den Forbindelse, hvorefter den
stræber, lide under deres ophøiede Stillings mange Ubeqvemmeligheder.
Endnu kan ikke let nogen Fyrste handle som et naturligt og dygtigt
følende Menneske i sin Families Kreds; ja saaledes, som Tingene have
staaet, er det næsten et Under, at vi i vor Tid have i Kongeslottene
seet Egteskaber saa lykkelige, som kun Nogen i Borgerhusene kan ønske
sig dem. Det er vistnok i denne Henseende, som i saa mange andre
Henseender, en stor Lykke, at der dog i Europa er eet Land, hvori en
Mængde smaae Fyrster have beholdt deres Fyrsteværdighed. Dog vil vel
engang den Tid komme, da en Fyrste af Aandsklarhed og Aandskraft vil
udfinde Middelet til at forskaffe sig, hvad ethvert sundt Menneske
trænger til, og at lære Verden, at man ei behøver at betale den Byrde at
være Fyrste med den Forret at være udelukket fra at nyde en reen
Menneskeexistents. Iøvrigt er det ganske naturligt, om den lavere Stand
stundom endnu stærkere skyer Forbindelsen med den høiere, end om denne
stundom er uvillig til at tænke sig i Forbindelse med hiin. Thi den
lavere har her mere at vove. Især gjelder Dette, naar en Pige af lavere
Stand skal forbindes med en Mand af den høieste; hvorimod det er mindre
betænkeligt, at en Pige af høi Adel forbindes med en Mand af
Middelstanden; forudsat iøvrigt, at de Tvende indbyrdes ere fuldkommen
visse paa hinanden. Thi da Manden drager Qvinden hen til sin Plads langt
mere, end han nærmes til hendes, saa er i det sidste Tilfælde
Bortfjærnelsen fra den sundere Middelstilling for Den af de To, der
hørte til den, ei saa stor, som i første Tilfælde.

Endnu er Alderens Forskjel at tage i Betragtning. Een saadan Forskjel
møder her, som altid er betænkelig, den, der medfører en Forskjel i
Henseende til den Periode i Livet, hvori begge med Hensyn til Elskovens
Følelser staae. Naar den Ene allerede har prøvet og nydt Elskovs hele
Fylde og Herlighed, eller dog paa een eller anden Maade er kommen
saavidt, at Elskoven ei mere hos ham kan have den ungdommelige
Enthusiasmes Fyrighed og Livelighed, men nu det andet Individ derimod
endnu staaer i Livets første Vaar, da kan det ofte med Hensyn til dette
sidste Individ være ret tungt, at Forbindelsen er saa ulige fra begge
Sider; thi da fører det let med sig, at den Yngre ei kommer til at finde
den ungdommelige, livelige, friske, frodige Elskovs Vaar, den dog var
bestemt til at nyde.

\begin{center}\rule{0.28\linewidth}{0.4pt}\end{center}

Uagtet al Kjærlighed, og tilmed ogsaa Elskov, efter sin indre Natur
vækker Aandens høiere Øie og gjør seende, saa kan den dog ogsaa ofte
blænde og gjøre, at Mennesket ikke seer, hvad der er at see, fordi den
Elskende saaledes fæster Øiet paa Det, der især tiltrækker ham, at han
overseer det Øvrige: en Eensidighed, hvori Elskende ofte endog med Lyst,
ja ligesom med Flid sætte sig fast, som om det var en religio, en hellig
Nødvendighed, for dem, at finde Alting skjønt og indtagende hos
hinanden. Men ligesom Elskov i det Hele gjør seende, saa stemmer det
ogsaa fuldkomment med dens Natur at gjøre seende i det Enkelte. Netop
den dybe Omhu for et andet Væsen, den inderlige Lyst til at vaage over,
at dette kan have alt paa det Bedste, aabner ogsaa Øiet for Manglerne;
hvorfor da og en ægte kjærlig og omhyggelig Moder pleier at have et
skarpt Øie for sine Børns Feil. Elskovens Inderlighed og Styrke vil ei
strax lide ved denne Erkjendelse af den Elskedes Mangler; den vil ei
lide derved, naar netop denne Erkjendelse tages som Noget, der forhøier
Mildhedens og Overbærelsens og en vis sjælelig Baretægts og Omhues dybe
og søde Følelser. Et kraftigt følende Væsen vil endog ikke, at det
skulde synes for den Elskedes Øine at være Mere eller Andet, end hvad
det i Sandhed er. Paa Sand bygger man ei engang en Hytte, end sige et
Pallads, og det vilde dog Den, der vil bygge sit hele Livs Lykke paa en
Vildfarelse. Derfor er det just Dem, der føle med ægte Kraft og Klarhed,
magtpaaliggende, ei at være deres Elskede Mere, end de i sig selv ere.
Sandhed er her, ligesom overalt, det alene Paalidelige, og at være
sanddru baade mod sig selv og Andre er, ligesom Gudsfrygt, nyttigt
overalt. Og netop under Elskovens ungdommelige Enthusiasmus er det
bedst, at Begge med reen Oprigtighed see, hvad der er at see og engang
vil blive seet; netop da er det allerede Tid, at den Ene er oprigtig for
den Anden og imod den Anden. Thi netop da lære Begge bedst at finde sig
i hinandens Mangler. Netop da kunne de, om ei lettest, dog med mindst
Smerte og Misnøie, lære, hvad de maae lære, at det herlige Gode, der er
skjænket dem i deres Elskov, er et jordiskt Gode, og at de, i hvor
naturligt det er, at den Ene er som en Gud for den Anden, ei for heftigt
maae ville, at det virkelig skal være saa. Vel overtrædes intet Bud
lettere og oftere end det første Bud; ja i Elskovs første Varme
overtrædes det endog paa en smuk hjerteqvægende og henrykkende Maade,
idet her i en vis Henseende virkelig den Ene bliver, hvad saa ofte
Mennesket maa være Mennesket, en Gud for den Anden. Men man vogte sig
dog endogsaa her! Vi see ei sjeldent heftigt elskende, men derhos
heftigt fordrende Individer, der med en sælsom Forvirring yttre sig og
handle, som om de vilde, at den Elskede virkelig skulde være, hvad den
syntes dem, et overjordiskt Væsen. Til Egteskabet skride Ingen med store
Forventninger eller Fordringer, men lære først ret at fatte og huske
paa, at vi alle ere Herrens Skyldnere, og alle med hinanden ere paa
Veien til hans Domstol, og alle haabe der at finde Naade. Og Enhver, der
underveis tager sig en uadskillelig Ledsager, tage denne, som den er,
med alle Mangler, og skjule disse ikke for sig.

Dernæst betragte man Egteskabet ei blot som en Lykke, der skal nydes,
men og som et Kald, der skal tilfredsstilles. Især er det vigtigt for
Qvinden at gaae ind i Egteskabet, som en retskaffent følende Mand gaaer
ind i det Kald, Statssamfundet tildeler ham, hvis Tjener han betragter
sig at være, og for hvis Skyld han anseer sig at existere, medens han
kalder det en stor Lykke at have fundet det, naar det svarer til hans
Hjertes Hu. Hvad dette Kald i Statssamfundet er for Manden, det bringer
Egtestanden Qvinden: den sande af Naturen bestemte, til en complet
Existents fornødne, Virkekreds. Og ingen Qvinde høre op med, naar hun nu
virkelig træder ind i denne Virkekreds, at betragte sig som Det, hvad
den ømt og dybt følende Elskerinde saa gjærne drømmer sig at være,
Mandens omhyggelige Tjenerinde i Alt. Ogsaa glemme Ingen af dem, at
Kjærlighed ikke blot constitueres ved en inderlig Tiltrækning, men og
tillige ved den Fjærnhed, hvorved Begge holdes ude fra hinanden, og
hvortil den Blyhed og Fiinhed svarer, der i Forbindelsen maa holdes i
Hævd. Agtelsens og den til denne hørende Opmærksomheds Følelser bevare
sig derfor og vorde nærede! Derfor glemme de Forbundne ei heller, at de
have Pligter mod hinanden! Ordet Agtelse og Ordet Pligt klinge vistnok
for et inderligt og stærkt følende og elskende Væsen ofte som sælsomme
Udtryk, der kun synes at have lidet hjemme og at være uden Betydning,
hvor Alt ved de sødeste og saligste Følelsers uimodstaaelige Styrke
synes at skee af sig selv, ei blot uden Modstræben, men og uden
Paamindelse. Det er ogsaa godt, og netop som det skal være, naar saavel
her, som overalt i hele Livet, et høit, helligt Livs ustandsede Strøm og
Aande saaledes bevæger og besjæler det jordiske Væsens Hjerte og dets
udadgaaende Virken, at nu hele Livet af sig selv leves, som det bør, med
fuld Kraft og inderlig Nydelse. Men i vor jordiske Tilvær overhovedet
ere, og ligesaa ere i Egteskabet, ei alle Øieblikke gudopfyldte eller
aandfulde og livfulde, endog for et gudopfyldt eller aandfuldt Individ.
Der komme matte Øieblikke, der kommer al Slags Misnøie frem under Livets
Uro. Da er det godt, at i Egteskabsforholdet, ligesom i Livet
overhovedet, Pligtfølelsen træder til, værnende og bevarende, paa det at
dog ei de matte Øieblikke skulle endog forvanske eller opholde de senere
aandfuldere, ja forbærve hele Forholdet. Mennesket maa overalt være
aarvaagent og vogte sit Hjerte og vaage over sin Elskedes tillige og
bære Omhu for dette. Egtefæller skulle jo bære Livets Byrder og Sorger
med hinanden, ja selv disse see unge Elskende med Lyst og Mod imøde,
fordi ogsaa disse skulle virke til at knytte dem fastere og inderligere
til hinanden; thi Dem, som frygte Herren, maa Alt komme til Gode. Men
til Det, Egtefæller skulle bære med hinanden, høre ei blot Livets
vigtigere Sorger og den Nød og Fare, mod hvilket et modigt Sind netop
føler sig opfordret og oplivet til kraftig Virksomhed eller Udholdenhed;
men dertil høre ogsaa Livets smaae Drillerier og Ærgerligheder eller
Kjedsommeligheder, af hvilke et Gemyt kun alt for ofte lader sig opægge
til en vis Egensind eller Uregjerlighed, der, naar det ei er aarvaagent
over sig selv, let kan gribe om sig. Hine større Sorger bæres lettere,
end disse Smaating. Mangen Kone finder sig lettere i at vaage syv Nætter
ved Mandens Seng eller ved et Barns Vugge, end hun finder sig i, at en
Naal er kommen bort i et ubeleiligt Øieblik. Men det var dog for Meget,
om en bortkommen Naal skulde skille dem ad, der have holdt Uger ja
Maaneder ud ved hinandens Sygeleier. Derfor glemme de ikke, at de skylde
hinanden, hvad Enhver af dem føler at skylde det første det bedste
Menneske, der kommer ind ad Døren! Men hvad de skylde hinanden, yde de
hinanden da, som de, der yde det af inderlig og aarvaagen Omhu for
hinanden! Idelige smaae Irriteringer kunne tære mere paa Livets og
Elskovens Fylde, end Nogen i Øieblikket troer. Og De, der ofte irritere
hinanden, vide sig jo aldrig sikkre. Men eet Sted idetmindste maa Enhver
have, hvor han veed sig sikker og føler sig vel, at han altid kan med
Lyst ile derhen, og dette ene Sted maa være hans Hjem.

At den Forening, hvortil Elskov fører, ei skal være en Forening i denne
eller hiin Henseende alene, men i alle Henseender, og at den ei skal
medføre et partielt, men et totalt Samliv, Dette føler i Ungdommens
skjønne Periode enhver lykkelig Elskende, og medens han selv stræber
efter at tilegne sig den Erkjendelsens og Sjælenydelsens Fylde, der har
opladt sig for ham, og hvorved et opvakt og ædelt ungdommeligt Sind
glæder sig saa meget, attraaer han tillige af al Magt at leve og nyde
det rige Liv med sin Elskede. Begge finde da tusinde Berøringspuncter,
og tusindfoldig Meddelelse finder Sted imellem dem, saafremt de kun
ellers høre sammen. Men endog De, hvis Elskov faldt i denne skjønne
ungdommelige Tid, og endnu oftere De, der først senere fandt Elskovs
Lykke, glemme saa ofte, naar den fulde Forening er indtraadt og Livets
øvrige Krav og Kald indfinde sig igjen, den lykkelige Skikkelse, deres
ungdommelige Elskov havde. Alt for ofte bliver endog Elskov taget som en
Overgang: de tvende forbundne Individer træde efterhaanden i sjælelig
Henseende tilbage i deres forrige Maade at existere paa, og komme Tid
efter anden saavidt, at begge tilsidst leve, hver for sig, deres
isolerede Liv ved Siden af hinanden, istedetfor at leve et, begges hele
Liv omfattende, Samliv med hinanden. Istedetfor at de skulle stræbe at
gaae ind i hinandens Individualiteter og knytte sig under mangfoldig
Meddelelse til hinanden i mange Puncter, indskrænker sig Samlivet til
nogle faa Puncter, hvori begge føle sig som sammenvoxede; udenfor disse
er det, som om de ei hørte sammen. Elskende ere paa mange Maader og i
mange Henseender Incitament for hinanden. Men til hvad Side Elskoven
fornemmelig kommer til at vende sig, beroer paa, hvorledes de tage den.
Og saa skeer det da, at ofte stor Eensidighed faaer Overhaand, og at
hele Regioner i deres Samliv komme til at ligge øde. Ofte indtræder
denne Tilstand nu vistnok, fordi begge i sig selv virkelig ikke hørte
sammen, fordi den ene ei kan gaae ind i den andens Ideekreds; stundom
derimod af en vis overstor passiv Tilfredshed i Begyndelsen, og af
Ladhed eller dog Mangel paa Opmærksomhed og rask Aandskraft siden: en
Ladhed eller Inertie, der endog ofte fører til at bornere Omdømmet og
det frie Blik paa Livet, saa at de Elskende endog tabe Sandsen for det
mangfoldige Herlige, der, forskjelligt fra Det, de skue i hinanden,
bevæger sig omkring dem, og nu derfor tabe Ansporelsen til egen
Fremadstræben og Udvikling. Imidlertid give de sig da dog gjærne hen til
en vis Art af de Bevægelser, hvori Elskov sætter. Der gives Mange, hvis
hele Elskov synes at blive staaende ved den Lyst, hvormed de beskue
hinanden og gjøre af hinanden; og nu, for dog at bringe en vis
Mangfoldighed og Afvexling ind i Forholdet og fordrive en
Kjedsommelighedsfølelse, hvilken Kjærtegn ei altid kunne forjage, søge
de idelige nye Opvækkelser, der dreie sig om det samme Punct, og Aanden
kan derover ei komme op. Heri forsynde sig mange Elskende og glemme, at
selv Eros's Tjeneste ei maa vanhelliges, glemme, at Elskov skal føre til
en Sjælemeddelelse og et aandeligt Forhold, og tabe nu Livets uendelige
Nydelse over en høist indsnævrende, der nu ei engang ret bliver en
Nydelse af Hjertens Grund, og let kan blive en aandsfortærende. Eller og
skeer det, at man lader Livets Torne og Tidsler faae saa megen Overhaand
i det aandelige Samliv, at de tilsidst underkue det fremspirende Frø.
Undertiden var vel ogsaa et af den i saa mange Henseender ødelæggende
Egoismus udspringende Overmod, en selvraadig og herrisk Lyst, til
Fordærvelse her, naar Den, der havde givet sig ganske hen, dog nu var
for stolt til at ville give sig ganske hen. Men ei sjeldent lider endog
et fra først af frodigt Samliv, og falder eller visner bort
efterhaanden, fordi der fattes en tilstrækkelig Næring, indtil tilsidst
Opmærksomheden paa den Næring, der endda kunde frembyde sig, ogsaa
svækkes. Omgivelserne bidrage da stundom ei lidet til at forstørre
Kløften mellem Begge. Der vises ofte for ringe Opmærksomhed eller for
liden Sands for Egteskabets Hellighed, og betænkes ei altid hvad hver En
ogsaa i denne Henseende skylder Andre. De kjæreste Mennesker ere os ei
altid kjærkomne, og kunne ei være det, naar de træde forstyrrende ind i
vor Kreds. For at nære et ægte Samqvem med hinanden, udfordres blandt
Andet ogsaa undertiden Eensomhedens Ro. I en Kreds, hvori et udbredt
selskabeligt Samliv let danner sig, lide ikke sjeldent de
fortreffeligste Mennesker meest i denne Henseende, fordi de
fortreffeligste snarest indeslutte saadanne Følelser dybt ind i sig
selv, der ved at komme frem til Dagen let kunde profaneres.

Egtefæller leve overhovedet alt for ofte ikke nok med og for hinanden.
Men for at et fuldeligt sjæleligt Samliv skal kunne finde Sted mellem
dem, maa Samlivet vistnok have et ydre Substrat, en Basis, hvorpaa det
kan hvile, og hvoraf det kan hente Næring. Hvad der udfordres til alt
Liv, er ogsaa fornødent til dette Samliv. Men alt Liv overhovedet træder
frem som levende derved, at det træder virkende og kraftigt ind i noget
Andet, end det selv, hvilket Andet nu bliver Basis eller Substrat for
dets levende Yttringer. Saaledes træder Sjælen, som levende Princip,
frem i Legemet, idet den former, uddanner og vedligeholder dette Legeme,
hvilket imidlertid dog er noget ganske Andet, end Sjælen, der udgjør den
uafbrudte Leven deri. Saaledes træder da ogsaa Alt, hvad der levende
bevæger sig i Menneskets Aand og Hjerte, den hele Fylde af høiere
Indflydelser, der under Ideers og Følelsers Form røre sig i Menneskets
Indre, alt Dette træder frem i Mennesket i et Indbegreb af Forestillinger
og Sindstilstande, hvilke ere noget ganske Andet, end hiin indre Fylde
af aandelig Leven. Forestillingerne og Sindstilstandene ere det
Øieblikkelige, Vexlende og Ydre, hvilket kun er og betyder Noget ved
hiint indre aandelige Væsen, som træder frem deri. De ere det Substrat,
det Grundlag for det aandelige Liv, i hvilket dette, ved ideligen at
fremkalde, bestemme og uddanne det, viser sin indre Kraftighed og
Væsenhed; heri træder det frem som aandelig Leven. Og saaledes maa da nu
ogsaa det Samliv, hvori tvende Elskende i en skjøn og inderlig Forening
skulle voxe sammen, have sit Substrat eller sin Basis. Ei blot i
udvortes og borgerlig, ogsaa i indvortes, reen sjælelig Henseende maa
Egteskabet have et Grundlag for sin Subsistents, en bestandig Kilde til
Næring. Der maa være Noget, hvorom Begges forenede Stræben, Virken og
Leven saaledes dreier sig, at det kan blive dem et bestandigt Vehikel for
Hjerternes Meddelelser, og ligesom et Materiale, ved hvilket Livets indre
Kraft og Gehalt bestandigen vækkes til at yttre sig. Mennesket skal i
aandelig Henseende voxe hele Livet igjennem. Denne Væxt skeer, idet
Livets indre Kraft bestandigen mere og mere uddanner den Ideekreds og det
Indbegreb af Følelser og Beslutninger, som Mennesket bærer og bevarer i
sit Indre. Men Elskende skulle voxe med hinanden. De maae altsaa i
Forening med hinanden uddanne de Ideekredse, hvori Enhver af dem
udtrykker og nyder sit indre Liv, saa at derved begges aandelige Liv paa
mange Maader kommer til at gribe ind i hinanden.

Lykkeligst er i denne Henseende Deres Stilling, hvis Kald knytter sig
saa umiddelbart til deres huuslige Liv, at Konen umiddelbart maa dele
Mandens hele Virkekreds og hans deraf udspringende Interesse og Arbeide.
At prise er i denne Henseende især Landlivet, hvori Mandens hele Virken
og Konens hele Virken gribe i hinanden. Det er usigelig sødt at arbeide
i Selskab med et elsket Væsen. Men denne Lykke, der i den lavere Stand
ofte haves, er i den høiere langt sjeldnere. Her ligger som oftest
Mandens Kald i en ganske anden Sphære, end hans huuslige Existents;
begge Egtefæller kunne vel arbeide for, men ikke med hinanden.
Imidlertid er det overalt naturligt at dele med den Elskede alt Det, der
levende bevæger og fylder Hjertet. Om derfor end Mandens Kald og Syssel
ligger endog fjærnt fra den Ideekreds, hvori Qvindens Liv alene pleier
at bevæge sig, er det dog naturligt, at han attraaer at dele med hende
den Interesse og de Følelser, der i dette Kald bevæge og opfylde ham;
ligesom ogsaa hun maa attraae at deeltage deri, for at kunne glæde sig
ved Det, der glæder ham, og inderligen kunne nyde denne Glæde med ham,
og at kunne gaae ind i hans Bekymringer og fatte dem tilfulde. Saa holde
da ogsaa Intet de Elskende tilbage fra at leve saaledes med hinanden, og
dele Alt, hvad der griber ind i det Liv, enhver af dem lever! Er end
Qvinden lidet skikket til visse mandlige Forretninger, saa kan hun dog
vel være i Stand til at fatte deres Værd og Betydning, til at forstaae
og med levende Interesse betragte deres Formaal. Overhovedet er der i
den qvindelige Naturs Eiendommelighed Intet, der hindrer i, at jo
Qvinden kan have Sands for alt Det, idetmindste i det Hele taget, hvad
Manden er istand til at erkjende, men hvad han vistnok tillige maa være
skikket til at udtænke og udføre. Egtefællers Samliv kan altsaa
forsaavidt meget vel være et ægte, begges hele Liv omfattende Samliv.

Næst efter Samfundet med Gud og en Forbindelse med andre menneskelige
Væsener, for hvilke et Menneske kan aabne sit Hjerte, er der neppe for
Den, der er kommen til en vis aandelig Modenhed, nogen aandelig
Fornødenhed mere væsentlig og mere indtrængende, end den, at have en
bestemt Virkekreds, et Kald, som svarer til hans eiendommelige Natur og
Tilbøielighed, og hvilket han er istand til at udfylde. Qvinden er i
denne Henseende saare lykkelig; thi naar hun først finder Den, hun kan
give sit hele Hjerte hen til, saa finder hun, paa eengang og i Eet,
Fyldest for tvende af de væsentligste Fornødenheder. Den Forening,
hvortil Elskov fører, aabner tillige den skjønneste Virkekreds for
hende, den, hvortil hun netop allermeest er kaldet af Naturen, den, som
er i Stand til paa den vederqvægeligste Maade at udfylde Sind og Hjerte.
Noget maa, som Schiller treffende lader Choret synge i Bruden fra
Messina, Noget maa Mennesket have at sørge for fra den ene Dag til den
anden, paa det han kan udholde det lange Liv. Men sikkert gives der
blandt denne Mennesket saa fornødne Sørgen og Syslen og Virken ikke
megen, der, naar den er, som den kan være, er skjønnere og mere
qvægende, end en Qvindes uafbrudte Omsorg for, at en elsket Mand og
kjære Børn kunne have Alt paa det Bedste. Livet er overhovedet i desto
høiere Grad en ægte Leven, jo mere det er en ægte Elsken; men at elske
er at give sig hen for noget Andet, end sig selv, og at nyde Livet og
glæde sig ved Livet i den Hengivenhed, der dreier sig om den elskede
Gjenstand. Derfor er en ægte Leven altid en ægte Kjærlighedens Tjeneste,
en inderlig Omhu for, at Det, til hvis Tjeneste Mennesket har givet sig
hen, kan nyde sin Ret og sin Fyldest og være paa det Bedste. Saaledes
elsker enhver ordentlig Mand sit Kald og føler sig i dets Tjeneste med
Lyst, og lader det være sig af Hjertet magtpaaliggende at tilfredsstille
det, nydende sit eget personlige Liv i sit Kalds trolige og hjertelige
Opfyldelse. Den Lignelse, at en Mand er viet til sit Studium og sit
Kald, at det er ham som en Brud, han inderlig elsker og drager Omsorg
for, er en Lignelse, der er af Betydning og der udtrykker Mandens
Forhold til sin Virkekreds og til den Idee, han ledes af, og for hvilken
han virker, som dette Forhold skulde være, men som det af mange Grunde
vistnok ei altid er. Men hvor skjøn er især Qvindens Stilling i denne
Henseende? Hvilken Eenhed, Holdning og Energie kan ikke hendes Kald give
hendes Liv? Hvad Manden saa ofte maa leve og besidde og nyde i tvende
forskjellige Kredse, det lever og besidder og nyder hun i een og samme
Kreds. Netop medens hun lever for sit Kald og lever for Ideen og tjener
Herren, lever hun umiddelbart tillige for de Væsener og til de Væseners
Tjeneste, der ere hendes Hjerte de nærmeste og dyrebareste. Hende er det
givet, medens hun netop opfylder sin nærmeste Bestemmelse, som Deel af
det Hele, at kunne concentrere sin Virken om den ene Kreds, hun har
givet sit Hjerte hen til. Men vistnok gjør iøvrigt det Samme, hvorved
Qvindens Lod kan være en saa heldig, tillige, at hendes Stilling i en
anden Henseende er misligere; thi ligesom begge saa væsentlige aandelige
Fornødenheder finde deres Fyldestgjørelse med hinanden og i Eet, saa
maae ogsaa begge undværes med hinanden, og det Qvindehjerte, der ei er
saa lykkeligt at kunne give sig ganske hen til en elsket Mand, maa nu
tillige forsage det ved Naturen hende bestemte Kald: en Naturbestemmelse,
der kun halvt kan erstattes ved noget Andet, det være da ved et ret
decideret Talent til et af de i Menneskesamfundets Hele indgribende
Virksomheder.

At Qvinden lever for at tjene og det at tjene personligen og til Andres
personlige Fornødenheders Tilfredsstillelse, og at denne Tjeneste er en
skjøn Tjeneste, og det en dobbelt skjøn, naar den dreier sig umiddelbart
om høiagtede Mænds Personligheder, det føler letteligen ethvert opvakt
og fyrigt qvindelige Væsen, hos hvem den naturlige Menneskelivvær ei har
lidt for meget under en forvendt Opdragelses eller uheldige
Standsforholds skadelige Indflydelse. Unge Fruentimmer af Sands og
Følelse finde stor Glæde i personlige Haandsrækninger, og lære ogsaa
efterhaanden uden Vanskelighed at erkjende, at det, da kun en ægte
Virken er en ægte nydelsesfuld Leven, er en Lykke, at deres Virken kan
og skal umiddelbart dreie sig om en elsket Kredses personlige
Velbefindende, og at Agtelsens og Hengivenhedens Følelser have saa megen
Ret og saa megen Leilighed til at udtale sig strax i Handling, under
personlig Omsorg for deres Gjenstand. Det Goethiske: „Qvinden lære at
tjene“ er let efterlevet, idetmindste efterlever den sundt følende
Qvinde det i det Hele taget med Lyst. Men da indfinde sig i den
egteskabelige Forbindelse ogsaa det Andet, han tilføier! netop ved at
tjene, komme hun til den Raadighed, der tilkommer hende i Huset! Ingen
hersker bedre og kommer bedre til et Herredømme, end Den, der først selv
levende og kraftigen beherskes af en sand Idee eller af en reen Følelse;
men Den, der besjæles heraf, nyde da og Raadigheden og Friheden i sin
Kreds! Fremfor Alt nyde Konen den da i Huset! Enhver sundt følende Kone
begjærer Dette, naar hun først er sig sin Dygtighed bevidst. Og det med
Rette. Hun kan kun tjene tilfulde, naar hun kan raade tilfulde. Og er
den sande Følelse der, den ægte, dybe Kjærlighed, der gjør aarvaagen og
gjør seende og gjør samvittighedsfuld, saa er ogsaa for Manden Intet
naturligere og beqvemmere, end Dette. Men man glemme da ikke, at, da
først Egteskabet aabner Qvinden Det, hvad Manden allerede har funden
andetsteds, den eiendommelige Virkesphære, saa være Qvinden forundt,
hvad Manden i sin Sphære saa gjærne ønsker at finde, Ret og Frihed til,
i at udføre, hvad der skal skee, medens en almindelig gjeldende Orden
deri redeligen og troligen iagttages, dog at gjøre det paa sin Viis, og
det heller nu og da mindre beqvemt, inden hun ret faaer dannet sig til
sit Kald, end at den sig udviklende Sands og Eiendommelighed skulde
kues. Men saa glemme da ogsaa Qvinden lige saa lidet som Manden, at, ---
medens det vel med Føie kan kaldes et jordiskt Himmerige, hvad der er
beredet et lykkeligt Par ved en ægte og ungdommelig Elskov, naar de tage
den dem forundte Lykke saaledes, at de selv lade den blive en Himmel for
sig, --- saa er dog dette jordiske Himmerige ingenlunde et Jorderige.
Ligesom Christus for Dem, der med al Magt i ham vilde see en jordisk
Konge, ikke blev den himmelske Konge, han var, saaledes kan alt for let
det Skjønneste og meest Himmelske i Elskovsforholdet blive forvansket
eller forvirret og forstyrret for Dem, der tage Forbindelsen, som om det
var jordisk Herlighed, de nu skulde klæde sig ud med. Det er saa langt
fra, at jordisk Overflod og Pragt skulde fremme en begyndende
egteskabelig Forbindelses Lyksalighed, eller forhøie det sig dannende
Samlivs Tillokkelser, at tvertimod en vis Indskrænkning i Kaar, naar kun
Veien til at erhverve er fri og aaben, bedst stemmer med det nylig
sammenknyttede Pars hele Existents, eller, som man kalder det, bedst
convenerer dem. Netop denne Indskrænkning gjør det begyndende Samliv saa
hyggeligt. Hvorimod en aldeles complet og rigelig huuslig Forfatning og
Besiddelse, hvori de strax sættes ind, findende Alting allerede ordnet
for sig, berøver dem den skjønne Glæde, ved fælleds Stræben og fælleds
Beraadslagning Tid efter anden at fremkalde den hele huuslige Omgivelse,
og uddanne og fuldende den med en til enhver enkelt Deel af samme sig
knyttende Interesse, Omhu og Kjærlighed. Rigdom og Overflod er just ei
altid en Velsignelse. Hvorhos da ogsaa er at erindre, hvad Goethe i
Wilhelm Meisters femte Bog siger, ved at gjøre opmærksom paa det
Forkeerte i de Festiviteter, hvormed egteskabelige Forbindelser begyndes
og ligesom indvies: „der er ingen Fest, der mere skulde høitideligholdes
i Ydmyghed, Stilhed og Haab, end Bryllupsfesten.“ Og dog, hvor ofte
viser det sig ogsaa her, hvorledes smaalig Selviskhed eller smaalige
Hensyn ere rede til at komme frem ved al Leilighed og at beherske
Menneskene. Især være Qvinden her varsom og høre ikke op at være
kjærlighedsfuld og ydmyg af Hjertet, fordi hun er bleven lykkelig. Vel
er det naturligt nok, at det qvindelige Væsen, der, medens hun endnu
saae den kommende Lykke i Haabet, maaskee viiste stort Sammenhold, at
dette samme Væsen, naar Lykkens Fylde er kommen, ikke mere saa vel kan
bevare Ligevægten. Lykkelig Elskov gjør Qvinden fri og myndig, den gjør,
at hun mere end forhen kommer til at føle sig som Den, der har en
Personlighed, har Værd og Betydning, og tilmed har en Villie. Det er vel
ogsaa som oftest godt, om det unge Hjerte, i hvilket nu en Verden af
Følelser kommer frem og arbeider sig frem, saaer Rum og Frihed til ret
af Hjertens Grund at slaae sig løs en Stund. Ingen fordre, at Det, der i
Dag endnu gaaer gjærende op og ned i Sindet, skal have sat sig til
Klarhed og Ro i Morgen. Imidlertid vogte Qvinden sig: Først for den
Egenvillie, de Yttringer af et herriskt Sind, som nu letteligen røre sig
og faae Lyst til at komme op; dernæst vogte hun sig for det Ustadige.
Det bedste Tilhold er her det nærmeste, det, der især er istand til at
bringe Sindet til at samle sig med Kraft, den hjertelige og trolige
Beredelse til et Kald, som virkelig er et Kald. Tillige ogsaa en
aarvaagen Lyst til ret af Hjertet at leve den Elskede til Behag og til
Glæde. Overhovedet betænke enhver Elskende, --- hvad dog saa mange ofte
synes at glemme, --- at Gjenkjærlighed ikke bevares paa anden Maade, end
den erhvervedes, og at man ei let kan blive ved at elskes, naar man
hører op at være elskværdig. Elskovs Løfte betragte Ingen som en Gjeld;
blot som en Gjeld ønske vi dog ikke, at Elskov skal ydes os. Ingen
fortryde paa, tvertimod være det Enhver kjært, at den Kjærlighed,
hvormed vi elskes, har sin Grund i en høiere, og at vi derfor maae vaage
over, at den bestandig kan næres derved, at tilligemed den ogsaa hiin
høiere finder Næring, saa at vi altsaa, for at bevare os den Kjærlighed,
der skjænkedes os, maae bevare og nære hos os Det, der gjør Mennesket
indtagende for Mennesket; og ei blot bevare Det, men og lade Det træde
frem og yttre sig saa, at det ret kan virke velgjørende udad. At elskes
med en blind Kjærlighed er ikke til Fromme for Nogen. Men at elske og
elskes saaledes, at en reen Kjærlighed til det Sande og Gode og til det
Skjønne, og en levende Attraa efter at udtrykke det i sin Person og i
sit Liv, kan befæstes og fremmes ved vor Elskov, det er at elske med
stor Fordeel.

Dog, de smaaligste Yttringer af en egoistisk Lyst ligge ei sjeldent, man
skulde neppe troe, hvor nær, ved Elskovs ellers Egoismen overvældende
Følelser. Alt for ofte yttre Mennesker sig, som om --- naar de neppe ere
komne saavidt, at de have opnaaet en Lyksalighed, der skulde gjøre dem
villige til med Lyst at give Slip paa Meget, hvilket ellers kunde være
tiltrækkende for dem --- som om nu, fremkaldede ved denne Lykke, alle
smaalige Lyster og Begjæringer myldrede frem af alle Hjertets Kroge, og
prøvede paa at fremmes og erholde nogen Næring ved det nye Liv, der nu
har aabnet sig. Og hvad ikke slige Lyster, hvad ikke alskens
Forfængelighed gjør til at besnære og forvirre, det gjøre utallige
Hensyn til Mode, og til Det, de Fleste pleie, eller til Disses og Hines
Dom, saa at det ei sjeldent er vanskeligt nok at komme til at nyde et
ligefremt, naturligt, sundt Liv paa egen Viis. Endog saa vidt gaaer det,
at Mennesker af alskens fremmede Hensyn eller Lyster undergrave Det, der
dog er Basis for deres Lykke, en tryg og paalidelig huuslig Forfatning i
Henseende til Udkomme, medens de just derved, istedetfor at faae nogen
ægte Nydelse derfor, forbærve det Sødeste og Qvægeligste i deres Samliv,
den ømme dybe Omhu for hinanden og den hyggelige Tilbagetrukkenhed og
kjærlighedsfulde Indskrænkning, der først gjør Samlivet til et sandt
idyllisk Liv, Noget, hvorved Egteskabet i sin Begyndelse saa meget
forskjønnes.

\begin{center}\rule{0.28\linewidth}{0.4pt}\end{center}

Men paa de værste, de meest aandsfortærende Maader kan det i sig selv og
i sine Hensyn hildede og besnærede Sind, der overalt virker saa
ruinerende, forstyrre og ødelægge en egteskabelig Forbindelse, naar det
fører til lutter Irriteringer, især til Skinsygens plagende Uro. I fyrige
Menneskers første ungdommelige Elskov komme de qvidefulde, Livet
formørkende og foruroligende Opirringer kun alt for let og ofte op.

Det er ikke strax og let fuldendt, hvad der skal skee, naar tvende
Individualiteter, der hver for sig dog altid i en vis Grad have fastnet
sig og antaget en bestemt Skikkelse, skulle løsnes, for at voxe sammen
til et nyt begge forenende Liv. En vis Kamp mellem Begge indtræder, der
hos heftige Sjæle let kan føre til al den Vexel af Ro og Uro, Fryd og
Qvide, Haab og Angest, hvori ungdommelig Elskov alligevel saa hyppigt
kommer til at yttre sig. Det er i mange Henseender en gjærende Bevægelse,
hvori Elskov sætter Gemytet. Vel er Elskov, om den end er en Lidenskab i
den høiere Betydning af Ordet, hvori enhver Henrevenhed, ogsaa den ægte
geniale, kaldes saa, idet man og taler om en guddommelig Lidenskab, dog
ingenlunde efter sit indre Væsen nødvendigvis en egentlig Lidenskab i
Ordets sædvanlige lavere Betydning; thi ingenlunde ligger det i dens
Natur, at Aandens frie Liv og sunde Fremadstræben skulde lide ved den;
ingenlunde hører til dens Væsen en saa eensidig og heftig Attraa, at den
skulde, ganske betagende Individet, omtaage Sindet, eller dog gjøre det
uklart, og, i Strid med andre Følelser, ophæve den indre trygge Ro og
Ligevægt i Aanden. Men i dens ungdommelige Begyndelse er fyrig Elskov
dog næsten altid lidenskabelig, og kan da forsaavidt føre noget Sygeligt
med sig, skjønt vistnok hos ellers sunde og med ægte Kraft elskende
Væsener intet andet Sygeligt, end hvad enhver af de Kriser kan bevirke,
hvilke høre til et sig udviklende Livs naturlige Gang. Fyrig Elskov
river naturligvis en Tid lang saaledes hen, at Individet, rykket ud af
sin forrige Bane, sværmer omkring i den nye Region, uden fast Holdning i
Henseende til Livets øvrige mange Bevægelser, ja endog uden sikker og
tryg Kraft i Henseende til Elskovens egne Yttringer. Vel løsnes Sjælens
Vinger og udfolde sig i den opgaaende Elskovs varmende Sol, men de
forvikles og hildes ogsaa nu letteligen, idet de kraftigen stræbe at
bevæge sig. Især er herved at lægge Mærke til, hvorledes de samme
Følelser, der med saa megen Kraft kunne rive Mennesket ud af en selvisk
Existentses snævre Kreds, dog i en anden Henseende saa let opægge og
fremskynde Individets Lyst til at haandhæve sin særegne Bestaaen og lade
sin personlige Eiendommelighed nyde Fremme. Da Elskov ikke blot er en
uendelig Hengivenhed, men ogsaa selv drager heftigen til sig og attraaer
af al Magt at møde den samme Hengivenhed, saa er det naturligt, at den
Elskende, søgende og forventende denne Hengivenhed hos den Anden, nu
ogsaa gjør sin Personlighed gjeldende. En vis idelig gjenkommende
gjensidig Modstræben, en vis Resistance, indtræder her naturligvis under
den fremadskridende Sammenvæxt, ja den hører med til Forholdets
Uddannelse, paa det at ikke Harmonien, istedetfor at vorde en ægte
Samstemming, der svarer til Begges Hu og eiendommelige Natur og yder dem
Begge sand Fyldest, skal udspringe af mat Eftergivenhed, og mere vorde en
ydre end en indre, mere en tilsyneladende end en virkelig. De om det
egne Jeg sig dreiende indre Sindsbevægelser og Bestræbelser kunne endog
slet ikke hvile aldeles, da Livet kun danner sig og griber om sig i og
med en sig fremhævende Personlighed. Endog for at nyde sin Lykke
tilfulde og ret at gaae ind i den med Lyst og Kraft, maa det egne Selv
op i Sjælen. Det nye Liv kan ikkun komme Individet fuldelig tilgode,
naar det føler sit inderste Jeg, sin eiendommeligste personlige Natur
derved opklaret, hævet og fremmet, formedelst en sand og dyb Samstemming
imellem sit eget Hjerte og det nye sig aabnende Liv. Foreningen med et
elsket og tilbedet Væsen sætte Ingen i Tvedragt med sig selv og sin
Natur! Derfor være og den Eftergivenhed og Føielighed, der vistnok i
mangt et Punct maa indtræde, ingen svag og hjerteløs, men en hjertelig
og kjærlighedsfuld, der giver efter, fordi Individet af Hjertet vil det,
af Hjertet vil lade sin Kjærlighed have denne Magt hos sig, idet det
ikke føler noget Sandt og Ægte i sit Væsen derved kuet, men netop i
Eftergivenheden føler Qvægelse for sin Hu, og villig giver sig hen til
den Løsning i Sindet, uden hvilken dog det egne Selv, den friere Sjæl,
ikke kan komme til at føle sig hævet og udvidet ved den gjennemtrængende
Kjærligheds Magt. Men ikke strax eller let findes og fastholdes, under
Individualiteternes nu ofte indtrædende Sammenstød, de rette
Holdningspuncter for de saa bevægelige Sindsstemninger og Følelser, de
faste Puncter, hvorom Samlivet kan danne sig til en sand Harmonie og en
fast og ægte Bestaaen. Megen Usikkerhed, Vaklen, Ustadighed og blind
Eensidighed kan nu komme frem, og det viser sig ofte ret tydeligt, at
Eros ikke er af himmelsk Udspring alene. Da den Lyksalighed, Hjertet har
fundet, altid forsaavidt ikkun er en jordisk, som dens Besiddelse ikke
er ophøiet over al Omveltning, saa spiller under Følelsernes Vexel ofte,
snart en vis Angest, snart en vis fristende Lyst til at føle og skue sin
Indflydelse bekræftet, og forvisse sig om sin Lykkes uforandrede
Bevaren, saaledes ind i Samlivet, at derved mangen Irritering opstaaer,
hvilken efterhaanden let fører til den Vexel af Opirringer og
Forsoninger, hvis Betænkelighed vi allerede i det Foregaaende ere komne
til at omtale, en ofte høist aandsfortærende, den dybe Ro og Tryghed
dræbende Vexel af Aandsspændinger og Stramheder og tilbagevendende
Mildninger og Hensmeltninger. Endog til et reent Spil med Følelserne
udarter ofte denne Tilstand, saa at det er, som om de Elskende, om end
halvt uvitterligen, fremkalde de hyppige Elskovstrætter med et Slags
Lyst, eller af en Slags Uvane, blot for under hiin Følelsernes Vexel ret
stærkt at fornemme og nyde deres Elskov, og derved at bringe en vis
Bevægelse ind i et Samliv, der, af Mangel paa ægte Meddelelse og kraftig
Livelighed og en omfattende gjensidig Interesse, undertiden kan blive
døsigt, ja kjedsommeligt, og reent gaae i Staae. Men ofte ligger Grunden
til deslige idelige Opirringer dybere. Ei sjeldent i en oprindelig i
Individualiteterne grundet Disharmonie mellem Gemyter, der i sig selv ei
vel kunne stemme sammen, fordi de ei svare til hinanden, hvor meget end
begge føle sig virkelig forelskede, dragne til hinanden ved visse
Tillokkelsers Magt. Og er der end ingen saadan Disharmonie, kan dog en
vis Strid imellem de selviske Sindsbevægelser og de sympathiske, imellem
Jeget og den Hengivenhed, hvormed det skal tabe sig hen i et andet
Væsen, vedvare længe, inden Individet kommer paa det Rene. Ja, som vi
allerede forhen have havt Anledning til at bemærke, netop Følelsen af
Elskovsbaandet som Baand, kan her blive hildende og betingende og derved
fristende og æggende. Begge blive da, imedens de i høi Grad ere
Incitament for hinanden til de stærkeste Elskovsfølelsers idelige
Oplivelse og Opbrusning, dog nu tillige ogsaa Incitament for hinanden
til lutter Forbittrelse; ofte, ret som om de ogsaa trængte hertil, og
just ogsaa ved denne Incitering droges til hinanden: en Tilstand
imidlertid, hvori der ingen Qvægelse og Husvalelse er. Kun alt for ofte
glemme Elskende i Tide at agte paa, hvorledes de staae med hinanden, at
gaae med Opmærksomhed og Interesse ind i hinandens Naturer, og nu,
sigende sig selv redeligen, hvad de ville, at lade falde, hvad der maa
falde, og drage frem, hvad der skal drages frem, og holde fast paa Det,
der maa fastholdes, paa det at ikke Forvirring og Ufred skal faae
Overhaand, hverken i det egne Indre eller i Samlivet. Dog kan vistnok
hos ungdommelige Elskende ikke enhver Elskovstrætte, enhver Forbittrelse
af den Art, hvilken Elskov alene gjør mulig, udeblive, endog ikke let
enhver ugrundet. Nidkjærheden vilde endog ei synes stor nok, naar den ei
kunde føre til dem. Imidlertid vogte man sig altid for et Følelsernes
Spil, der, medens det tærer paa Aandskraften, kan virke til bestandigen
mere at befæste, ja at forstørre Kløften mellem Begge, der nu istedetfor
at voxe sammen, lade de skarpe Kanter, hvorved de støde an mod hinanden,
fastnes mere og mere.

Vi møde her Jalousien i den første af dens Skikkelser, som en til
Lidenskabelighed stigende, Sindet hildende og besnærende, Nidkjærhed i
det egnes Jegs Haandhævelse i Forhold til, og i Modsætning til, den
Elskede, en Lidenskabelighed, der imidlertid forudsætter Elskoven, og
ikkun i og ved denne antager det Præg, den har, ideligen paa Ny næret og
optændt under Tilbøielighedens stærke Følelser. Thi ikke tale vi her om
det blotte selvraadige og herriske Sindelag, der med Tilsidesættelse af
den Anden, som om denne var et blot Middel, koldt og frossent ikkun vil
have sin Villie frem, og ei spørger om Samstemming. Den under Elskovs
heftigste Følelser ofte indtrædende ivrige Kæmpen for det egne Selv, der
optænder Strid og Hede, vil og kræver netop Samstemming, og det en
Samstemming af Kjærlighed, kræver den Andens udeelte Opmærksomhed og
Beredvillighed og Eftergivenhed overalt og i Alt, som Tegn og Pant paa
Kjærlighedens uafbrudte Bevaren og usvækkede Styrke. Netop med Elskovs
Tvivl og Qvaler, Ængstelser og Forbittrelser sønderriver sig her det
oprørte Sind. De i alle Mennesker saa let fremtrædende, til at yttre sig
saa rede og snare, selviske Sindsbevægelser kaste sig paa Elskoven, og
gribe denne som Det, hvori de sætte sig fast, og rase nu med dobbelt
Styrke, istedetfor at Kjærlighed skulde faae sin fra smaalig Egoisme
befriende Magt og Indflydelse i den Elskendes Hjerte. Den egoistiske
Paaholdenhed paa sig selv og paa den egne Værd, Indflydelse og Betydning,
som en saadan, der nu idelig ligger Individet i Tankerne, fordi den er
dets egen, denne omfatter nu den hele Elskov med. Egoismens Djævel raser
med dobbelt Styrke. Den urene Aand, der en Tidlang havde forladt
Mennesket og var faret omkring i Ørkener og paa tørre Steder og havde
søgt Hvile og ikke fundet den, vender nu, som der staaer skrevet, om
igjen, og finder Huset, den forlod, ryddeligt og udpyntet, og gaaer hen
og tager syv Aander med sig værre end den selv, og gaaer ind og boer der
igjen; og det Sidste bliver for det samme Menneske værre end det Første.
Det Hellige og Skjønne i Elskoven bliver nu selv til Næring for den
Overhaand tagende Syge. Netop den usigelige Ømhed, den grændseløse
Hengivenhed, hvormed et Hjerte af Kjærlighed synker ganske hen til et
andet, netop denne fremstiller sig nu saaledes for det forvirrede Sind,
at just Tanken paa denne, der synes med usvækket Styrke i hvert Øieblik
at skulle vise sig og komme Fordringen imøde, ja ei engang lade det
komme til en Fordring, ansporer til de mange harmfulde Bevægelser, til
de idelige Tvivl og Forbittrelser, og synes ei blot at berettige til
dem, men endog ligesom at helliggjøre dem. Saa befæster Gemytet sig da
endog, ret som om det var med Flid og Iver, i den næsten rasende
Indtændthed og Heftighed, der efterhaanden, som den forbærver Elskovens
egen Skjønhed og gjør Sindet øde, nu ogsaa forbærver og ødelægger det
hele Samliv. Just Individets Lykke fører det ind i den ulykkelige
Forfatning, og jo skjønnere den Himmel, der syntes beredet for den
Lykkelige, viiste sig at være, desto heftigere optændes, ved Tanken paa
den, det frygtelige Helvede. Til Vrede og de fiendtligste Fornemmelser,
til Had og Grusomhed kan nu Elskov føre, naar tilsidst den Tilbedte
synes ikke mere at være værdig til at leve. Ved alle Leiligheder oprøres
nu paa Ny den hele Indeharme i den bestandigen paa sin Betydning og
Indflydelse tilbageskuende, ustadig frem og tilbage kastede Sjæl, der
altid er tilmode, som om den aldrig var tryg og var sin Sag vis, men
idelig spurgte sig og ideligen vilde forfare, om den ogsaa hos den
Elskede endnu gjelder, hvad den skal og vil gjelde. Og i hvilken
Forfatning sættes nu det elskede Individ? Er dette af mildt og føieligt
Sindelag, tilbøieligt til Ro, saa kues det let, eller bliver dog let saa
varsomt og frygtsomt, at Samlivets Friskhed, Livelighed og Munterhed
derved lider, ja at det tilsidst gaaer i Staae, eller dog, om Bevægelse
blive deri, ikkun bevæger sig ængsteligen under idelige Spændinger. Til
saadan Ukjærlighed kommer da tilsidst den heftige skinsyge Kjærlighed,
at den selv virker til at nedslaae og ruinere, hvad der netop maatte
være den magtpaaliggende at holde i Hævd, hvad der just maaskee
fornemmeligen gjør den Elskedes Person indtagende og tiltrækkende, ja
maaskee endog fremskyndte Elskoven selv. Hvilket alt ogsaa indtræder,
men hvortil endnu mangen heftig Strid, Vrede og Bitterhed, Uro og Ufred,
og tilmed mangen ødelæggende Uorden i Gemyterne og i Samlivet kommer,
naar det elskede Individ har for megen Charakteerfasthed og Aandsstyrke
til at lade sig nedkue. Et blot Vedhæng ved en Andens Individualitet,
som om det kun sattes i Bevægelse ved den Andens Liv og Sjæl, ganske
afformet og passet efter denne Anden, et saadant blot Tillæg til det
Liv, hvori et andet Væsen bevæger sig, uden eiendommelige, frie, kraftige
Livsyttringer, kan et energiskt følende Væsen ei udholde at være, og vil
ei være det. Alle Hensyn kan det ei lade gaae til Grunde i det ene
Hensyn til den Elskede, alle Livsbevægelser kan det ei lade betage af
den Elskede, der da vilde være en besnærende og sygelig. En Gud, for
hvem alt Muligt skulde staae tilbage, Alt skulde opoffres, kan intet
Menneske, selv ikke det høiest elskede og agtede, blive for et andet,
saa meget det end i en baade taabelig og formastelig Forkeerthed bærer
sig ad, som om det vilde være det, ja mener at skulle være det.

Vi see saaledes, endnu uden at en Tredies foruroligende Tiltræden virker
med, uden alt Hensyn til en omgivende Omverden, en Skinsyges ustadig
omkringfarende Ild brænde i et Gemyt, der ei lader Elskovens
Hengivenhed faae Kraft til at overvinde og dæmpe de saa let opæggelige
selviske Sindsbevægelser. Og denne forstyrrende Heftighed er ganske nær
beslægtet med den mere egentlige saakaldede Skinsyge. Ved denne, ved den
lidenskabeligen forvirrede og blinde Jalousie over den Elskede, i
Modsætning til og under indtrædende Forhold til, tredie Person, udvides
kun den Sphære, hvori de samme Følelser rase. Kun Anledningerne
forandres. Ogsaa her aabenbarer sig, hvor der ikke i Sandhed er Fare og
Grund til Aarvaagenhed, et i det egne Selv hildet, sit Jeg aldrig
glemmende, men ideligen for dette frygtende og higende Sind. Kun at det
her er en tredie Person, mod hvem tillige Skinsygens Heftighed optændes,
saasnart den Elskede yttrer Forekommenhed og et kjærligt Hjerte imod en
saadan Tredie, eller finder Behag i denne Tredie, der nu bliver
Anledning til Uro og Misnøie og til idelig nye Tvivl om
Gjenkjærlighedens fulde Styrke og Tilforladelighed. Ogsaa her arbeider
Lidenskaben imod sig selv og imod det Bedste og Skjønneste i det hele
Kjærlighedsforhold.

Er ikke Frihed det Kjærlighedens Element, uden hvilket den ei kan komme
op og røre sig som sand, dyb og inderlig Kjærlighed? Ja er ikke Det,
hvad den indtil Jalousie heftige og brændende Elskende vil møde og vide
sig at eie, netop en fri, alene ved Hjertets uimodstaaelige
Tilbøielighed fremkaldet og derfor just ogsaa uanfægtelig Elskov fra den
Elskedes Side? Ingen Elsker, men mindst den i denne Grad ivrige, er
tilfreds med blotte ved en Pligtfølelse fremkaldede Kjærlighedsyttringer;
en umiddelbar alene af Hjertets frie Bevægelse udspringende Hengivenhed
attraaer han, og den Tanke, at Det, der skjænkes ham, alene ydes ham
ifølge en Pligtforestilling, vilde netop opirre og oprøre ham. Men
hvilken Forkeerthed da, at tilintetgjøre eller arbeide imod den
nødvendige Betingelse for Det, han attraaer, den Betingelse, uden
hvilken Elskoven ei kan være den, han vil, eller uden hvilken han
idetmindste aldrig kan vide, om den er en saadan, nemlig en aldeles fri?
Hans Følelsers naturlige Styrke og Conseqvents maatte føre ham til at
ville, at alt Nødende skulde bortfjærnes, selv den allerringeste Tanke
paa Pligt og Rettighed, som om her hverken var en Skyldighed eller en
Fordring. Enhver Tvang, enhver Mindelse om en Fordring og et Maa vilde
jo stride mod det Ønske, at elskes af fri Tilbøielighed alene. Dette
føler ogsaa selv det med sygelig lidenskabelig Heftighed elskende
Menneske ofte, og, tilstaaende at det Intet har at fordre, er det kun
mod sig selv, det bitterhedsfuldt og grusomt vil rase, naar Skinsygen
overfalder det. Men ogsaa da sætter det jo den Elskede i Angst, og
paalægger ogsaa da piinlige og indsnævrende Baand, og forjager den
Frihed, der alene kan forvisse det den rene, trygge og totale Følelse
af, at det i den Elskedes Hjerte virkelig eier et sig med fri Lyst
ganske hengivende Hjerte. Netop Elskovens naturlige Conseqvents maatte
føre til den ægte Frihedssands. At forvandle alt Det, der er en
Nødvendighed, til en Frihedens Sag, og lade Aanden og Aandsfriheden faae
Rum overalt, det gjør Den, der har Sands og Hu og Gave til at leve Livet
paa den ægte Viis; men endog at tage Det, der er en Frihedens Sag og
alene kan have Værd ved sin Frihed, saaledes, at det behandles som en
Gjenstand for bindende Nødvendighed, det er en stor Daarlighed. Og dog,
hvor ofte trække Menneskene Det, hvori den frieste Hjertelighed vil
aabenbare sig, ned i en betingende Fordrings Sphære, og lade (for at
anvende et Sted af Goethe i hans Aus meinem Leben) hinanden bøde nok
for, at de engang have været saa lykkelige at kunne være hinanden til
Glæde.

Og hvor lidet er ogsaa endnu i mangen anden Henseende hiin sygelige
Elskov en ægte Kjærlighed. Ikke at tale om, at den saa ofte allerede er
rodfæstet i det Sandselige, paa en lav og uværdig Maade: selv hvor den
som dybere Følelse, som egentlig Elsken, gaaer ud paa det Sjælelige, er
den, om end i Rod og Udspring en ægte og energisk, dog ei trængt
saaledes igjennem, at den er kommen til at udfolde sig med en sand
Kjærligheds alsidige og dybe Kraft. Hvor lidet tillader den idelige
Kæmpen og Higen, Tvivlen og Forsagen, den idelige Reflexion paa sig selv
og paa sin Besiddelse, som sin, ret som om Haandhævelsen af det egne
Jegs Betydning aldrig kunde glemmes, hvor lidet tillader den de sødeste
Følelser i Elskoven at komme op og at virke med omfattende Styrke! Den
sande sig selv forglemmende Omhues og Ømheds, Mildhedens og
Overbærelsens søde Følelser komme ei til at virke. Den hildede Sjæl
kommer ikke engang til ret at skue eller at nyde, hvad der er givet den.
Hvad der skulde være den Elskede til sand Glæde, hvad der netop med Lyst
skulde fremkaldes, vedligeholdes og næres, den Elskedes aabne Sands for
Livet og den Hjertelighed, hvormed et reent og frit Gemyt, en lykkelig
Sjæl, gjærne kommer ethvert varmt Menneskehjerte imøde og aabner sig for
det, Dette bliver nu for den ulykkelige Lidenskab en aldrig standsende
Kilde til Pine og Uro. Hvilken Kraft har ikke en sand og usvigelig
Kjærlighed, netop den Kjærlighed, Enhver attraaer at finde, til at aabne
og ligesom løsne og frigjøre Sindet, at det nu med frisk Livelighed og
hjertelig Munterhed giver sig hen til ethvert skjønt Indtryk, til enhver
sjælefuld Livsbevægelse og til enhver opvaagnende Følelse af den almene
Kjærlighed til alle Mennesker. Men alt Dette arbeider den utæmmede
Lidenskab i blind Vildelse paa at standse og fortrænge, eller dog at
svække og forrykke hos den Elskede, at dennes omspændte Gemyt tilsidst
af Ængstelighed ei tør yttre sig og udvide sig med Frihed. Paa en øde Ø
gad det under Jalousiens Heftighed livende Menneske leve med den
Elskede, udelukket fra al Verden, paa det at Begge alene kunde aldeles
være hinanden nok, og at for Alting Ingen af dem skulde finde en anden
eller en friere og mere udvidet Himmel, end den, hvilken den
lidenskabelige Sjæl da mener at ville berede sig og sin Elskede: en
Himmel, der i Sandhed vilde blive snæver og fattig nok, om der endda
kunde komme nogen frem, og ikke snarere Skinsygen alligevel vilde rase
under hiin dens første Skikkelse, hvortil ingen Tredies Tiltræden
behøves. Saa formasteligt er det i denne Sindsforvildelse hildede
Menneske, og saa rede til at ruinere den lykkeligste og meest
indtagende Natur hos den Elskede. Og nu lader det tillige i sin egen
Elsken Det gaae til Grunde, der ellers er det Skjønneste i Kjærligheden,
den inderlige rene og derfor rolige og trygge, Sindet udvidende og
frigjørende, usigelige Tillid, en sand og fuldelig Hengivenheds rene
Følge, den umiddelbare i den rene Elsken alene hvilende Forvisning.
Saavidt kunne Jalousiens sælsomme Forvildelser gaae, at noget høist
Uselt kan træde istedetfor det Høieste og Herligste, og at den Elskende,
fortabt i en sig selv sønderrivende Vanvid, hellere gad leve i Tilliden
til Eunucher, end i den dybe hjerteqvægende Forvisning, den uendelige
Tillid til den Elskedes hjertelige Kjærlighed og Troskab. Og hvad er dog
en Kjærlighed værd, hvori ikke disse Tillidens dybe Følelser leve? Hvor
vanhelliget bliver her Nydelsen af det Hellige! I Sandhed skal Elskov
føre til Mistillidens eller de ideligen paa Ny opæggede Tvivls
uophørlige Qvide og Plage og de aldrig hvilende ængstelige Selvhensyn,
og skal derved det Sødeste, Ædleste, Frieste og Bedste i Elskoven
bestandigen kues og aldrig komme til at besjæle Hjertet med fuld Kraft,
saa er det dog mere sundt ikke at elske. Bedre, end at lide og hildes saa
meget under sin Elskovs Lyksalighed, er det dog at rense sig for slig en
Aandsqvide ved at rykke sig dens Rod udaf Hjertet, den heftige
Lidenskab, fra hvilken al denne Sindslidelse kommer, eller dog at lade
den lunknes saavidt, at en taaleligere Existents, just fordi den er
lunknet, kan indtræde. Ja Den, der lider af hiin aandsfortærende Syge,
er maaskee sin Lyksalighed nærmest, naar han kan bringe det saa vidt, at
han forsager paa den heftige Attraa med et religiøst Sindelag, og,
bekjendende sig som svag, er bered til at finde sig i det Smerteligste;
thi da tør maaskee snarest den onde Dæmon vige, saa at Herrens Kraft kan
vise sig stærk i den Svage. Klar Erkjendelse og rolig Betragtning,
hvorved det Unaturlige i Skinsygen, det Ødelæggende og mod sand
Kjærligheds Væsen aldeles Stridende deri, fattes og sees fra alle Sider,
maa paa den ene Side virke; men paa den anden Side den under Aandens
Barndom her paa Jorden alligevel saa ofte fornødne Tvang, Individet
paalægger sig selv. Lider du saa Meget under din qvidefulde Elskov, under
de ideligen tilbagevendende Tvivl og Tanker, der ingen Ro og Hvile unde
dig, saa beslut hos dig selv, at nu vil du, saa heftigt du end elsker,
dog finde dig i, om det skulde saa være, at Gjenkjærligheden hos din
Elskede ikke var stærkere eller tryggere, end den i de bittre Øieblikke
forekommer dig. Forsøg engang, hvad vel en from Resignation kan formaae!
Staaer Alt i Herrens Haand, saa lad ogsaa din Elskovs Lykke staae i den.
Maaskee det da ogsaa vil gaae i Opfyldelse paa dig, at Den, der
retteligen forstaaer at tabe sig selv, han skal vinde sig selv. Ægte
resolut Hengivenhed i Guds Villie er overalt af underfuld Kraft. I den
høieste Kjærlighed ligger Det, der kan give al anden Kjærlighed Hvile og
Fred. Men ogsaa her betænke man tillige, at, hvad der skal komme frem og
voxe i Aandens Liv, ligesaavel, som hvad der voxer i den ydre organiske
Verden, kun efterhaanden naaer sin fornødne Styrke. Ingen vente, at et
Onde, der har faaet Magt i hans Indre, altid viger strax, fordi han
fordømmer det. Det kan komme mange Gange tilbage, inden det aldeles
synker ned. Derfor blive Den, der har Noget at arbeide paa i sit Indre,
ikke utaalmodig eller modløs, men lade ufortrødent Stræben til det Bedre
virke roligt og langsomt!

Den frygtelige Tilstand, i hvilken Jalousie i sin høieste Grad kan føre,
er just ikke saa hyppig. Men Sindsrørelser, som spille ind deri,
Følelser, der ligesom smage af den qvalfulde Pine og af den uregjerlige
Vildhed, Opbrusninger, hvori de sagte Begyndelser til hiin
Sjæleforvirring og hiin Bitterhed røre sig, ere ei saa sjeldne, og ligge
vel de fleste med ungdommelig Fyrighed elskende Sjæle nær nok for
letteligen at kunne opvækkes. Det viser sig her, hvad der saa ofte sees,
at den opnaaede Besiddelse af et med Heftighed attraaet jordiskt Gode
endnu just ikke altid strax er en Himlens Velsignelse. Netop som
jordiskt Gode ligger Elskovs Lykke saa nær ved hine kjærlighedsløse
Sindsbevægelser; thi som jordiskt Gode er den Gjenstand for udelukkende
Besiddelse, derfor ogsaa for Fare og Frygt, ja for brændende Higen og
Jagen, for Tvivl og Ængstelse, naar den aldrig synes tryg og
tilforladelig nok, eller aldrig høi, reen og guddommelig nok. Det viser
sig her, hvor farligt det er at have andre Guder foruden Herren. Just
Det, saa aldeles at sætte Eet og Alt i sin Elskov, eller give den saa
høi en Plads og Betydning, at alt Andet i Livet synes at ville fortæres
af den, og nu at drive det til en Forgudelse, der ei lader Hjertet Rum
til dyb Kjærlighed til Gud og til hjertelig Fromhed, som om den Elskede
var Guddom nok, just en saadan overstadig Tilbedelse af den Elskede
fører let til, at man nu, villende og fordrende det Samme af denne, ogsaa
sætter sig selv som en Gud, og aldrig synes at finde udeelt Kjærlighed
og Opmærksomhed nok. Saa høit maa overhovedet Elskov ikke sættes, saa
hellig maa den ikke agtes, at den skulde føre til at standse det indre
Livs frie Udvikling og lænkebinde den Aand, der er bestemt til paa
Jorden at modnes en Evighed imøde, og hvis Gud Ingen kan være uden den
sande Gud, til hvem hele Livet skal stræbe at komme.

\begin{center}\rule{0.28\linewidth}{0.4pt}\end{center}

%% ============================================================
%%  TREDIE BOG  (printed pp. 130--164; PDF 146--180)
%%  Fraktur in the original. Transcribed and image-verified.
%% ============================================================
\chapter*{Tredie Bog.}
\markboth{Tredie Bog}{Tredie Bog}

I Sandhed Eros maa vel siges at være af blandet Udspring og Natur, og
Alle have ikke Grund til at prise den mægtige Dæmon. Usigelig
velgjørende kan Elskov vorde, naar den i en inderlig Forening under
heldige Kaar sammenknytter tvende Gemyter, der i skjøn Harmonie vide af
Hjertens Grund at nyde den dem forundte Lykke. Men hvor ofte virkede
Elskov til at hilde Aandens Kraft og fremadstræbende Liv, og at holde
det Bedste i Mennesket fangen i Uro og Forvirring. Selv de aandrigste og
elskværdigste Væsener kom undertiden ikke til at leve et skjønt,
aandfuldt og nydelsesrigt Liv med hinanden, alene fordi en ulyksalig
Forelskelse traadte forstyrrende ind i deres Forhold. En reen, Sjælen
fyldende og qvægende Beskuelse af en herlig Individualitets Dybde og
Fylde, en inderlig aandelig Meddelelse og Deeltagelse, sød Fortrolighed,
Tillid og Hengivenhed, og den fulde Forvisning om, i den Andens Hjerte
at besidde en tilforladelig Tilflugt under alle Livets Hændelser: hvilke
Kilder til usigelig Glæde i et indholdsrigt Samliv. Og dog, hvor ofte
kom disse skjønne Goder det Menneske ikke tilgode, de dagligen stode
aabne for, eller bleve ham kun til en ureen, uklar og halv Nydelse, blot
fordi de kaldte Elskovs higende, brændende Attraa med et uroligt Sinds
stormfulde Bevægelser frem! Være det nu, --- thi paa mange Maader kan
Elskov volde Forvirring og Ufred, --- være det, at den opæggede det
Egoistiske i Hjerterne og Skinsygens Heftighed og Brynde; eller at den
drog tvende Individer til hinanden, der kun paa sørgelige og vildførende
Veie kunde komme til Fyldest for en Begjær, som, vakt og næret, men
utilfredsstillet, idelig forstyrrede og foruroligede; eller at den ikkun
opflammedes i det ene af de tvende Hjerter, og nu kastede dette ind i en
haabløs Elskovs aandsfortærende Vexel af lidenskabelige Sindsstemninger,
og i de ustadigen omkringtumlede Følelsers idelig fornyede Oprør. Ogsaa
saavidt kan det komme med Elskovs Vildfarelser, at den, fremkaldet ved
visse Tillokkelsers inciterende Magt, knytter Naturer til hinanden,
hvilke egentligen og i Grunden aldeles ikke harmonere, ja at den
fængsler et elskværdigt Væsen til en aldeles Uværdig. Føie vi hertil
endnu, at Elskov stundom som en fiendtlig Dæmon træder forstyrrende ind
i tidligere kjære, fredelige Forhold, og bringer Mennesket til for sin
Elskedes Skyld at forlade Fader og Moder i en Betydning af Ordet, hvori
det bliver tungt og trykkende for det følende Hjerte! Imod alle
saadanne sørgelige Forviklinger, hvori en ellers i Menneskelivet saa
herlig, betydningsfuld og velgjørende Kjærlighed kan føre Mennesket ind,
er den Elskov næsten ikke at kalde en ulyksalig, hvis ydre Skjæbne er
saa ulykkelig, at de Elskende i jordisk Henseende rives fra hinanden,
medens dens indre Væsen, dens indre Reenhed, Skjønhed, Styrke og Fylde
bliver uanfægtet, og det aandelige Eie derfor kan forblive ufortabt,
reent og klart, saa meget end Hjertet kan lide under det smertelige
jordiske Stød.

\begin{center}\rule{0.28\linewidth}{0.4pt}\end{center}

Hvis nogen Elskov er ulyksalig, saa er den det vist, ved hvilken tvende
Individer, som enten aldeles ikke stemme sammen, eller og over lutter
uordentlige Sindsbevægelser ikke kunne komme til en ægte fredelig
Harmonie, bringes til at indgaae en Forbindelse, hvori de bestandigen
blive Anledning og Incitament for hinanden til tusinde Følelser, der
afmatte og udmarve Aandens Kraft. Det er vanskeligt at sige, hvad der i
slige Forhold meest tærer paa Sjælen og omsnærer eller betynger den,
enten at en virkelig Elskov bliver ved at virke med Heftighed, eller at
de Følelser, der knyttede Forholdet, lunknes og døe hen. Det Første
bliver aandsfortærende ved de mange indtrængende og brændende
Inciteringer, de mange Aandsspændinger og piinlige Ængstelser, den
bestandige Uro og de uophørlige Hensyn, imellem hvilke Sjælen, som paa
et oprørt Hav, veiløs kastes frem og tilbage under en urolig Vexel af
Hede og Kulde. Er derimod Elskoven forsvunden, hvilken Mathed, Tomhed og
Lede indfinder sig let, hvor aandssløvende og dødende bliver Samlivet,
naar Det, der skulde sætte det i Bevægelse, er borte, og andre friske
Livsbevægelser bestandigen standses, hildes eller nedtrykkes ved den
aldrig dulmende Følelse af en Lænke, der hænger fast og hviler tungt paa
Sjælen. Dobbelt angribende for Aand og Hjerte bliver Tilstanden, dersom
den omsnærede Sjæl nu endda kommer til at erkjende, at den er bunden til
en Uværdig, være det, at Tillokkelserne desuagtet --- muligt er det ---
endnu udøve deres forrige Indflydelse, eller at det er ude med
Kjærligheden. Den vedvarende Forbindelse bliver uden Fred eller uden
Fylde, og selv Befrielsen fra det ængstende eller trykkende ydre Forhold
vil, om den forundes, endnu ei befrie den hildede Sjæl, der med halve og
sønderslidte Følelser, maaskee under indre Tvedragt med sig selv, søger
at redde sig ved at rive sig løs.

Mangfoldige ere de forskjellige Tilfælde, som her møde os, naar vi
spørge, hvorvidt slige Baand skulle holdes i Agt eller sønderrives.
Ethvert særegent Tilfælde maa decideres for sig efter dets eiendommelige
Beskaffenhed. Der gives Forhold, i hvilke et reent og bestemt følende,
redeligt Gemyt med Fasthed erkjender, at det ei kan eller tør tænke paa
et Brud; der gives andre, hvori den bestemte Følelse af Adskillelsens
Nødvendighed, endog om Hjertet maaskee endnu ei har revet sig løs, er
lige saa fast og paalidelig. Lykkelig er Den, i hvis Indre under
saadanne ulyksalige Forhold en reen og kraftig, ikke ved fremmed
Tilbøielighed og Lidenskab fremlokket, men af sand Agtelse for den
høiere Kjærlighed og sand Troskab mod den Evige, udspringende Følelse,
der i sin Rolighed og Fasthed bærer Vidne om sin Paalidelighed, taler
med afgjørende Stemme, og føder og nærer en urokkelig Beslutning, den
monne nu byde at forblive eller at flye. Men ikke altid bliver saaledes
snart og let ved en høiere Dom det mellem saa mange Hensyn frem og
tilbage trukne Sind fæstet til et bestemt Punct og det foruroligede
Hjerte hævet op over de bange Tvivl og den ængstelige Vaklen. Da er det
især vigtigt ei at haste, men at tøve og lade de første Bevægelser
dulme, og ei at ile fra Følelse til Villie, fra Forsæt til Handling, paa
det at langsomt den nærede Følelse kan forsøge sig under Livets Fremgang
og mangfoldige Sammenstød, om den vil holde ud under dem. Ogsaa
Phantasien maa have Rum til forud at udmale sig den Fremtid, hvortil
Menneskets Hu begynder at staae, om denne Fremtid vel vil blive ved at
vise sig i den Skikkelse, den først syntes at ville antage. Ofte lærte
Mennesker først ved at adskilles, hvor Meget de vare for hinanden: saa
gjøre Sindet da et Forsøg med sig selv, ved at sætte sig levende ind i
den Forestilling, at Adskillelsen allerede var skeet.

Alle de Overveielser, hvilke Spørgsmaalet om Skilsmissers større eller
ringere Tilladelighed fremkalder, naar de betragtes fra Lovgivningens
Standpunct i en Stat, især i en bestemt Tidsperiode, ville vi her ikke
gaae ind i. Ikkun Det, hvad der under slige Forviklinger bevæger sig i
ædle og rene Menneskehjerter, hvad der for et redeligt og kraftigt, af
en høiere Kjærlighed besjælet Sind frembyder faste Holdningspuncter for
de maaskee vaklende Beslutninger, ikkun Dette ville vi gjøre til
Gjenstand for saadanne almindelige Betragtninger, hvilke i givne
Tilfælde snart kunne give Udslaget til den ene, snart til den anden
Side.

Foreløbigen bemærke vi i Almindelighed: Uagtet Egteskabet i sig selv
altid bør at betragtes som en efter sin indre Natur uopløselig
Forbindelse, fordi den altid maa være knyttet med den Tanke, at den
skal gjelde for hele Livet, saa tør man dog i Sandhed ikke gaae saavidt,
at man skulde modsætte sig enhver Skilsmisse, hvilken ikke grove Brud
paa Egteskabet have gjort nødvendig, og tør overhovedet ikke indskrænke
Tilstedelsen af en i Mindelighed besluttet Afskillelse inden de
snævreste Grændser. Der gives mangen Skilsmisse, over hvilken enhver
Sindig, der staaer nær nok for at see, hvad der fremkaldte den, kan
glæde sig. Er det end vel i og for sig ikke glædeligt, at Skilsmisser
overhovedet existere, og er deres Hyppighed end et sørgeligt Tidens
Tegn, saa er dog ikke dette Tegn strax at ansee for det Sørgelige, men
meget mere Det, hvorpaa det er et Tegn. Er Skilsmissen ikke glædelig,
saa var ofte Egteskabet det endnu mindre. Det forholder sig hermed som
med Retsstridigheder eller Processer, hvis tiltagende Mængde ei er et
Beviis paa en Stats tiltagende Sundhed, men hvilke dog hellere maae
tiltage, end at upaaankede Retskrænkelsers Antal skulde stige, og
Uretfærdigheden gaae sin frie Gang. Directe fremmes vist ikke
egteskabelige Forbindelsers Hellighed derved, at der holdes Hævd over
saadanne Egteskaber, hvilke ligesaa lidet leves til Behag for Gud og
Mennesker, som de i nogen Henseende komme enten de forbundne Personer
tilgode eller deres Børn; endog for Børnenes Skyld kan ei sjeldent
Forældrenes Afskillelse være det Raadeligste. En rolig, fredelig, tryg
Existents er Basis for al ægte Fremvæxt i det Gode; under lutter
Aandsspændinger og urolige Sindsbevægelser eller i en qvalm og trykkende
Omgivelse kan Intet trives i Menneskets Aand. Lader man Det, hvad den
legemlige Sundheds Bevarelse kræver, være af saa stor Vægt, saa betænke
man ogsaa den sjælelige. Der er i Menneskeverdenen ikke saa stor
Overflod paa Aandelighed, at man jo maa være aarvaagen mod alt Det,
hvorved der stormes ind paa Sjælesundheden, og man ansee det ei for en
Bisag ved Betragtningen af et Forholds Beskaffenhed, om en opvakt og
livelig Aand deri tæres bort og forkommer, eller næres og trives.
Bringer ofte Overraskelse og Svaghed, Letsindighed og et vist
ungdommeligt Overmod, eller Temperamentets Heftighed Mennesker, førend
de endnu ere komne til Klarhed over dem selv og over Livet, til at
knytte Forbindelser, der aldrig burde have været indgaaede, saa bør,
naar Sligt nu for sildigt blives vaer, en kjærlig Omhu, --- og en saadan
skal dog finde Sted overalt, --- ei lade dem lide derfor et heelt
Menneskeliv igjennem, og det lide paa en Maade, der kan nedkue deres
hele Fremadstræben til høiere Maal. Vel maa Mennesket have for megen
Ærefrygt for det Hellige til, let og snart at lade sig rive hen til at
ønske hellige Baand løsnede, fordi de synes at trykke, eller fordi noget
Tækkeligere synes at tilbyde sig, og Ingen ansee sig for kuet eller
omsnæret, fordi et Forhold betinger og begrændser og fører Pligter med
sig, der ei altid ere lette og behagelige. Men i et Liv, hvori selv den
Bedste daglig maa bede: Lad os ikke falde ud i Fristelse! gives der dog
ei saa meget sjeldent Tilfælde, hvori Mennesker saaledes dagligen friste
hinanden til utallige fiendske, onde, bittre, mørke Tanker og Følelser,
at det er bedre for dem, at de skilles ad. Pauli Ord i 1 Kor.\ 7, 5: paa
det at Satanas ikke skal friste Eder formedelst Eders Uregjerlighed (διὰ
τὴν ἀκρασίαν ὑμῶν), have Anvendelse i mangt et Tilfælde udenfor det, han
omtaler.

Imidlertid er det saaledes vistnok blot for vore Synders Skyld, at Noget
maa tilstedes og kaldes det Bedste, hvad der derfor endnu just ei kan
kaldes godt. Og saa maa da tillige erindres, at netop ogsaa for disse
vore Synders Skyld mangt et Baand maa binde stærkere, end det med Hensyn
til det enkelte Tilfælde alene vilde behøves, om ei den almindelige
Letsind og Uregjerlighed var. Mangen Disharmonie, der nu fører til et
Brud, hvorved der ingen Velsignelse er, vilde maaskee ei være kommen
saavidt, dersom den ei var bleven holdt fast og forfremmet, fordi Tanken
paa Afskillelsens Lethed gjorde mindre agtsom og omhyggelig. Naar
Egteskabsbaandet først viser sig som let opløseligt, viser det sig ogsaa
snart som mindre helligt, og efterhaanden kommer det da saavidt, at
allerede ved Egteskabets Indgaaelse en større og større Letsindighed
indtræder. Derfor maa da for at haandhæve Grundsætningerne og deres
Virkning i det Hele, for at haandhæve Egteskabernes Hellighed i
Almindelighed, stundom Det nægtes i det enkelte Tilfælde, hvad der med
Hensyn til de tvende Individer, betragtede for sig alene, maaskee vilde
være det Raadeligste. Dog glemme man ikke, at til det Hele ogsaa de
tvende Individer høre! Nødvendigheden af at vedligeholde Ærbødighed for
hellige Baand kan fordre Meget indirecte; men dog betænke man altid, at
Egteskabet, ligesom Sabbaten, er til for Menneskets Skyld, ei Mennesket
for Egteskabets, og at Individualiteten har en stor Ret. Bedre end
Forbud, der afskjære al fri Selvbestemmelse og egen Beslutning, turde
vel ogsaa saadanne Foranstaltninger virke, der, ved at sætte en hellig
Ærefrygts og Religiøsitetens Følelser i Bevægelse, kunne bringe
Gemyterne til med Frihed at træde tilbage. Det vilde sikkert i denne
Henseende være af Virkning, om Egteskaber ligesaa vel opløstes ved en
kirkelig Act, som de indgaaes ved den. Som da overhovedet Skilsmisse vel
er en Handling af den Natur, at den bør skee med en Høitidelighed, der
kan minde Mennesket om, at det ved Alt bør have Gud for Øie, og gjøre,
hvad det gjør, med et fromt Sindelag. Mennesker, der engang saa
høitideligen ere blevne knyttede til hinanden, og have været saa Meget
for hinanden, bør dog ikke gaae fra hinanden for bestandigt, uden med et
ved religiøse Forestillinger til Eftertanke og Alvor og til Fredelighed
stemmet Sind.

Saameget foreløbigen i Almindelighed. For Egtefæller, hos hvilke Tanken
paa Skilsmisse kommer op, og for Den, der staaer nær nok, for at være
kaldet til at virke paa deres Sind, er endnu Meget at overveie. Først:
Hvad det vil sige, fra nu af at skulle tænke sig som fremmed for Den, ja
maaskee at skye at see Den, eller dog uden fordums fortrolige Tiltale at
maatte gaae forbi Den, hvem man har været saa inderlig nær forenet med,
og hvem maaskee i tidligere Tider Hjertet mangen Gang har banket imøde
af Lyst og Længsel. Især maa det ikke være et kun nogenlunde dybt og
fiint følende qvindeligt Væsen let, uden at komme i indre Uro og
Forvirring og føle sig ængstelig tilmode, at forene Tanken paa den
Fjærnhed, der nu skulde indtræde, med en levende Mindelse om Alt, hvad
det engang har været for den Mand, det nu vil rive sig løs fra,
hvorledes det har været forbundet med ham, hvor nøie og hvor nær. Det
kan vist ikke let, naar det udmaler Forholdets hele Beskaffenhed for
sig, udholde den Tanke,
at en saadan Hengivenhed, et Liv, hvori det saaledes har tilhørt en
Anden, nu skulde agtes for en frygtelig Drøm, det maatte skye at tænke
paa eller høre omtale. Og nu til de ruinerede eller dog forbittrede
Erindringer udaf en maaskee meget skjøn Fortid endnu de mange med det ene
Forhold tillige sønderrevne eller forvirrede øvrige Forhold; fremfor alt
Det, at, om Børn ere der, ogsaa disse nu ret som andre Ting skulle deles
og rives bort fra den Ene eller den Anden, og, under en unaturlig
Forstyrring i deres Indre, enten blive fremmede for Hiin eller Denne,
eller og befinde sig i en uklar og spændt Stilling imellem dem Begge. I
Sandhed der maa dog Meget til, for at kunne ville og beslutte Noget, der
er saa unaturligt og stødende for Følelsen.

Og dog gives der Tilfælde, hvori et Menneske, for at bevare Troskaben
mod sig selv og mod noget Høiere, maa ville dette Unaturlige. Men saa
ville Mennesket det da ogsaa kun, hvor en høiere Troskab fordrer det. Ei
om en Ret være her Tale, men om en Pligt, ei om det tør, men om det bør
rive sig løs, være Spørgsmaalet. Kun naar en hellig Pligtfølelse bevæger
Sindet og driver til Beslutningen, lade Mennesket denne faae Rum og Magt
hos sig. Til de i moralsk Henseende vilkaarlige og blot tilladte
Handlinger, ved hvilke, saasnart kun den udvortes Leilighed tilbyder sig,
Lyst og Tilbøielighed tør raade over Valget, kan dog en Skilsmisse
ingensinde henhøre. Og heri have vi da Noget, som kan orientere og lede i
mange Henseender. Ikkun naar Individet føler sig kaldet og tilskyndet ved
en af Kjærligheden til det Gode og af en trofast Vedholdenhed ved dette
udspringende Følelse af Adskillelsens Nødvendighed, gribe det til at
beslutte den. En med Bestemthed og Klarhed sig forkyndende Følelse af en
hellig Forpligtelse kan her alene frembyde et fast Holdningspunct for de
maaskee ellers ustadige Tanker og forvildende Ønsker. Derpaa prøve her
Hjertet den opstaaende Attraa, om det vel har Mod til at kalde Det, den
attraaer, en Pligt for sig. Men det være heri redeligt mod sig selv, og
vogte sig for at skuffe sig. Hvad der ei er ægte og godt, kan ei føre til
sand Fred. Det sande Godes redelige Overholdelse er Grundvolden for en
lyksalig Tilværelse, først i indre, siden i ydre Henseender. Hellige
Love, naar vi overtræde dem, vende sig som fiendtlige mod os selv;
indvortes lade de os ingen Ro, og udvortes tør vi ikke længere fremdeles
i Livet haabe Beskyttelse af dem.

Men naar der nu alene er Tale om et Bør og om Det, hvad en hellig
Nødvendighed fordrer, saa betænke man for det Første, at den Troskab, der
hører til Kjærlighedens Væsen, selv er noget i høi Grad Helligt.

Om end Det, der fra først af knytter tvende menneskelige Væsener til
hinanden, er Noget, som er høiere end begge, og de elske hinanden, fordi
de elske det Gode og Skjønne, saa knytter Elskov sig dog ganske til
Personen, som Person, og fordi denne Person var os og er os, hvad den var
og er os, føle vi os bundne til den. Derfor troe da Følelsen sig ikke
strax berettiget til at træde tilbage og maaskee at vende sig til andre
Kanter, fordi den Elskede begynder at synes mindre indtagende eller
mindre elskværdig! Om ogsaa den personlige Tiltrækning og det personlige
Velbehag tager af i Inderlighed og Sødme, saa aftage dog ingenlunde
Trofastheden og den egentlige Kjærlighed og Omhu! Netop da føle Hjertet,
og nære den Følelse, at det jo dog ikke var for Intet, at det gav sig til
dette bestemte Individ, at det jo ikke blot vilde modtage og glædes, men
ogsaa vilde meddele og vise sig virksom i Omsorg, og i Alt vilde dele
Skjæbne med den Elskede. Mindst af Alt tør fremmed Tilbøielighed, Lyst
til Dette eller Hiint, der synes mere indtagende, eller opstaaende ny
Lidenskab, faae Rum og Indflydelse. At lade Deslige faae Magt over sig,
det er at kaste sig ind i en Hvirvel, der tilsidst kan kaste os mod de
frygteligste Klipper eller kaste os ud paa de ødeste Strande; det er at
fordre alt Helligt op imod sig. Endog Det, der meest af Alt maatte kunne
synes at være opfordrende, Overbeviisningen om den Medforbundnes moralske
Upaalidelighed eller moralske Uværd, kan endnu ei strax være et kraftigt
og sundt følende Hjerte Nok, for at ville skilles fra Den, det paa saa
fast og betydningsfuld en Maade, som ved et Egteskab, er knyttet til. Saa
længe det endnu har den Magt over det andet Hjerte, at det kan virke
heldigen paa dette og kan holde sig selv reent, saa vil det netop føle
sig opfordret til at blive, ja vilde forekomme sig som Den, der vilde
undflye et helligt Kald, om det tænkte paa at rive sig løs. Her gjelde
Pauli Ord i 1 Kor. 7, 16: „hvor veed Du, Hustru, om Du ikke kan frelse
Manden, eller hvor veed Du, Mand, om Du ikke kan frelse Hustruen?“ Kun
naar Mennesket føler sig selv ude af Stand til at virke, og nu tillige
maaskee i Fare for at komme i fristende og tunge Kampe, hvoraf det
maaskee ei kan redde sig med et reent og heelt Sind, eller naar det seer
sig nedværdiget eller mishandlet, eller uden Gavn for Nogen bundet og
begrændset i sin Stræben til det Bedre, kan det Tilfælde indtræde, da det
føler sig kaldet til at fordre Befrielse. Hellige Love fordre ogsaa
derved deres Hellighed haandhævet, at Mennesket ikke skal, medens det
selv taaler det Uværdige, taale det Helligste at trækkes ned og
vanhelliges. Til saadanne Vanhelligelser kan da ogsaa undertiden den
heftige Jalousie henregnes, naar den i sin Vildhed ødelægger eller
fordærver det Bedste i Kjærligheden. Der kunne gives Tilfælde, hvori det,
naar ellers Intet vil hjelpe mod den fortærende Sygdom, kan blive en
hellig Pligt baade for Den, der lider under sin egen Jalousies ustyrlige
Heftighed, og for Den, der lider under den Elskedes, og det for Enhver af
dem Pligt baade mod sig selv og mod den Anden, at unde sig selv og
hinanden Hvile, ved at opløse reent det Forhold, der volder dem saa megen
Qvide og Plage, og er dem en saa stor Anledning eller Fristelse til, at
saare Meget opægges i deres Indre, hvoraf det aldrig er godt for
Mennesket at opfyldes. Hvad Omsorgen for de fælleds Børn angaaer, om der
ere saadanne, saa kan denne i mange Tilfælde lige saa bestemt fordre
Egteskabets Opløsning, som den i andre Tilfælde kan give den ene eller
den anden af Forældrene Kraft og fast Villie til at finde sig i Alt for
at opretholde det.

En i Mindelighed fuldbragt Skilsmisse er just ei altid, om end Begges
udvortes Samtykke haves, skeet med Begges sande indre Samstemming, især
ikke altid fra først af. Hvor derimod Begge ret af Hjertet ønske og
begjære Skilsmissen, synes der, med Hensyn til de to Individer i og for
sig, ikke at være nogen Betænkelighed ved den, og der synes intet
Spørgsmaal at kunne være om Troskabspligt, hvor Begge med lige Lyst
fritage hinanden for den. Imidlertid: Troskaben har her i mange
Henseender sin Grund i en høiere, og hænger paa saa mange Maader sammen
med denne høiere, at man ikke maa troe Alting at være klart, fordi Begge
med lige Lyst og Villighed komme hinanden i Møde i deres Ønske, og fordi
der heller ikke er nogen Tredie, som derved kunde lide. Mennesket være og
sig selv tro og sine Beslutninger, og holde troligen fast ved sin
Ungdoms bedste Erindringer og Følelser, som ved en trofast Ven. Hvad der
i dets Livs bedste Øieblikke opfyldte det, det have det Respect for i de
mange matte og mørknede. Et Brud med Den, der engang var Hjertet kjær,
endog det lempeligste og mildeste Brud, gjør en forstyrrende Rivt ind i
hele Livet, og det kommer Mennesket i mange Henseender ikke tilgode at
bryde sit Livs Følge af, eller at kaste et heelt forløbet Liv fra sig. En
saadan Forandring er altid betænkelig og farlig, og kan lettelig for
Fremtiden gjøre Livet, der nu er kommet ud af sin forrige Retning, uvist,
ustadigt og vaklende, dersom ikke faste og kraftige Følelser af høiere
Art bestemte til Beslutningen og derfor træde ind istedetfor de fortabte
og give Livet ny Continuitet og Holdning. Af blot Lyst og Tykke lade man
ogsaa her ikke det Hellige blive afhængigt; ogsaa her kan kun en vis
Nødvendighed, en Tilskyndelse af moralsk Natur, medføre Beføielsen. Det
er sikkrere at beherske et Sind, der uregjerligt vil gribe efter Dette
eller Hiint, end at lade det Tøilen. Ere vi alle her paa Jorden som paa
en Vei med hinanden til en bedre Verden, saa maae vi holde ud med
hinanden, naar vi først have besluttet troligen at følges ad, og ei strax
gaae fra hinanden til Høire eller Venstre, blot fordi det tykkes
beqvemmere.

Iøvrigt er det ikke ligegyldigt, af hvad Beskaffenhed det Forhold er og
var, der brydes, hverken hvad Art af Forbindelse det er, ei heller
hvorledes den knyttedes og begyndte. Den blot foreløbige Forbindelse,
Forlovelsen, binder naturligvis ikke saa stærkt, som Egteskabet, og et
Forhold, der vitterligen blev indgaaet uden egentlig Kjærlighed, tør
snarere løses, end et, hvilket virkelig Forelskelse knyttede.

Naar det endnu ei er kommet til Egteskab, falde naturligvis en Deel af de
Betragtninger bort, der gjøre Egteskabets Opløsning saa betænkelig. Endnu
have Begge ikke givet sig ganske hen til hinanden og forladt Livets
forrige Gang og Art i den Grad, som den sluttelige Forbindelse medfører
det; endnu er det sidste og afgjørende hellige Løvte ikke givet.
Ophævelsens større Lethed er allerede en ligefrem Følge deraf, at den
foreløbige Forbindelse kan ansees som en Art af Prøvelse, om end de
Forbundne ei betragtede den saaledes, eller vilde have betragtet den
saaledes af Andre. Et saa helligt Forhold, som et Egteskab, skal dog
heller ikke indgaaes med et halvt Hjerte, og det vilde det dog, om Lysten
til at blive fri derfor allerede havde reist sig. Den, der føler denne
Lyst, kan dog, saalænge den ei er slaaet ned igjen, ikke give det hellige
Løvte, som det skal gives, af fuldt Hjerte uden indre Modstræben; saa at
Trofasthed og Redelighed selv synes at fordre, at han ikke fordølger sit
Hjertes sande Følelser og dermed faaer Udseende af at love, hvad han ei
er i Stand til at yde. Men med alt det vil et ordentligt og kraftigt
følende og villende Gemyt, der ikke lader enhver opkommende
Sindsbevægelse eller Sindsstemning strax faae Rum til at raade eller og
kun til at tale, være langsomt til at træde tilbage eller yttre
Tilbøieligheden dertil, hvor det har bundet sig ved et Løvte, som, var
det end kun et foreløbigt, dog var et helligt, og i Hjertet allerede et
afgjørende. Det maa ogsaa her føle, at det maa blive det eengang udkaarne
Væsen og blive sit Livs gode Øieblikke tro, saalænge ingen Følelse af
høiere Natur fordrer Ophævelsen. Den Bevidsthed at have grebet
forstyrrende ind i et andet Menneskes indre Sjæletilstand og ydre
Forfatning er for trykkende, til at Livet kan leves frit og let, naar man
bærer den med sig. Imidlertid kunne under Bevæggrundenes Modstrid her
vistnok de, der fordre Forholdets Bevarelse, ikke virke saa stærkt, som
ved Egteskabet. Og er det først klart, at Begge ikke harmonere og ikke
vel kunne haabe, at de ville komme til at leve et lyksaligt Liv med
hinanden, saa er det det Bedste, at Begge forene sig om at adskilles, paa
det at de ikke skulle synes at ville friste Herren, deres Gud, ved dog at
kaste sig ind i Forholdet og at give deres Livs Fred til Priis. Ogsaa den
Forskjel bliver altid imellem Egteskabet og Forlovelsen, at, saalænge
Sjælen opfyldes af Tvivl, eller dog ingen tilforladelig Bestemmelsesgrund
er kommen til moden Kraft i Sjælen, saalænge bør der ei handles, saalænge
bør altsaa et Egteskab hverken hæves, naar det allerede er der, eller
knyttes, om det ei er indgaaet; Forlovelsen bør altsaa saalænge ikke gaae
over til egteskabelig Forbindelse. Beslutter imidlertid den Ene for den
Andens Skyld, uagtet den indre Modstræben, at holde fast, for ei at
betynge sig med Følelsen af, at have forstyrret den Andens Fred og Vel,
saa komme denne Beslutning først til Modenhed og Kraft, og hvad der er
besluttet, skee nu med Hengivenhed i Gud og med et heelt Hjerte.

En Forbindelse, der aldrig var en virkelig indre Sjæleforbindelse, hvori
aldrig et Hjerte af ganske Hjerte hengav sig til et andet, tør
naturligvis allersnarest blandt alle Forbindelser ophæves, saasnart det
tillige var vitterligt for Begge, at gjensidig Elskovs Følelser ikke
droge dem til hinanden. Da her kun yderligen og i ydre Henseender lovedes
Tro, saa brydes her altsaa kun ved Ophævelsen en saadan ydre Tro, det er
ingen virkelig ægte. Naar et Væsen ved Overraskelse eller en Art af
Overvældelse i tidlig ungdommelig Alder blev et Bytte for et andet
Individ, der vel vidste, at det bemægtigede sig hiint andet, enten uden
Gjenelskov, eller dog, om en vis Gjenelskov blev ægget frem, bemægtigede
sig det i en Alder, hvori dette slet ikke endnu kunde være kommen til det
Modenheds Punct, det er hellig Pligt at oppebie, inden et helligt Løvte
begjæres: naar et Væsen saaledes er blevet en Andens Bytte, er det
naturligt, at det, allerede blot derfor alene, senere hen, naar det
kommer til Klarhed over sig selv og Livet, kan fordre Forholdets
Ophævelse, tilskyndet til denne Fordring ved en retfærdig Indignation,
paa det at dog ikke den listige, eller voldsomme, eller dog egenmægtige
og altid uretlige og formastelige Bemægtigelse skal blive ved at binde
den uværdigen Behandlede, og at gjelde for retlig og god, fordi den
lykkedes. En Fremfærd, som i sig var uretlig og formastelig, bør vel
binde Den, der brugte den, men dog ei binde Den, der var Gjenstand for
den uværdige Behandling. Imidlertid vil dog et ædelt følende Hjerte endog
i et saadant Tilfælde være langsomt til Beslutning og vaersomt. Thi
saameget end Det, der skal skee, er Frihedens og Overlæggets Sag, saa
faaer dog Det, hvad der er skeet, forsaavidt det er skeet, for Den, der
lider derunder, Udseende af en høiere Tilskikkelse. Og det Hjerte, der
seer sig fænslet, kan ei tænke sig som det, der uden en Styrelse, det maa
respectere, er kommet ind i sin nuværende Forfatning. Endog dets tidligere
Blindhed eller Mangel paa Kraft kan det nu regne herhen; og derfor maa
det ogsaa her vogte sig for, at det ikke beslutter Noget, hvorved det
rokker den Grund, hvorpaa dets Existents hviler, kaster sig ind i en
ustadig og uskikker%[FLAG: "uskikker" as printed p.149 — perh. "uskikkelig"/"uskikket"; verify]
Livets Gang, og sætter sig i Tvedragt med sig selv eller sin Natur. Det
oppebie derfor det Tidspunct, da en reen og hellig Beslutning kommer til
Modenhed og Fasthed, og da det føler sig vis paa sig selv og paa sin
Beslutnings Reenhed og Godhed.

Overhovedet blive, under hvilkensomhelst af de mange ulyksalige
Forviklinger, hvori et Gemyt kan finde sig stedet, ingen Beslutning
fattet i Heftighed, eller medens Sindet endnu uroligen arbeider med sig
selv! Hvor vigtige Beslutninger skulle fattes i Anliggender, der dreie
sig om det egne Vel, især hvor det gjelder det egne Jegs Haandhævelse i
Modsætning til andre Mennesker, komme letteligen, selv om det er klart,
hvad der bør skee, de selviske eller egoistiske Sindsbevægelser under
alle Slags Skikkelser op med, og gjøre Sindet uklart, og besmitte eller
forurene de Motiver, der alene skulde besjæle Mennesket. Et redeligt
Gemyt kan meget vel fornemme, om Noget, der ikke er godt og reent,
spiller med ind i Følelserne, og Enhver være rede til at tilstaae sig, at
han forvirres og anfægtes af Det, der giver os alle Nok at gjøre med i
vort Indre, og hvad Ingen kan sige sig sikker for. Det er ikke for Intet,
at vi dagligen skulle bede: Frels os fra det Onde. Dersom vi vare
ophøiede over alle dets Angreb, hvorfor skulde vi vel bede saaledes? Saa
bringe da Sindet først sig selv paa det Rene og komme til indre Klarhed
og Ro, inden det begynder udad til at virke og handle, for at faae
fuldbragt, hvad det føler at være ret og nødvendigt! Det er ei
ligegyldigt, med hvad Sind, og tilmed ei heller i hvad Tidspunct, Det
skeer, hvad der maaskee i al Fald skal skee.

\begin{center}\rule{0.28\linewidth}{0.4pt}\end{center}

Forladende en Gjenstand, ved hvis Betragtning man ikkun forsaavidt kan
dvæle med nogen Lyst, som man kan haabe, ved at oplyse den, at forebygge
eller løse Forvirringer, ville vi strax gaae til en anden Gjenstand, om
hvilken det Samme gjelder: det ulykkelige Forhold eller Sammenstød af
Forhold, under hvilke Hjertet ikke kan følge sin rene, frie Tilbøielighed
og efter Lyst raade for sin Fremtid i Henseende til et Individ, der
attraaer en Forbindelse med det, uden at støde an mod de Følelser, det
længe roligt nærede mod andre Personer, især imod Forældre, og at gjøre
et forstyrrende Brud paa Forholdet til disse. Ogsaa under saadanne
Sammenstød i Livet skeer det ofte, at Sindet kastes frem og tilbage
imellem modsatte Forestillinger, og hverken paa den ene eller paa den
anden Side seer sig i Stand til strax at komme til en reen, klar, Ro og
Fred bringende Forfatning.

To forskjellige Tilstande og Forhold er det her af Vigtighed at adskille.
Naar det af Elskov bevægede Hjerte med Heftighed attraaer at knytte en
Forbindelse tvertimod Deres Ønske og Villie, af hvilke det fra Ungdommen
af har været vant til med hellig Ærefrygt at lade sig lede, er Forholdet
et ganske andet, end det er, naar der er Tale om, fordi disse ønske det
og anspore dertil, at skulle indgaae en Forbindelse, uden at elske. I
sidste Tilfælde er Modstanden langt sikkrere, end i første.

Allerede den egne Personlighed maa være Mennesket for hellig til, at det
skulde lade sin aandelige Velfærd, maaskee i de høieste og vigtigste
Henseender, sættes i Vove for lavere Beregningers og Hensyns Skyld. Men
endnu mere maa en kraftig Følelse for den egteskabelige Forbindelses
Hellighed give Kraft til Modstand, hvor det gjelder, om vi skulle taale,
at Menneskehedens skjønneste Eie skal trædes ned eller forvanskes i vor
Person. Ethvert Menneske kan og maa i saadanne Tilfælde føle sig som
Menneskehedens Repræsentant.

Med Undertrykkelse af Hjertets bedste Følelser at indgaae en Forbindelse,
der ifølge sin Natur bør knyttes af en fri og total Hjertets
Tilbøielighed, er i sig selv altid en Profanation af noget Helligt, en
Nedværdigelse af Det, der overalt bør holdes i Agt og Respect. Den
egteskabelige Forbindelse maa agtes for høit til, at den skulde blive til
Gjenstand for nogen Art af Tvang, eller af saadan fremmed Paavirkning,
der binder eller hilder. Dertil kommer, at, hvor de Personer, for hvis
Skyld et Menneskehjerte skulde offre sig, blot ledes af saadanne Hensyn,
der ere af egennyttigt eller forfængeligt Udspring, eller ledes af alt
for store Ængstelser for den Fremtid, der dog altid staaer i Herrens
Haand og med Hensyn til hvilken dog Ingen selv kan være sig sit Forsyn,
er det endog for disse Personers egen Skyld betænkeligt at give efter. At
føie sig efter Andre for at understøtte dem i lave jordiske Begjæringer
og Planer, ja i Det, der maaskee engang kan falde tungt og trykkende paa
deres Samvittighed, er i Sandhed dog ikke af Hjertens Grund at ville dem
vel. Hvad der ikke er af det Gode, er ikke heller til sand Fromme for
Nogen. Denne Betragtning kan ogsaa finde Plads der, hvor det just ikke
ere lave Hensyn eller smaalige Betænkeligheder, men mere naturlige og
menneskelige Følelser, der ere i Bevægelse. Ulyksalige Børn ere dog, som
vi allerede eengang have havt Leilighed til at gjøre opmærksom paa, saa
stor en Ulykke for kjærlige eller blot menneskeligt følende Forældre, at
netop Børnene selv fra denne Betragtning kunne hente stærke Motiver til
fast Modstand, for i Tide at sikkre elskede Forældre for hiin
hjerteskjærende Ulykke. Imidlertid: I et Liv, hvori den Sjæl, der meest
syntes kaldet til at røre sig frit i en høiere Sphære, og meest længes
derefter, ei sjeldent trykkes af aandeligen trange Kaar, i et Liv, hvori
Menneskehjertets skjønneste og bedste Attraa og Stræben saa ofte ikke
finder Rum og Frihed til skjøn og harmonisk Udfoldning, ja hvori det
stundom vel endog kan være ængsteligt ved Siden af den Nød, i hvilken
Andre befinde sig, at leve i Skjønhed, Frodighed og Fylde, i et saadant
Liv kunne der vel gives særegne Tilfælde, hvori Hjertet ei engang tør
holde fast paa Det, hvad det ellers i Almindelighed er hellig Pligt at
haandhæve og holde i Agt. Saaledes som Mennesker nu eengang maae leve med
Mennesker, saaledes som vi Alle nu eengang maae lide Tryk og Trængsler
tilsammen med hinanden, og ikke ansee nogen menneskelig Nød for fremmed
for os selv, kunne der vel indtræde Forhold, hvori elskede Væseners Nød
kan træde et Hjerte saa nær, at det kan føle sig kaldet og tilskyndet af
de bedste Følelser til at opoffre sit indre sjælelige Velbefindende, for
at redde Dem, hvis Vel maa være det i høi Grad magtpaaliggende. Midt
under den fulde Erkjendelse og Følelse af, hvad det er, som man
bestemmer sig til at bringe til Offer, kan dog i saadanne Tilfælde dette
Offer være det, der meest svarer til en from Kjærligheds dybe Følelser.
Men altid høre dog slige særdeles Tilfælde til de Undtagelser, der
sjeldent kunne indtræde, og det være aldrig Svagheden, men den moralske
Styrke, som bringer til at give sig saaledes hen. Hvad der bør være
helligt for hele Menneskeheden, være det og for enhver Enkelt i hans
individuelle Stilling. Sit aandelige Velvær give Mennesket ei for let og
snart hen til de Misligheder og ulyksalige Forvirringer, et Egteskab kan
have til Følge, hvilket er indgaaet imod Hjertets Tilbøielighed. Ethvert
Menneske betragte sig som en Plante i Herrens Urtegaard, og holde Hævd
over sig selv som en saadan, at den kan voxe og blomstre til Velbehag og
glædelig Beskuelse for Gud og Mennesker. Der gives Opoffrelser, hvilke,
om de end synes at forstyrre Væxten og Fylden, dog netop, naar det skeer
retteligen, fremme Livets Frodighed og Skjønhed, ja de smerteligste ere
stundom ikkun Forflyttelser til en bedre Jordbund, hen under en skjønnere
Himmel; men der kunne ogsaa gives Opoffrelser, der, om de end skee med
god Villie, dog ødelægge eller kue, fordi de angribe Existentsens inderste
Rod.

Af ganske anden Art, end de nu omtalte Tilfælde, er det Sammenstød af
stridige Bevæggrunde og Tilskyndelser, som indtræder, naar Talen er om at
knytte den for hele Livet afgjørende Forbindelse tvertimod elskede
Forældres bestemte Ønske og Villie. Hvad der hist i Almindelighed kan
opfordre til fast Modstand, vil her derimod snarere lede til at give Slip
paa Hjertets Attraa, idetmindste gjøre det langt mere betænkeligt at
følge Tilbøieligheden: netop Følelsen for den egteskabelige Forbindelses
Hellighed og høie Betydning. Just fordi den er en saa skjøn Forbindelse,
er det naturligt og en Følge af inderlige Følelsers gjennemtrængende
Styrke, hellere at forsage, end at knytte den, saalænge ikke Alt er paa
det Rene, men noget Foruroligende er tilbage, der vil forstyrre dens rene
Fred og Skjønhed. Den er for hellig til, at den skulde indgaaes, medens
netop derved andre kjære, hellige og betydningsfulde Forhold forvirres,
og Sjælen ei udeelt og ganske kan føle sig vel og give sig hen med fuld
Lyst. Vel gives der Stillinger, hvorunder det indtrædende Brud i de
forrige Forhold ei betyder saa Meget; thi ikke altid have i de enkelte
Tilfælde i Virkeligheden Forholdene den Betydning og Værd, de efter deres
oprindelige Natur skulde have. Men hvor den attraaede Forbindelse vilde
blive virkelig hjerteskjærende for agtede og agtværdige Forældre, der
maaskee, dersom deres Børn saaledes forlode dem, vilde ansee dem for
tabte, bør man dog ikke alene agte deres Modbydelighed eller Advarsel,
men ogsaa agte selv den begjærede Lyksalighed for meget, til at kjøbe den
saaledes, kjøbe den i en saa ufuldkommen Skikkelse, at den dog kun halvt
blev Det, hvad Hjertet begjærede, og hvad der alene kan fyldestgjøre det.
Ogsaa her gjelder det, at noget Helligt profaneres, naar den
egteskabelige Forbindelse indgaaes med Undertrykkelse af Hjertets bedste
Følelser. Ved at forsage gjøres der ikke heller just strax en afgjørende
Bestemmelse for hele Livet; er Kjærligheden ægte og paalidelig fra begge
Sider, saa kan den indre Sjæleforbindelse endnu forblive, og med denne,
der er det Væsentligste, kan ogsaa endnu Haabet holdes fast. Hvorimod det
Skridt, der gjøres, naar Forbindelsen, uagtet alt Det, der staaer den
imod, dog indgaaes, binder for bestandigt, uden at være ledsaget af den
Velsignelse, der er saa velgjørende for Sind og Hjerte. Mennesket vil saa
gjærne raade for sin Skjæbne alene, men ved de vigtigste Skridt være han
ikke for hastig og for selvraadig, men betænke, at det Gode maa være os
beskikket ved en høiere Styrelse, naar det skal vorde os et Gode, og at
Mennesket maa tage vigtige og betydningsfulde Hindringer som Vink til
Betænksomhed. Skal overhovedet ethvert Egteskab, der skal blive lyksaligt
paa Jorden, være besluttet i Himmelen, som vi forhen et Par Gange have
mindet om, saa bør især her den indre høiere Samstemming saa fuldkomment
indtræde, at den lader os med fuldkommen Fred og Tryghed gjøre Det, hvad
vi gjøre.

\begin{center}\rule{0.28\linewidth}{0.4pt}\end{center}

Vi komme nu til den Art af ulykkelig Kjærlighed, paa hvilken man ved
dette Udtryk snarest pleier at tænke: den haabløse, der ikke finder
Gjenkjærlighed, ja ofte endog ikke tør fordre eller begjære den, men dog
vedbliver, og nu volder Uro og Smerte. Thi selv da, naar det vilde være
Uret at begjære virkelig Gjenkjærlighed, og denne Uret ogsaa virkelig
erkjendes og føles, ja naar det endog er Kjærlighedens egen Sandhed og
Dybde, der fører med sig, at den Forelskede ikke kan ønske den, selv i
saadanne Tilfælde næres dog Tilbøieligheden let endnu og den søde Lyst og
Glæde ved Tanken paa den Tilbedte. Under al den Qvide og Smerte,
hvorunder det i sig selv sig indesluttende Sind lider, kan det dog ikke
rive sig løs fra Kilden til al denne Lidelse, men holder fast ved sin
Tilbøielighed, og gjemmer den som et usigeligt sødt Eie. Den dybe Glæde
ved et skjønt Væsen, den uendelige Hengivenhed og Meddeeltagelse, aabner
og udvider Sind og Hjerte, og ligefor Indgangen til et Himmerige kan Øiet
ikke lade af bestandigen at skue ind i det, om end Indtrædelsen ikke
forundes. Den frugtbare Spire til et heelt nyt stort og herligt Liv vil,
om den end ikke kan komme til den Udfoldning, hvortil den sigter, ja
netop da endog desmere, fængsle Sind og Betragtning og fylde Sjælen med
Forventning og Længsel. Den virkede ogsaa, saa mislig og i Almindelighed
lidet ønskelig den end ofte er, dog stundom med usigelig velgjørende
Kraft i en Menneskesjæl, og hvor det er et Menneskes Tilskikkelse, at
dets aandelige Liv skal næres og varmes ved Lidelse og Smerte, kan vist
især Den, hvis Lidelse især udspringer af en haabløs Elskov, komme til at
fornemme „disse Smerters, dybt i Hjertet, hemmelig dannende Magt“, som
Goethe kalder dem i et af sine Smaadigte. Dog holde Den, der bærer en saa
frugtbar, med saa megen Himmel befrugtet Smerte og Qvide i sit Indre, sig
selv fast, og flye Sindsadspredelsen eller en omkringfarende Griben efter
Dette og Hiint, og concentrere med Kraft sine Følelsers Fylde. Han lære:
ikke at forsage sin Elskovs indre Gehalt, dens sande indre Væsen, om han
end maa forsage Tilbøielighedens jordiske Stræben! Det, der i al Elskov
er høiere end den til sig trækkende Begjær, Glæden ved det Skjønne i et
andet Væsen, vil et sandt philosophiskt, det er et i Sandhed Ideen
elskende og for Ideen levende Gemyt holde fast midt under Tabet. Men for
at saaledes den høiere Kjærlighed kan yde sin Himmel, maa Begjæringen
forstumme ganske, og ikke Lysten efter at trække til sig og besidde
idelig incitere. Hvad der bestandigen inciterer og irriterer, og,
klyngende sig fast til alle Forestillinger under deres Omløb i Sindet,
lig en fix Idee, idelig vender tilbage igjen, det betynger derved og
trykker Sindet og levner ikke engang Rum og Frihed til at nyde den høit
tilbedte Gjenstand, naar den byder sig frem for Beskuelsen til fuld
aandelig Nydelse. Den saaledes hildede Psyche kommer da ikke engang til,
midt i sit rette Element, naar den mildeste Luft vifter og lokker, at
udfolde de sammentrykkede Vinger. Bestandigen ægget til at gribe til,
uden at kunne det, aabner Sindet sig ikke for den vederqvægende Beskuelse
og en reen Meddelelses rene Glæde. Skal Hjertet ikke forkomme ynkeligen i
det Herliges glædelige Nærhed, midt under dets i sig selv himmelske
Indflydelse, men aabent modtage det fulde, rene Indtryks qvægende Kraft,
saa maa Lystens Braad ikke bestandigen stikke og brænde. Med en kraftig
Selvbestemmelse og Selvbeherskelse maa den undertrykkes eller rives ud af
Sjælen. Hvortil da vistnok ogsaa mangt et ydre Middel er nødvendigt,
eller dog hjelpende og fremmende. Stundom hjalp allerede en fri og frisk
Sindets Henvendelse paa Livets øvrige mangfoldige herlige og skjønne Eie;
thi Livet er rigt og kan paa mange Maader sætte i nye Bevægelser, naar
der retteligen gribes ind i det, og det frembyder da Stof nok til, at
Sindet, opfyldt deraf, ogsaa i det elskende Væsens Nærhed selv kan bevæge
sig i friere, herligere, oplivende Regioner, hvorved Nydelsen af det
elskede Væsens Skjønhed, Aand og Hjerte netop kan forhøies og blive
renere og skjønnere, medens Aanden nu løsner sig fra de mere smaalige
Tanker, der forhen ideligen paa Ny kom op og hildede Sindet. Af sig selv
skeer Dette nu vistnok ikke, men kun ifølge den tiltrædende kjække og
kraftige Beslutning, der nok tør vove, at lade det skjønne Himmelske faae
sin fulde velgjørende Indflydelse, uden at den dermed forbundne Lyst skal
faae Raaderum eller
spille Mester. Det er noget saa usigeligt Qvægende, Oplivende, Noget,
der med saa velgjørende Varme gjennemtrænger Sjælen, at fryde sig ved et
elsket Væsens Skjønhed og Ynde, Hjertelighed og Livelighed og aandelige
Indtagenhed, at man just ei strax skulde forsage den Glæde, ligesom at
bade Sjælen i saa mildt og oplivende et Sollys, om end Lysten vil vise
sig uregjerlig, men at man tvertimod skal tvinge denne saa, at den ei tør
forvirre eller forurene hiin Sjæleglæde. Kun at den ei lister sig for
meget frem. Den har altid en forræderisk Tendents. Lig en gridsk Orm vil
den undergrave og udhule det skjønne Liv og de søde Følelser, naar den ei
kues. Den er en Ild, der fortærer og øder, naar den faaer Magt til at
blusse op. Den vil kaste ind i et Liv fuldt af Uro og Ufred; ja naar
Veien til Tilfredsstillelse for Lysten kun gaaer gjennem Uret, vil den
kaste ind i et Liv fuldt af forstilt, falsk og løgnagtig Fremfærd, og dog
ei komme til sit Maal eller naae Fyldestgjørelse; thi Nydelsen, der ei
kan blive en Nydelse af Hjertens Grund, vil, forstyrret, halv, og uden
Totalitet, ligesom forsvinde under Hænderne og blive ved at brænde og
tære, hvor den skulde synes stillet. Saadan Tilfredsstillelse, der endog
paa en forbryderisk Maade gaaer ud paa at ruinere hvad der for Følelsen
just skulde vise sig som noget Uforkrænkeligt hos den Elskede, uden
hvilket Elskoven selv ei kunde ville bestaae, en saadan Tilfredsstillelse
er dog i og for sig allerede en kummerlig Tilfredsstillelse. Og derfor
tage man sig da ogsaa Styrke til at undflye ganske, om det maa saa være;
thi ogsaa dette Middel er ofte nødvendigt og i sig selv sikkrere, end
hiin en mere philosophisk Aandskraft fordrende Udvei. Men det tilbedte
Billede tage man dog altid med som en qvægende, sødt henrykkende
Besiddelse, der endog ved Bortfjærnelsen fra Gjenstanden kan vinde i
Indtagenhed og Skjønhed, fordi dets dybe Grundtræk kunne træde renere
frem for det mere roligt henskuende og overskuende Blik. Det er vistnok
ikke let at modstaae stærke Tillokkelsers inciterende Magt. Det Sikkreste
er altid at flye for dem. Adskillelse løser vistnok ikke let et Baand
mellem Hjerter, ei heller vil den saa let, selv hvor ingen gjensidig
Tilbøielighed knyttede et fælles Baand, ophæve Kraften i den Tillokkelse,
hvorved Nogen drages hen til en Anden formedelst Erkjendelsen af Hjertets
sande, dybe Værd hos denne Anden. Det skal Adskillelsen da ikke heller,
men kun ophæve det Foruroligende, Ængstende og Piinlige i hiin
Tillokkelses Magt, og gjøre i Stand til mere at have Godt af hiin herlige
Erkjendelse. Fra en usund Tilbøielighed derimod, der ei engang
fremkaldtes og næredes ved en saadan Erkjendelse, kan Adskillelsen vel
befrie. Hvor det ikkun vare de blot ydre Tillokkelser, være det
legemlige eller sjælelige, hvilke fængslede og hildede Sjælen, virker
Bortfjærnelsen stundom mere og hurtigere, end man skulde ventet. Thi har
den elskede Skikkelse ikke ægte Gehalt i sig, saa vil Billedet af den,
naar de øieblikkelige Opfriskninger af Lysten ikke mere vende
bestandigen tilbage, ei kunne holde sig i Beskuelsen, fordi det ei er i
Stand til at fylde denne. Dog er, endog efter lang Tids Forløb, en igjen
indtrædende Nærmelse altid betænkelig, og det bekjendte, at gammel
Kjærlighed ikke ruster, kan kun alt for let bekræfte sig. Hine
Tillokkelser virkede saa stærkt, fordi de gik i Forbund med Noget i
Menneskets Indre, der altid er stærkt og henrivende, Skjønhedsideen og
Følelsen for det Skjønne; thi denne vaktes ved hine tillokkende Fortrins
Beskuelse, og disse kunne nu ved Tilbagemindelser og ved at kalde denne
Skjønhedsidee frem paa Ny, let igjen udøve deres forrige Magt. Især kan
dog den paa Ny indtrædende Tilnærmelse være farlig da, naar det elskede
Billede i sig selv var skjønt og herligt og værdt at elskes saa høit. Dog
dette Billede selv lade man sig ikke berøve, især ikke forvanske eller
forkomme ved nogen Selvbedøvelse eller Bitterhed.

Endog Bitterhed og Harme fremægges let ved den ulykkelige Elskov, ikke
blot mod en mere lykkelig Tredie, men ogsaa mod det elskede Væsen selv.
Forbittrelse og Vrede og en vis, endog stundom til grusomme Tanker
stigende, Lyst til fiendtligen at antaste, ja ligesom at sønderrive det
Væsen, der ikke kan bevæges ved saa megen Tilbedelse, griber let om sig
ved Siden af Elskoven, fordi, naar den stærke Begjæring og Higen kommer
op, saa er det Egoistiske i Bevægelse, men dette er efter sin Natur lige
saa nær til at yttre sig fiendtligt og bortstødende, som til at vise
Tilbøielighed. Det heftige Incitament pirrer og ægger, men irriterer og
forbittrer derfor ogsaa, naar det, fordi Tilfredsstillelsen ei kommer,
bliver piinligt. Saaledes træder da let i en ny Skikkelse Jalousiens Hede
frem, ei mindre fordærvelig, end i dens forhen betragtede Skikkelser.
Dette er det Værste, hvortil det kan komme, og rører noget Saadant sig op
i Sindet, da er det vistnok høi Tid at betænke, at det er bedre at flye,
end at lide saadan Brynde. Elskoven viser atter her sin dobbelte Natur,
hvorledes den ikke blot er af guddommelig, men ogsaa af en aldeles
jordisk Herkomst. Lysten til udelukkende Tilegnelse, til at besidde og
kalde sit, hvad der er Gjenstand for Velbehag, fordærver atter her det
skjønne Eie, der, naar der elskedes med mere philosophisk Kraft og Dybde,
kunde nydes og bevares. Derfor tænke den i den haabløse Elskov fængslede
Sjæl allerede tidligen paa at bevare sig det skjønne, fulde Billede, den
har optaget i sig. At bære et reent og herligt Billede af en elsket
Individualitets Skjønhed og Fylde og hele Eiendommelighed i sig, og netop
at bære det i sig med en reen og total Hengivenhed, fri for Begjær, er
noget usigeligt Qvægende, Oplivende og Sjælen inderligt Nærende. Just af
en saadan Natur skal ogsaa den Hvile være, hvortil den lykkeligt Elskende
skal komme ved en lykkelig, al Attraa stillende Forbindelse; og for den
haabløs Elskende skal saaledes af Kjærlighedens egen Fylde den høiere
Erkjendelse og Glæde ved det Skjønne og Guddommelige gaae frem, der kan
føre den en Tid lang i en snæver Kreds hildede og fangne Sjæl tilbage
igjen til en friere Tilvær og vorde den frugtbare Spire til et nyt i
høiere Regioner sig med ny Kraft bevægende Liv. Saa at det igjen maa gaae
i Opfyldelse, at Den, der veed at lade det Guddommelige komme til
Herredømme i sit Indre, ham maa Alt komme tilgode. Hvortil iøvrigt en vis
bestemt ret dybt i Livet indtrængende Virksomhed, enten det saa er som
Studium og Ideers Udvikling og Betragtning, eller paa en mere praktisk, i
det ydre Menneskeliv mere umiddelbart indgribende Maade, er en nødvendig
Betingelse. Især i Stand til at frigjøre Sindet og fremme den nye Livets
Kraft er en saadan Virksomhed, der staaer i nær Forbindelse med
Kjærligheden til det Evige og frembyder mangfoldig Anledning til, at høie
Ideer og Følelser for det Store og Skjønne sættes i Bevægelse paa en
Maade, der river Individet udaf sit individuelle Livs Tankekreds midt ind
i den store Natur og det almene overalt rige og betydningsfulde Liv.

\begin{center}\rule{0.28\linewidth}{0.4pt}\end{center}

Men naar vi saaledes ved den haabløse Elskov vise hen paa Kjærlighedens
egen, af personlig Erhvervelse og Besiddelse uafhængige Gehalt og Fylde,
som paa det Væld, hvoraf et nyt kraftigt Liv er i Stand til at rinde
frem, vistnok ei uden Smerte og Resignation, men dog uden at qvæle og
fortabe Virkningen af det Guddommelige (det θεῖον), der kom det bevægede
Sind nær og berørte det: hvor meget mere maae vi da vise hen paa det
Samme, naar vi spørge, hvad Den har tilbage, der, efterat have henlevet
Elskovens skjønne Vaar i hjertelig Glæde og Hengivenhed og sød Meddelelse
og Fortrolighed, nu ved en udvortes Tilskikkelse aldeles bliver berøvet,
vel ikke berøvet Gjenkjærligheden, og derfor ei heller berøvet den
Elskede, men dog dennes Nærværelse og umiddelbare Besiddelse; hvorved vi
naturligvis især tænke paa den Opløsning af Elskendes Samliv, der skeer
ved Døden, skjønt vel endnu andre næsten lige saa afgjørende Adskillelser
ere tænkelige.

Lige saa naturlig som heftig er Smerten, naar det elskende Hjerte rives
ud af den trygge rolige Besiddelse, og nu, følende sin hele Tilvær
sønderrevet, stirrer ud i den mørke skumle Fremtid. Hvad der gav Livet
saa megen Betydning og Tillokkelse, hvad der udgjorde det Middelpunct,
hvorom Livet i saa mange Henseender samlede sig, og hvorved saa mange
Tanker og Følelser, saa megen Lyst og Virksomhed sattes i Bevægelse, er
borte, og Livet synes, berøvet sit bedste Indhold, tomt og øde. Just i
den Livets Deel og Sphære, der gjennemtrængte Sindet med
Evighedsfølelser, greb Tilskikkelsen tilintetgjørende ind, og naar nu det
bespændte Hjerte vil i Tanken flytte sig hen til
den Elskede, og spørger: hvor? hvorhen?, og derfor ogsaa let kommer til
at spørge: til hvad Ende dog og hvorfor det maatte lide det smertelige
Tab, og nu ei synes sig at have et fast Punct, hvortil det kan fæste
Tanken, hvor nær ere da Fortvivlelsens Følelser og let rede til at komme
op og gribe om sig? Men hvad der lettest er i Stand til at løsne det i
Begyndelsen ofte ligesom forstenede Gemyt, saa at Smerten gjør sig Luft i
Taarer og derved mildnes: den levende Mindelse af den bortkaldte
Elskedes Godhed og Kjærlighed, hvorved dets egen dybe Kjærlighed kommer
til Liv og Bevægelse igjen, og derved ogsaa alle de skjønne Sindsrørelser,
som ligge i Kjærligheden eller følge med den --- Dette er ogsaa Det,
hvorved et nyt Middelpunct for Livet igjen kan danne sig og Tilværelsen
kan erholde nyt Indhold igjen. Ikke hvad man kalder Adspredelse eller
Distraction, er Det, der trøster eller hjelper; i det Høieste kan den yde
en Bihjelp til at skaffe Hvile, naar Smertens Udbrud ere alt for
voldsomme og angribende. Det er ikke Forglemmelse eller Sindets
Bortvendelse fra den Elskedes Billede, men netop en energisk Fastholdelse
af dette Billede, en ægte og dyb Mindelse og Ihukommelse, der er i Stand
til at opretholde og bringe ny Levekraft og Levelyst tilbage. Hiin finder
ogsaa kun liden Indgang i det sorgfulde Hjerte, der snarere støder den
fra sig. Ethvert sandt og godt Ord over den bortrykkede Elskede, enhver
levende Udmaling af sammes Elskværdighed, er derimod trøstende og
husvalende. Hvilken utrøstelig Fortrøstning er dog ogsaa den, at med
Tiden nok de stærke og inderlige Følelser skulle tabe deres Magt, og
tusinde nye Indtryk nok skulle fortrænge det nu saa heftige, idet de
skulle trænge Tanken paa den Elskede mere og mere tilbage i den Elskendes
Indre. Hvilken utrøstelig Fortrøstning, at Smerten nok skal tabe sin
Kraft, naar til det Tab, vi have lidt, endnu et nyt Tab skal komme, Tabet
af Det, vi endnu have tilbage, den herlige, Sjælen ganske udfyldende og
hentagende Erindring, hvorved vi, saa meget den end bringer Smerten til
at vælde frem paa Ny, dog saa gjærne dvæle, fordi vi ved den, midt under
Smerten, nære vor Kjærlighed og bevæge denne i vort Indre. Hvor naturligt
er det derimod netop, naar den Elskede er frarevet os, at holde
Kjærligheden og dens Væsen fast, og med den tillige at bevare den
uendelige Glæde ved den Elskede, og ligesom ganske at opfylde Sjælen med
dennes herlige Billede. Naar Sjælen med elastisk Kraft fortrænger hvad
der vil overvælde og formørke den, saa ville Elskovens Følelser selv gaae
styrkede frem af den indre Qvide, som af en rensende Ild, og
gjennemtrænge Sjælen igjen med ny Inderlighed, og fylde den igjen med
rene, skjønne Evighedsfølelser. Jo væsentligere Tabet var, jo mere
velgjørende den Elskedes Nærhed, jo fuldeligere Nydelsen af Elskovens og
Samlivets hele Skjønhed, desto Mere vil den Elskende beholde tilbage, og
desto rigere vil i Sjælens Inderste Samfundet med den Fraskilte endnu
kunne blive. Det at have eiet er allerede et herligt Eie; og det er desto
herligere, jo Mere vi eiede. Som den Elskedes Billede uformørket holdes
og bevares i en trofast Phantasie, friskt og blomstrende, saaledes kan
ogsaa Kjærligheden holde sig frisk og blomstrende i det trofaste Hjerte.
Ja nu er den netop nærmest ved at kunne forklare sig til den høieste
Reenhed, idet alt det Smaalige, hvad der ellers mangen Gang forvirrer,
falder hen, og den forstummende Begjær lader en reen Beskuelse aldeles
faae Rum. Elskoven smelter nu allerlettest sammen med Fromheds og
Gudhengivenheds rene Følelser, med Tanken paa Evigheden og med en
Fornemmelse af Herrens Aand, som den, i hvem vi alle leve og røres, og
der derfor ligesom udgjør det Medium, hvori alle De, som elskes, ere i og
hos hinanden. Adskillelsen viser her en forklarende og helligende Magt.
Den hæver den elskede Skikkelse bort, udaf Øieblikkets Trængsler og
Anfægtelser, hen over i en Region, hvor Alt staaer med rene, store Træk,
og bevarer sig i ungdommelig Friskhed. Hvad der er optaget i Erindringens
Helligdom, det beholde vi evindeligen uformørket. Og nu bringes det
derved, jo renere og klarere det staaer for os, desmere tilbage til sit
evige Udspring, at dette dets evige Udspring ret maa træde frem deri. Naar
den Gudopfyldte ellers overalt i Det, der opfylder ham med dyb
Kjærlighed, elsker sin Guds uendelige Nærhed, saa vil han nu med dobbelt
Kraft føle denne midt i sine Længsler efter den Bortrykkede og i sin
endnu vedvarende aandelige Tilknytning og Hengivenhed til denne.

Og naar nu det endnu med den forrige Inderlighed og Dybde elskende Hjerte
ret beslutter at holde sin Kjærlighed fast og at blive ved at leve for
den Elskede, saa leve det da saa fuldkomment for og med denne, at det
lader ny Livelighed igjen faae Rum i sig, og aabner sig paa Ny for Glæden
ved Livet! Netop for fuldkomment at leve den Elskede til Behag, om denne,
skjønt revet bort ved Døden, dog endnu kan --- og Dette tænker sig jo det
elskende Væsen saa gjærne --- skue til det, lade det den Levelyst virke
igjen, der i og for sig er saa oplivende og glædelig for Beskueren, og
uden hvilken alt Skjønt og Godt ikke kan leve og voxe og blomstre i
Sindet. Eller skal under den maaskee lange Venten efter Gjenforening det
trofaste elskende Væsen forkues og forkomme netop formedelst sin trofaste
Kjærlighed, og just for dens Skyld ligesom ikke kunne træde den Elskende
imøde i fuldkommen Livsfylde? Og dersom det haabende gudhengivne Hjerte
ikke kan Andet, end tænke sig Tilstanden efter Døden som den, hvori den
Gode skuer og lever i en Klarhed, Aandelighed og Salighed, hvortil vor
høieste Skuen og vor høieste Sjæleglæde ikkun forholder sig som et Glimt
til det fulde Lys, saa give det ogsaa Plads for det Glædelige i den Tanke,
at den Elskede er hengangen til hiint Aandelighedens og Herlighedens
Rige! Er det Kjærlighedens Natur, at den bringer et Individ til, saaledes
at give sig hen til et andet og glemme sig selv i og for det andet, at den
Elskedes Vel i høieste Grad ligger den Elskende paa Hjertet, og at Tanken
paa den Elskedes Lyksalighed er høist lyksaliggjørende, saa vise
Kjærligheden da her den samme Kraft, og, med den vedvarende og nu sig
endog mere samlende og forstærkende Kjærlighed, vinde ogsaa Glæden over,
at den Hedengangne nu har det godt, Rum og Magt i det elskende Gemyt.
Endog om det er Døden, der har skilt de Elskende fra hinanden, kunne dog
efterhaanden de Ord begynde at finde Anvendelse, hvilke Goethe i sin
Tasso lader den ædle Prindsesse Eleonore sige, at med en fjærn Ven, naar
vi vide ham lykkelig, leve vi endnu sammen paa en ret særdeles venlig
Maade. Det knytter overhovedet det følende og haabende Menneske paa en
særdeles inderlig Maade til det Hisset, vi Alle gaae imøde, at vide
elskede Væsener der, i tryg Behold, uanfægtede af Jordelivets
Omvæltninger og Forvirringer. Naar den elskelige Skikkelse virkelig
indeholdt Det, der var værdigt til at udfylde et dybt følende Væsens
Kjærlighed og dets Længsel, og den nu lever saaledes fort i Erindringens
Himmel, saa lad Glæden over, at et saadant Væsen har levet, træde til
den, at det nu er i tryg Behold og i sit Hjem.

Og saaledes kan da ogsaa Smerten over den Elskedes Tab vise en helligende
Magt og igjennem den en skjøn Vielse til et reent og gudshengivent Liv
beredes Den, der tager Tilskikkelsen, som enhver Tilskikkelse skal og kan
tages, saaledes, at den kan og maa komme ham tilgode.

Men for at tage den saa, maa man ikke deri skue blot en Skjæbne, men
ogsaa en Tilskikkelse, nemlig en Tilskikkelse fra den Gud, hvis Veie ere
urandsagelige, men hos hvem der er beredet den Gudfrygtige et Sted, hvor
en Herlighed skal aabenbares paa ham, mod hvilken dette Livs Lidelser ere
Intet. Som blot og bar Skjæbne betragtet er Lidelsen kold, haard og
iisnende, som et Forsyns Tilskikkelse betragtet trækker den det af Sorg
bevægede Hjerte hen til et uendeligt rigt Hjerte, hos hvilket alle
Skabninger, under den ydre Naturnødvendigheds jernhaarde, ubøielige Gang,
have en Tilflugt, og hos hvilket Alt er i trygge Hænder, men fra hvilket
ogsaa enhver Skjæbne er at tage som en Stemme og en Paamindelse. Tages
den saaledes, da sætter den Hengivenheds og Ydmygheds dybe Følelser i
Bevægelse, og naar kun Mennesket ikke i sin Fortvivlelse vil gaae i Rette
med Herren, som om det havde lidt Uret, fordi det ei nød den Gave, der
altid var en reen Naade, længe nok, saa kan selv en saa bitter, det
Aandeligste i Livet saa stærkt berørende, Lidelse blive, hvad Lidelsen
saa ofte var i et Liv, hvori Lidelsen endnu hører med til det Bedste, vi
have: et Vækkelsesmiddel og en Kilde for et nyt Liv, der endnu med dybere
Følelse og inderligere Kraft, end forhen, leves som et kæmpende og sig
rensende Liv, som et Liv i Gud og et Liv for Evigheden. At døe --- saa
lærte allerede Oldtiden, og Sokrates kalder hos Platon i sine sidste
Timer denne Tanke frem for sine Venners Betragtning, og atter
Christendommen viser paa mange Maader hen paa det Samme --- at døe, det er
Noget, den Vise hele Livet igjennem stræber efter og øver sig i. Nemlig
at skilles ret fra det Jordiske og Sandselige og renses derfor og at
bringe det saavidt, at Aanden ret kan leve reent aandeligt for sig selv
og i sig selv, d. e. døe af fra det Jordiske og fra al Begjær. Af en
saadan Død vælder nyt Liv frem, ja hos det sunde Menneske ny frisk Glæde
ved Livet og Lyst til at gribe ind i og at tage Deel i dets ydre
Mangfoldighed og at nyde ogsaa Øieblikket og hvad Tingenes idelige Vexel
bringer med. Dette kan just Den nyde renest og friest, der, fri for al
Lidenskab og al brændende Begjær, ikke har fæstet Hjertet til de vexlende
Ting, eller gjort dem til Grundlag og Sæde eller til Middelpunct for sit
Livs Glæde, men under al Livets Vexel stræber efter at leve og røres i
det Uomskiftelige. Men til at komme saavidt, at al særlig Begjær og Higen
med dens Hede og Uro gaaer op i en udeelt Glæde ved det Evige og i en
reen henstrømmende Virksomhed, bringes Mennesket paa mange Veie. Hvortil
Lidelse og Smerte skulle føre Hine, til det Samme skal netop den fulde
Tilfredsstillelse og Fyldestgjørelse føre Dem, den forundtes i sin fulde
Udstrækning: thi ogsaa hos dem skulle derved Livets Inciteringer stilles,
og et roligt, trygt, lidenskabsfrit Liv indtræde. Og det kan ofte være
uvist, naar der skues hen paa Livet i det Hele, hvis Lod der var meest at
prise, enten Dens, der nød Samlivet med den Elskede uafbrudt sit hele Liv
igjennem, eller Dens, der, fordi det afbrødes tidligt, nu i en klar og
reen Erindring tog det --- i dets fulde Blomstren og dets reneste
Skjønhed, uden at det blev givet til Priis for Jordelivets Angreb, --- med
sig ind i sit øvrige Liv, og nu deri eiede Noget, hvoraf en ikke hvilende
Stræben efter at betragte det Evige og dvæle i dets Betragtning og virke
til dets Udbredelse gik frem i hans Aand. Thi ikke den blotte Besiddelse
er Det, der lyksaliggjør, men den Aand, hvoraf Livet sættes i Bevægelse,
og den indenfra kommende Livsfylde, der lige snart ved Besiddelse og Tab
kan kaldes frem til at udfolde sig. Kun da kan overhovedet et ægte Livet
besjælende og dannende Princip gaae frem af Elskoven, naar den forklarer
sig til en Aandelighed og Reenhed, i hvilken Glæden ved den Elskede
smelter sammen med Glæden ved det Eviges i Enhver tilstedeværende
Herlighed og Fylde, med Glæden ved den Gud, i hvem vi alle ere, leve og
røres.

\begin{center}\rule{0.28\linewidth}{0.4pt}\end{center}

\end{document}
